% Learning Quantitative Finance in Rust
% Template: kaobook
% Compilation: remote via Podman (latex-remote skill)
%
% Structure:
%   book/*.cls, *.sty        - kaobook template files (root for LaTeX compatibility)
%   book/chapters/*.tex      - chapter content
%   book/figures/*.tex       - TikZ figures
%   book/build/              - compiled output

\documentclass[
	fontsize=10pt,
	twoside=false,
	secnumdepth=2,
	open=any,
]{kaobook/kaobook}

% Language configuration
\usepackage[english]{babel}
\usepackage[style=english]{csquotes}

% Essential math packages
\usepackage{amsmath, amssymb}
\usepackage{mathtools}

% TikZ graphics
\usepackage{tikz}
\usepackage{pgfplots}
\pgfplotsset{compat=1.18}
\usetikzlibrary{
    arrows.meta,
    positioning,
    calc,
    shapes.geometric,
    patterns,
    decorations.pathreplacing,
    fit,
    backgrounds,
    matrix,
    datavisualization
}

% Tables
\usepackage{longtable}
\usepackage{booktabs}
\usepackage{pdflscape}
\usepackage{multirow}
\usepackage{array}

% Fix \refeq conflict before loading kaorefs
\let\refeq\relax

% kaobook packages (from kaobook/ folder)
\usepackage{kaobook/kaotheorems}
\usepackage{kaobook/kaorefs}
\usepackage{kaobook/kaobiblio}
\addbibresource{references.bib}

% Image and figure paths
\graphicspath{{images/}{figures/}{./}}

% TikZ figure path (for \input{figures/...})
\makeatletter
\g@addto@macro\input@path{{figures/}}
\makeatother

% Index
\makeindex[columns=3, title=Index, intoc]

%----------------------------------------------------------------------------------------
%	CHAPTER FORMAT
%----------------------------------------------------------------------------------------

\renewcommand*{\chapterformat}{\thechapter\enskip}
\renewcommand*{\chaptermarkformat}{\thechapter\enskip}

%----------------------------------------------------------------------------------------
%	FINANCE AND STATISTICS COMMANDS
%----------------------------------------------------------------------------------------

% Probability and statistics
\newcommand{\E}{\mathbb{E}}
\newcommand{\Var}{\text{Var}}
\newcommand{\Cov}{\text{Cov}}
\newcommand{\Corr}{\text{Corr}}
\newcommand{\Normal}{\mathcal{N}}
\DeclareMathOperator*{\argmin}{argmin}
\DeclareMathOperator*{\argmax}{argmax}

% Returns
\newcommand{\simplereturn}{\frac{P_t - P_{t-1}}{P_{t-1}}}
\newcommand{\logreturn}{\ln\left(\frac{P_t}{P_{t-1}}\right)}

% Moments
\newcommand{\skew}{\text{Skew}}
\newcommand{\kurt}{\text{Kurt}}
\newcommand{\excesskurt}{\text{ExcessKurt}}

% Time series
\newcommand{\acf}{\text{ACF}}
\newcommand{\pacf}{\text{PACF}}
\newcommand{\adf}{\text{ADF}}

% Finance metrics
\newcommand{\sharpe}{\text{Sharpe}}
\newcommand{\maxdd}{\text{MaxDD}}
\newcommand{\vol}{\sigma}

% Crate and code references
\newcommand{\crate}[1]{\texttt{#1}}
\newcommand{\module}[1]{\texttt{#1}}
\newcommand{\func}[1]{\texttt{#1()}}

%----------------------------------------------------------------------------------------
%	COLORED BOXES
%----------------------------------------------------------------------------------------

\usepackage{tcolorbox}
\tcbuselibrary{skins,breakable}

% Rust code listings
\usepackage{listings}
\usepackage{xcolor}

\definecolor{rustbg}{RGB}{250,250,250}
\definecolor{rustcomment}{RGB}{128,128,128}
\definecolor{rustkeyword}{RGB}{0,0,200}
\definecolor{ruststring}{RGB}{0,128,0}

\lstdefinelanguage{Rust}{
    keywords={fn, let, mut, pub, struct, enum, impl, trait, for, while, loop,
              if, else, match, return, use, mod, crate, self, Self, where,
              async, await, move, ref, const, static, type, dyn, Box, Vec,
              Option, Result, Ok, Err, Some, None, true, false},
    keywordstyle=\color{rustkeyword}\bfseries,
    sensitive=true,
    comment=[l]{//},
    morecomment=[s]{/*}{*/},
    commentstyle=\color{rustcomment}\itshape,
    stringstyle=\color{ruststring},
    morestring=[b]",
    morestring=[b]',
    backgroundcolor=\color{rustbg},
    basicstyle=\ttfamily\scriptsize,
    breaklines=true,
    showstringspaces=false,
    tabsize=4,
    frame=single,
    framerule=0.5pt,
    rulecolor=\color{gray!50},
}

\lstset{language=Rust}

% Companion code box
\newtcolorbox{companioncode}[1][]{
    enhanced,
    breakable,
    colback=blue!5!white,
    colframe=blue!60!black,
    fonttitle=\bfseries,
    fontupper=\small,
    coltitle=white,
    title={Companion Code: \crate{#1}},
    left=5pt, right=5pt, top=5pt, bottom=5pt
}

% Key insight box
\newtcolorbox{keyinsight}[1][]{
    enhanced,
    breakable,
    colback=yellow!10!white,
    colframe=yellow!60!black,
    fonttitle=\bfseries,
    title={Key Insight},
    left=5pt, right=5pt, top=5pt, bottom=5pt,
    #1
}

% Warning box
\newtcolorbox{warning}[1][]{
    enhanced,
    breakable,
    colback=red!5!white,
    colframe=red!70!black,
    fonttitle=\bfseries,
    title={Warning},
    left=5pt, right=5pt, top=5pt, bottom=5pt,
    #1
}

% Formula margin note
\newcommand{\formulamargem}[2]{%
    \marginnote{%
        \small\textbf{#1}\\[0.5em]
        #2
    }%
}

% Intuition margin note
\newcommand{\intuicaomargem}[2]{%
    \marginnote{%
        \small\textbf{Intuition:}\\[0.2em]
        {\footnotesize #2}%
    }%
}

%----------------------------------------------------------------------------------------
%	BEGIN DOCUMENT
%----------------------------------------------------------------------------------------

\begin{document}

%----------------------------------------------------------------------------------------
%	BOOK INFORMATION
%----------------------------------------------------------------------------------------

\titlehead{quant-finance --- on-research}

\title[Learning Quant Finance in Rust]{%
    Learning Quantitative Finance in Rust\\[0.5em]
    \Large A Progressive Journey from Kaggle to Quant Research
}
\author[Marcelo Correa]{Marcelo Correa}
\date{\today}
\publishers{\href{https://github.com/bitscrafts/quant-lab-rust}{bitscrafts/quant-lab-rust}}

%----------------------------------------------------------------------------------------

\frontmatter

%----------------------------------------------------------------------------------------
%	COPYRIGHT PAGE
%----------------------------------------------------------------------------------------

\makeatletter
\uppertitleback{\@titlehead}

\lowertitleback{
	\textbf{About this book} \\
	This book documents the learning journey through quantitative finance,
    implemented from first principles in Rust. Each chapter corresponds to a
    crate in the repository and a phase in the curriculum.

	\medskip

	\textbf{Companion code} \\
    All code lives in the \crate{quant-lab-rust} repository at:\\
    \url{https://github.com/bitscrafts/quant-lab-rust}

	\medskip

	\textbf{Colophon} \\
	Compiled with \href{https://www.latex-project.org/}{\LaTeX} using the
    \href{https://github.com/fmarotta/kaobook/}{kaobook} class. Diagrams
    created with TikZ/PGFPlots.

	\medskip

	\textbf{Compilation} \\
	Compiled remotely via Podman on \today.
}
\makeatother

%----------------------------------------------------------------------------------------
%	DEDICATION
%----------------------------------------------------------------------------------------

\dedication{
	The market can stay irrational \\
    longer than you can stay solvent.\\
	\flushright -- John Maynard Keynes
}

%----------------------------------------------------------------------------------------
%	TITLE AND TOC
%----------------------------------------------------------------------------------------

\maketitle

%----------------------------------------------------------------------------------------
%	PREFACE
%----------------------------------------------------------------------------------------

\chapter*{Preface}

This book documents a learning journey: from data science fundamentals to
quantitative finance, implemented entirely in Rust. The goal is not just to
understand the concepts, but to internalize them through implementation.

\paragraph{Philosophy}
We start with beginner-friendly Kaggle projects and progressively build toward
advanced quant research tools. No black-box libraries --- the math is
hand-rolled, because the implementation IS the learning.

\paragraph{Structure}
The book is organized in five parts:
\begin{itemize}
    \item \textbf{Part I --- Foundations}: First Kaggle projects, basic stats
    \item \textbf{Part II --- Core Concepts}: Returns, volatility, backtesting
    \item \textbf{Part III --- Quantitative Methods}: Moments, time series, stationarity
    \item \textbf{Part IV --- Derivatives and Portfolio}: Options, Markowitz, factors
    \item \textbf{Part V --- Advanced Topics}: Microstructure, AFML techniques
\end{itemize}

\paragraph{How to use}
Each chapter has:
\begin{itemize}
    \item Learning objectives at the start
    \item Mathematical foundations with TikZ diagrams
    \item Rust implementation with code listings
    \item Margin notes with key formulas and insights
    \item Exercises and references
\end{itemize}

The companion code lives at \url{https://github.com/bitscrafts/quant-lab-rust}.
Clone and run examples with:
\begin{lstlisting}[language=bash]
git clone https://github.com/bitscrafts/quant-lab-rust.git
cd quant-lab-rust
cargo run -p qf-01-fraud --example fraud_analysis
\end{lstlisting}

%----------------------------------------------------------------------------------------
%	TABLE OF CONTENTS
%----------------------------------------------------------------------------------------

\begingroup

\setlength{\textheight}{230\vscale}
\etocstandarddisplaystyle
\etocstandardlines

\tableofcontents

\listoffigures

\let\cleardoublepage\bigskip
\let\clearpage\bigskip

\listoftables

\endgroup

%----------------------------------------------------------------------------------------
%	MAIN CONTENT
%----------------------------------------------------------------------------------------

\mainmatter
\setchapterstyle{kao}

%----------------------------------------------------------------------------------------
% PART I: FOUNDATIONS (Beginner)
%----------------------------------------------------------------------------------------

\pagelayout{wide}
\addpart{Foundations}
\pagelayout{margin}

% Chapter 0: Foundations (introductory, numbered as 0)
\setcounter{chapter}{-1}
\chapter{Foundations: Mathematics, Statistics, and Programming}
\label{ch:foundations}

This chapter establishes the mathematical, statistical, and programming
foundations needed to understand the rest of the book. If you are already
comfortable with these topics, feel free to skip ahead and use this chapter
as a reference.

\section{Why These Foundations Matter}

Quantitative finance sits at the intersection of three disciplines:

\begin{itemize}
    \item \textbf{Mathematics}: The language of precision---formulas, proofs,
      and rigorous reasoning
    \item \textbf{Statistics}: The science of uncertainty---data analysis,
      probability, and inference
    \item \textbf{Programming}: The tool of implementation---turning ideas
      into working systems
\end{itemize}

\marginnote{%
\textbf{The Quant Mindset}\\[0.3em]
A quantitative approach means making decisions based on data and models,
not intuition or emotion.
}

You don't need to be an expert in all three to start. This book teaches by
doing---each concept is introduced when needed and immediately applied.

\section{Mathematical Notation}

Throughout this book, we use standard mathematical notation. Here's a quick
reference:

\subsection{Sets and Numbers}

\begin{center}
\begin{tabular}{cl}
\toprule
\textbf{Symbol} & \textbf{Meaning} \\
\midrule
$\mathbb{R}$ & Real numbers (e.g., $3.14$, $-2.5$, $\sqrt{2}$) \\
$\mathbb{R}^+$ & Positive real numbers \\
$\mathbb{Z}$ & Integers ($\ldots, -2, -1, 0, 1, 2, \ldots$) \\
$\mathbb{N}$ & Natural numbers ($1, 2, 3, \ldots$) \\
$\in$ & ``is an element of'' (e.g., $x \in \mathbb{R}$ means $x$ is real) \\
$\forall$ & ``for all'' \\
$\exists$ & ``there exists'' \\
\bottomrule
\end{tabular}
\end{center}

\subsection{Summation and Products}

\marginnote{%
\textbf{Summation}\\[0.5em]
$\displaystyle\sum_{i=1}^{n} x_i = x_1 + x_2 + \cdots + x_n$
}

The summation symbol $\Sigma$ means ``add up'':
\[
\sum_{i=1}^{n} x_i = x_1 + x_2 + \cdots + x_n
\]

For example, if $x = [2, 5, 3]$, then $\sum_{i=1}^{3} x_i = 2 + 5 + 3 = 10$.

\marginnote{%
\textbf{Product}\\[0.5em]
$\displaystyle\prod_{i=1}^{n} x_i = x_1 \times x_2 \times \cdots \times x_n$
}

The product symbol $\Pi$ means ``multiply together'':
\[
\prod_{i=1}^{n} x_i = x_1 \times x_2 \times \cdots \times x_n
\]

\subsection{Functions and Calculus}

\begin{center}
\begin{tabular}{cl}
\toprule
\textbf{Symbol} & \textbf{Meaning} \\
\midrule
$f(x)$ & Function $f$ evaluated at $x$ \\
$f'(x)$ or $\frac{df}{dx}$ & Derivative of $f$ with respect to $x$ \\
$\frac{\partial f}{\partial x}$ & Partial derivative (when $f$ has multiple variables) \\
$\ln(x)$ & Natural logarithm (base $e \approx 2.718$) \\
$\exp(x)$ or $e^x$ & Exponential function \\
$\argmax_x f(x)$ & Value of $x$ that maximizes $f$ \\
$\argmin_x f(x)$ & Value of $x$ that minimizes $f$ \\
\bottomrule
\end{tabular}
\end{center}

\begin{keyinsight}
The natural logarithm $\ln$ is everywhere in finance. Key property:
$\ln(a \times b) = \ln(a) + \ln(b)$. This turns multiplication (compounding)
into addition, making calculations much easier.
\end{keyinsight}

\section{Statistics Essentials}

Statistics helps us understand and quantify uncertainty---essential when
dealing with financial markets.

\subsection{Measures of Central Tendency}

\marginnote{%
\textbf{Mean}\\[0.5em]
$\bar{x} = \dfrac{1}{n}\displaystyle\sum_{i=1}^{n} x_i$
}

The \textbf{mean} (average) is the sum divided by the count:
\[
\bar{x} = \frac{1}{n}\sum_{i=1}^{n} x_i
\]

The \textbf{median} is the middle value when data is sorted. It's more robust
to outliers than the mean.

\subsection{Measures of Spread}

\marginnote{%
\textbf{Variance}\\[0.5em]
$\sigma^2 = \dfrac{1}{n}\displaystyle\sum_{i=1}^{n} (x_i - \bar{x})^2$
}

\textbf{Variance} measures how spread out the data is:
\[
\sigma^2 = \frac{1}{n}\sum_{i=1}^{n} (x_i - \bar{x})^2
\]

\textbf{Standard deviation} is the square root of variance:
\[
\sigma = \sqrt{\sigma^2}
\]

Standard deviation has the same units as the original data, making it easier
to interpret than variance.

\begin{lstlisting}[language=Rust]
/// Calculate mean of a slice
pub fn mean(data: &[f64]) -> f64 {
    if data.is_empty() { return 0.0; }
    data.iter().sum::<f64>() / data.len() as f64
}

/// Calculate variance of a slice
pub fn variance(data: &[f64]) -> f64 {
    if data.is_empty() { return 0.0; }
    let m = mean(data);
    data.iter().map(|x| (x - m).powi(2)).sum::<f64>() / data.len() as f64
}

/// Calculate standard deviation
pub fn std_dev(data: &[f64]) -> f64 {
    variance(data).sqrt()
}
\end{lstlisting}

\subsection{Probability Distributions}

A \textbf{probability distribution} describes how likely different outcomes are.

\subsubsection{Normal Distribution (Gaussian)}

\marginnote{%
\textbf{Normal PDF}\\[0.5em]
$f(x) = \dfrac{1}{\sigma\sqrt{2\pi}} e^{-\frac{(x-\mu)^2}{2\sigma^2}}$
}

The normal distribution is the famous ``bell curve'':
\[
f(x) = \frac{1}{\sigma\sqrt{2\pi}} e^{-\frac{(x-\mu)^2}{2\sigma^2}}
\]

where $\mu$ is the mean and $\sigma$ is the standard deviation.

We write $X \sim \Normal(\mu, \sigma^2)$ to say ``$X$ follows a normal
distribution with mean $\mu$ and variance $\sigma^2$.''

\begin{warning}
Financial returns are NOT normally distributed! They have ``fat tails''---
extreme events happen much more often than the normal distribution predicts.
We explore this in Chapter~\ref{ch:moments}.
\end{warning}

\subsubsection{Uniform Distribution}

Every outcome in a range is equally likely:
\[
f(x) = \frac{1}{b - a} \quad \text{for } a \leq x \leq b
\]

\subsection{Expected Value and Variance}

\marginnote{%
\textbf{Expected Value}\\[0.5em]
$\E[X] = \displaystyle\sum_x x \cdot P(X = x)$\\[0.5em]
(discrete)\\[1em]
$\E[X] = \displaystyle\int_{-\infty}^{\infty} x \cdot f(x) \, dx$\\[0.5em]
(continuous)
}

The \textbf{expected value} $\E[X]$ is the average outcome you'd expect over
many trials:

For a discrete random variable:
\[
\E[X] = \sum_x x \cdot P(X = x)
\]

Key properties:
\begin{itemize}
    \item $\E[aX + b] = a\E[X] + b$ (linearity)
    \item $\E[X + Y] = \E[X] + \E[Y]$ (additivity)
    \item $\Var(X) = \E[X^2] - (\E[X])^2$
\end{itemize}

\section{Financial Concepts}

Before diving into quantitative methods, let's establish some fundamental
financial concepts.

\subsection{What is a Stock?}

A \textbf{stock} (or share) represents partial ownership of a company. When
you buy a stock, you own a small piece of that company and are entitled to
a share of its profits.

\marginnote{%
\textbf{Stock Price}\\[0.3em]
The price at which buyers and sellers agree to trade ownership.
}

The \textbf{stock price} is determined by supply and demand in the market.
It reflects what buyers are willing to pay and sellers are willing to accept.

\subsection{OHLCV Data}

Market data is typically recorded as \textbf{OHLCV}:

\begin{itemize}
    \item \textbf{O}pen: First price of the trading period
    \item \textbf{H}igh: Highest price during the period
    \item \textbf{L}ow: Lowest price during the period
    \item \textbf{C}lose: Last price of the period
    \item \textbf{V}olume: Number of shares traded
\end{itemize}

This data captures the essential price action within any time frame (day,
hour, minute).

\subsection{Returns}

\marginnote{%
\textbf{Simple Return}\\[0.5em]
$R_t = \dfrac{P_t - P_{t-1}}{P_{t-1}}$\\[1em]
\textbf{Log Return}\\[0.5em]
$r_t = \ln\left(\dfrac{P_t}{P_{t-1}}\right)$
}

A \textbf{return} measures how much an investment gained or lost:

\textbf{Simple return}:
\[
R_t = \frac{P_t - P_{t-1}}{P_{t-1}}
\]

\textbf{Log return} (continuously compounded):
\[
r_t = \ln\left(\frac{P_t}{P_{t-1}}\right)
\]

Log returns have a useful property: they add up over time.
\[
r_{t_1 \to t_n} = r_{t_1} + r_{t_2} + \cdots + r_{t_n}
\]

\subsection{Risk and Volatility}

\textbf{Risk} in finance is typically measured by \textbf{volatility}---the
standard deviation of returns.

\marginnote{%
\textbf{Volatility}\\[0.5em]
$\sigma = \text{std}(r_t)$\\[0.5em]
Often annualized: $\sigma_{\text{annual}} = \sigma_{\text{daily}} \times \sqrt{252}$
}

Higher volatility means more uncertainty---prices swing more wildly.

\textbf{Annualized volatility} scales daily volatility to a yearly figure:
\[
\sigma_{\text{annual}} = \sigma_{\text{daily}} \times \sqrt{252}
\]

(252 is the typical number of trading days per year.)

\subsection{The Risk-Return Tradeoff}

\begin{keyinsight}
Higher expected returns generally require accepting higher risk. There is no
free lunch in finance---if someone promises high returns with low risk, be
skeptical.
\end{keyinsight}

The \textbf{Sharpe ratio} measures risk-adjusted return:
\[
\sharpe = \frac{\E[R] - R_f}{\sigma}
\]

where $R_f$ is the risk-free rate (e.g., Treasury bills).

\section{Programming in Rust}

This book uses Rust for all implementations. Here's a quick introduction for
readers unfamiliar with the language.

\subsection{Why Rust?}

\begin{itemize}
    \item \textbf{Performance}: As fast as C/C++, suitable for
      high-frequency trading
    \item \textbf{Safety}: The compiler catches bugs before runtime
    \item \textbf{Modern}: Great tooling, package management, and community
    \item \textbf{Learning}: Writing correct Rust forces you to think
      carefully---great for learning
\end{itemize}

\subsection{Basic Syntax}

\begin{lstlisting}[language=Rust]
// Variables (immutable by default)
let x = 5;
let mut y = 10;  // mutable
y = 20;

// Functions
fn add(a: i32, b: i32) -> i32 {
    a + b  // no semicolon = return value
}

// Structs (like classes)
struct Stock {
    symbol: String,
    price: f64,
}

// Implementation block
impl Stock {
    fn new(symbol: &str, price: f64) -> Self {
        Stock {
            symbol: symbol.to_string(),
            price,
        }
    }

    fn value(&self, shares: u32) -> f64 {
        self.price * shares as f64
    }
}

// Using the struct
let aapl = Stock::new("AAPL", 150.0);
let portfolio_value = aapl.value(100);
\end{lstlisting}

\subsection{Traits (Interfaces)}

\marginnote{%
\textbf{Traits}\\[0.3em]
Rust's way of defining shared behavior. Similar to interfaces in Java or
protocols in Swift.
}

Traits define shared behavior:

\begin{lstlisting}[language=Rust]
// Define a trait
trait Analyzer {
    fn analyze(&self, data: &[f64]) -> f64;
}

// Implement for a struct
struct MeanAnalyzer;

impl Analyzer for MeanAnalyzer {
    fn analyze(&self, data: &[f64]) -> f64 {
        data.iter().sum::<f64>() / data.len() as f64
    }
}

// Use generically
fn run_analysis<A: Analyzer>(analyzer: &A, data: &[f64]) -> f64 {
    analyzer.analyze(data)
}
\end{lstlisting}

This book uses traits extensively to create modular, reusable components.

\subsection{Error Handling}

\begin{lstlisting}[language=Rust]
// Result type for operations that can fail
fn divide(a: f64, b: f64) -> Result<f64, String> {
    if b == 0.0 {
        Err("Division by zero".to_string())
    } else {
        Ok(a / b)
    }
}

// Using Result
match divide(10.0, 2.0) {
    Ok(result) => println!("Result: {}", result),
    Err(e) => println!("Error: {}", e),
}

// Or use the ? operator
fn calculate() -> Result<f64, String> {
    let x = divide(10.0, 2.0)?;  // propagates error
    Ok(x * 2.0)
}
\end{lstlisting}

\subsection{Iterators}

Rust's iterators are powerful and efficient:

\begin{lstlisting}[language=Rust]
let prices = vec![100.0, 102.0, 101.0, 105.0];

// Map and collect
let returns: Vec<f64> = prices.windows(2)
    .map(|w| (w[1] - w[0]) / w[0])
    .collect();

// Filter
let positive: Vec<f64> = returns.iter()
    .filter(|&&r| r > 0.0)
    .copied()
    .collect();

// Fold (reduce)
let sum: f64 = prices.iter().sum();
let product: f64 = prices.iter().product();
\end{lstlisting}

\section{Running the Examples}

All code examples in this book are available in the companion repository.

\subsection{Prerequisites}

\begin{enumerate}
    \item Install Rust: \url{https://rustup.rs/}
    \item Clone the repository:
\begin{lstlisting}[language=bash]
git clone https://github.com/bitscrafts/quant-lab-rust.git
cd quant-lab-rust
\end{lstlisting}
\end{enumerate}

\subsection{Running Examples}

Each chapter has a corresponding crate with examples:

\begin{lstlisting}[language=bash]
# Chapter 1: Credit Card Fraud
cargo run -p qf-01-fraud --example fraud_analysis

# Chapter 2: Loan Default
cargo run -p qf-02-loan --example loan_analysis

# Chapter 3: Stock Prices
cargo run -p qf-03-stocks --example stock_analysis
\end{lstlisting}

\subsection{Running Tests}

\begin{lstlisting}[language=bash]
# Run all tests
cargo test

# Run tests for a specific crate
cargo test -p qf-01-fraud
\end{lstlisting}

\section{Summary}

This chapter covered:

\begin{itemize}
    \item \textbf{Mathematical notation}: Summation, products, functions,
      logarithms
    \item \textbf{Statistics}: Mean, variance, standard deviation,
      distributions
    \item \textbf{Financial concepts}: Stocks, OHLCV, returns, volatility,
      Sharpe ratio
    \item \textbf{Rust programming}: Variables, functions, structs, traits,
      error handling
\end{itemize}

\vspace{1em}
\noindent\textbf{Next chapter}: We apply these foundations to detect credit
card fraud using anomaly detection.


\chapter{Hello Finance: Credit Card Fraud Detection}
\label{ch:fraud}

\begin{companioncode}{qf-01-fraud}
Code: \texttt{crates/qf-01-fraud/}\\
Dependencies: \crate{qf-common}\\
Run: \texttt{cargo run -p qf-01-fraud --example fraud\_analysis}
\end{companioncode}

\section{Learning Objectives}

In this introductory chapter, we build our first quantitative finance application:
a fraud detection system. Along the way, we learn:

\begin{itemize}
    \item How to load and parse financial CSV datasets in Rust
    \item Basic statistical measures: mean and standard deviation
    \item Z-score anomaly detection---a simple but effective baseline
    \item Evaluation metrics: confusion matrix, precision, recall, F1 score
    \item The trait-based architecture that will guide all future implementations
\end{itemize}

\begin{keyinsight}
This is not throwaway code. Every component we build follows a \textbf{library-first design}---composable, testable, and reusable in future projects like \crate{b3-swing-trading} and \crate{stock-sim}.
\end{keyinsight}

\section{Introduction: Why Fraud Detection?}

\marginnote{%
\textbf{Class Imbalance}\\[0.3em]
492 frauds in 284,807 transactions = 0.172\%. This extreme imbalance is typical in financial anomaly detection.
}

Credit card fraud detection is a classic binary classification problem in quantitative finance. Banks and payment processors must distinguish legitimate transactions from fraudulent ones in real-time, balancing two competing objectives:

\begin{itemize}
    \item \textbf{Precision}: Minimize false positives (legitimate transactions flagged as fraud). High false positive rates annoy customers and increase operational costs.
    \item \textbf{Recall}: Maximize detection of actual fraud. Missed fraud means direct financial loss.
\end{itemize}

We use the Kaggle Credit Card Fraud Detection dataset\footnote{\url{https://www.kaggle.com/datasets/mlg-ulb/creditcardfraud}}: 284,807 European transactions from September 2013, with only 492 frauds.

\section{Data Structure}

\subsection{The Dataset}

The dataset contains:
\begin{itemize}
    \item \textbf{Time}: Seconds elapsed since first transaction
    \item \textbf{V1--V28}: PCA-transformed features (original features anonymized for privacy)
    \item \textbf{Amount}: Transaction amount in euros
    \item \textbf{Class}: Label (0 = normal, 1 = fraud)
\end{itemize}

\marginnote{%
\textbf{Why PCA features?}\\[0.3em]
Banks anonymize raw features (merchant, location, time patterns) via PCA to protect customer privacy while preserving predictive power.
}

\subsection{Rust Data Model}

We define a \texttt{Transaction} struct in \crate{qf-common}:

\begin{lstlisting}[language=Rust]
/// Represents a single financial transaction.
#[derive(Debug, Clone)]
pub struct Transaction {
    /// PCA-transformed features (V1, V2, ..., V28).
    pub features: Vec<f64>,

    /// Transaction amount in dollars/euros.
    pub amount: f64,

    /// Class label: 0 = normal, 1 = fraud.
    pub class: u8,
}
\end{lstlisting}

\subsection{Loading Data}

The CSV loader handles errors gracefully using Rust's \texttt{Result} type:

\begin{lstlisting}[language=Rust]
pub fn load_transactions(path: impl AsRef<Path>)
    -> Result<Vec<Transaction>, CommonError>
{
    if !path.as_ref().exists() {
        return Err(CommonError::FileNotFound(
            path.as_ref().display().to_string()
        ));
    }

    let mut reader = csv::Reader::from_path(path)?;
    let mut transactions = Vec::new();

    for result in reader.records() {
        let record = result?;
        // Parse features, amount, class...
        transactions.push(Transaction { features, amount, class });
    }

    Ok(transactions)
}
\end{lstlisting}

\section{Z-Score Anomaly Detection}

\subsection{Statistical Foundation}

\marginnote{%
\textbf{Z-Score Formula}\\[0.5em]
$z = \dfrac{|x - \mu|}{\sigma}$\\[0.5em]
where $\mu$ is the mean and $\sigma$ is the standard deviation.
}

Our baseline detector uses \textbf{z-score outlier detection}. The intuition: if a transaction's feature values are ``too far'' from the typical values, it might be fraudulent.

For each feature $i$:
\begin{enumerate}
    \item Compute mean $\mu_i$ and standard deviation $\sigma_i$ from training data
    \item For a new transaction with feature value $x_i$, compute the z-score:
    \[z_i = \frac{|x_i - \mu_i|}{\sigma_i}\]
    \item Flag as fraud if \textbf{any} feature exceeds a threshold (typically $z > 3$)
\end{enumerate}

% TikZ figure temporarily simplified due to kaobook compatibility
% TODO: Add proper TikZ visualization of normal distribution with rejection regions

\begin{center}
\fbox{\parbox{0.8\textwidth}{
\textbf{Z-Score Detection Visualization}\\[0.5em]
Imagine a normal (bell) curve. The threshold at $z = 3$ corresponds to the tails
of the distribution---only 0.3\% of normally distributed data falls beyond
$\pm 3\sigma$. Transactions in these extreme regions are flagged as potential
fraud.
}}
\end{center}

\marginnote{%
\textbf{The 68-95-99.7 Rule}\\[0.3em]
For normally distributed data:\\
$\pm 1\sigma$: 68.3\%\\
$\pm 2\sigma$: 95.4\%\\
$\pm 3\sigma$: 99.7\%\\[0.3em]
So $z > 3$ catches only 0.3\% of normal data as false positives.
}

\subsection{Rust Implementation}

The \texttt{ZScoreDetector} follows our trait-based design philosophy:

\begin{lstlisting}[language=Rust]
#[derive(Debug, Clone)]
pub struct ZScoreDetector {
    threshold: f64,
    feature_means: Vec<f64>,
    feature_stds: Vec<f64>,
}

impl ZScoreDetector {
    pub fn new(threshold: f64) -> Self {
        Self {
            threshold,
            feature_means: Vec::new(),
            feature_stds: Vec::new(),
        }
    }

    pub fn fit(&mut self, transactions: &[Transaction]) {
        // Compute mean for each feature
        for tx in transactions {
            for (i, &value) in tx.features.iter().enumerate() {
                self.feature_means[i] += value;
            }
        }
        for mean in &mut self.feature_means {
            *mean /= transactions.len() as f64;
        }

        // Compute standard deviation
        for tx in transactions {
            for (i, &value) in tx.features.iter().enumerate() {
                let diff = value - self.feature_means[i];
                self.feature_stds[i] += diff * diff;
            }
        }
        for std in &mut self.feature_stds {
            *std = (*std / transactions.len() as f64).sqrt();
        }
    }

    pub fn predict(&self, transaction: &Transaction) -> bool {
        for (i, &value) in transaction.features.iter().enumerate() {
            let z = (value - self.feature_means[i]).abs()
                  / self.feature_stds[i];
            if z > self.threshold {
                return true;  // Anomaly detected
            }
        }
        false
    }
}
\end{lstlisting}

\begin{warning}
The implementation uses \textbf{population} standard deviation (divides by $n$), not sample standard deviation (divides by $n-1$). For large datasets this difference is negligible, but be aware of it.
\end{warning}

\section{Evaluation Metrics}

\subsection{Confusion Matrix}

Performance is measured via the confusion matrix, which counts:
\begin{itemize}
    \item \textbf{TP} (True Positives): Fraud correctly identified
    \item \textbf{FP} (False Positives): Normal flagged as fraud
    \item \textbf{TN} (True Negatives): Normal correctly identified
    \item \textbf{FN} (False Negatives): Fraud missed
\end{itemize}

\begin{lstlisting}[language=Rust]
#[derive(Debug, Clone, Copy)]
pub struct ConfusionMatrix {
    pub tp: usize,
    pub fp: usize,
    pub tn: usize,
    pub fn_: usize,  // `fn` is a Rust keyword
}
\end{lstlisting}

\marginnote{%
% Confusion Matrix Grid
% Usage: % Confusion Matrix Grid
% Usage: \input{figures/confusion-matrix}
\begin{tikzpicture}[scale=0.6]
    \draw[thick] (0,0) grid (2,2);
    \node at (0.5,1.5) {\tiny\textbf{TP}};
    \node at (1.5,1.5) {\tiny\textbf{FP}};
    \node at (0.5,0.5) {\tiny\textbf{FN}};
    \node at (1.5,0.5) {\tiny\textbf{TN}};
    \node[rotate=90] at (-0.3,1) {\tiny Pred};
    \node at (1,2.2) {\tiny Actual};
\end{tikzpicture}

\begin{tikzpicture}[scale=0.6]
    \draw[thick] (0,0) grid (2,2);
    \node at (0.5,1.5) {\tiny\textbf{TP}};
    \node at (1.5,1.5) {\tiny\textbf{FP}};
    \node at (0.5,0.5) {\tiny\textbf{FN}};
    \node at (1.5,0.5) {\tiny\textbf{TN}};
    \node[rotate=90] at (-0.3,1) {\tiny Pred};
    \node at (1,2.2) {\tiny Actual};
\end{tikzpicture}

}

\subsection{Derived Metrics}

From the confusion matrix, we compute:

\marginnote{%
\textbf{Metric Formulas}\\[0.5em]
Precision: $\dfrac{TP}{TP + FP}$\\[0.8em]
Recall: $\dfrac{TP}{TP + FN}$\\[0.8em]
F1: $\dfrac{2 \cdot P \cdot R}{P + R}$
}

\begin{itemize}
    \item \textbf{Precision}: Of predicted frauds, how many are real?
    \[P = \frac{TP}{TP + FP}\]

    \item \textbf{Recall} (Sensitivity): Of actual frauds, how many did we catch?
    \[R = \frac{TP}{TP + FN}\]

    \item \textbf{F1 Score}: Harmonic mean of precision and recall
    \[F_1 = \frac{2PR}{P + R}\]
\end{itemize}

\begin{lstlisting}[language=Rust]
impl ConfusionMatrix {
    pub fn precision(&self) -> f64 {
        let denom = self.tp + self.fp;
        if denom == 0 { 0.0 } else { self.tp as f64 / denom as f64 }
    }

    pub fn recall(&self) -> f64 {
        let denom = self.tp + self.fn_;
        if denom == 0 { 0.0 } else { self.tp as f64 / denom as f64 }
    }

    pub fn f1(&self) -> f64 {
        let p = self.precision();
        let r = self.recall();
        if p + r == 0.0 { 0.0 } else { 2.0 * p * r / (p + r) }
    }
}
\end{lstlisting}

\subsection{Evaluation Function}

The \func{evaluate} function runs the detector on test data:

\begin{lstlisting}[language=Rust]
pub fn evaluate(
    detector: &ZScoreDetector,
    transactions: &[Transaction]
) -> ConfusionMatrix {
    let mut tp = 0; let mut fp = 0;
    let mut tn = 0; let mut fn_ = 0;

    for tx in transactions {
        let predicted = detector.predict(tx);
        let actual = tx.class == 1;

        match (predicted, actual) {
            (true, true)   => tp += 1,
            (true, false)  => fp += 1,
            (false, false) => tn += 1,
            (false, true)  => fn_ += 1,
        }
    }

    ConfusionMatrix { tp, fp, tn, fn_ }
}
\end{lstlisting}

\section{Toward a Trait-Based Architecture}

\begin{keyinsight}
This implementation will evolve into \crate{quant-lib}. Before writing any struct, we ask: \textbf{what trait should this implement?}
\end{keyinsight}

\marginnote{%
\textbf{Future Traits}\\[0.3em]
\texttt{AnomalyDetector}\\
\texttt{BinaryClassifier}\\
\texttt{Evaluator<T>}\\
\texttt{DataSource}
}

While Phase 1 uses concrete types, future phases will extract traits:

\begin{lstlisting}[language=Rust]
// Future: trait-based design
pub trait AnomalyDetector {
    fn fit(&mut self, data: &[Transaction]) -> Result<(), Error>;
    fn predict(&self, tx: &Transaction) -> bool;
    fn score(&self, tx: &Transaction) -> f64;  // anomaly score
}

pub trait Evaluator<M> {
    fn evaluate(&self, preds: &[bool], actuals: &[bool]) -> M;
}

// Then any detector + any evaluator compose:
fn backtest<D, E>(detector: &D, eval: &E, data: &[Transaction])
where
    D: AnomalyDetector,
    E: Evaluator<ConfusionMatrix>,
{ /* ... */ }
\end{lstlisting}

This composability is why we invest in clean abstractions from the start.

\section{Results and Analysis}

\subsection{Running the Example}

\begin{lstlisting}[language=bash]
$ cargo run -p qf-01-fraud --example fraud_analysis
Loading transactions from data/creditcard.csv...
Loaded 284807 transactions (492 frauds)
Training detector with threshold 3.0...
Evaluating on training set...

Results:
  True Positives:  312
  False Positives: 1847
  True Negatives:  282468
  False Negatives: 180

  Precision: 0.1446
  Recall:    0.6341
  F1 Score:  0.2356
\end{lstlisting}

\subsection{Interpreting Results}

The results reveal characteristic tradeoffs:

\begin{itemize}
    \item \textbf{Recall = 0.63}: We catch 63\% of frauds. Not bad for a simple baseline!
    \item \textbf{Precision = 0.14}: Of our ``fraud'' predictions, only 14\% are real. This means many false alarms.
    \item \textbf{F1 = 0.24}: The harmonic mean reflects this imbalance.
\end{itemize}

\subsection{Limitations}

\begin{warning}
We evaluated on the training set---a cardinal sin! True performance requires a held-out test set. Exercise 3 addresses this.
\end{warning}

This baseline approach has several fundamental limitations:

\begin{enumerate}
    \item \textbf{Gaussian assumption}: Z-scores assume normally distributed features. Financial data often has heavy tails (more extreme values than Gaussian predicts).

    \item \textbf{Equal feature weighting}: All features are treated equally. Some PCA components may be more predictive than others.

    \item \textbf{Independent features}: We check each feature separately. Sophisticated fraud may only be detectable via feature \emph{combinations}.

    \item \textbf{Fixed threshold}: The threshold $z = 3$ is arbitrary. ROC analysis (Chapter 2) shows how to optimize it.
\end{enumerate}

\section{Exercises}

\begin{enumerate}
    \item \textbf{Threshold Sensitivity}: Try thresholds of 2.0, 2.5, 3.0, 3.5, 4.0. Plot precision vs. recall. What threshold maximizes F1?

    \item \textbf{Feature Analysis}: Modify \func{predict} to return which feature(s) triggered the alert. Which features are most discriminative for fraud?

    \item \textbf{Train/Test Split}: Split the data 80/20 chronologically (not randomly---why?). How does test-set performance compare to training-set?

    \item \textbf{Amount Feature}: The current implementation ignores the \texttt{amount} field. Add it as a feature. Does performance improve?

    \item \textbf{Trait Extraction}: Refactor \texttt{ZScoreDetector} to implement an \texttt{AnomalyDetector} trait. What methods should the trait require?
\end{enumerate}

\section{Key Takeaways}

\begin{itemize}
    \item \textbf{CSV loading} in Rust with proper error handling via \texttt{Result}
    \item \textbf{Basic statistics}: mean, standard deviation, z-score
    \item \textbf{Binary classification metrics}: confusion matrix, precision, recall, F1
    \item \textbf{Precision-recall tradeoff}: you can't maximize both simultaneously
    \item \textbf{Evaluation pitfalls}: never evaluate on training data!
    \item \textbf{Library-first design}: even simple code should be composable
\end{itemize}

\vspace{1em}
\noindent\textbf{Next chapter}: Feature engineering and logistic scoring for loan default prediction. We'll implement ROC-AUC and learn why it's often better than F1 for imbalanced datasets.


\chapter{Risk Basics: Loan Default Prediction}
\label{ch:loan}

\begin{companioncode}{qf-02-loan}
Code: \texttt{crates/qf-02-loan/}\\
Dependencies: \crate{qf-common}\\
Run: \texttt{cargo run -p qf-02-loan --example loan\_analysis}
\end{companioncode}

\section{Learning Objectives}

Building on Chapter 1, we now tackle a more nuanced classification problem:
predicting loan defaults. This chapter introduces:

\begin{itemize}
    \item Feature engineering for credit risk assessment
    \item Categorical variable encoding (one-hot)
    \item Feature normalization (min-max scaling)
    \item Linear scoring models with probability calibration
    \item ROC curves and AUC---a better metric for imbalanced data
    \item Trait-based architecture in action
\end{itemize}

\begin{keyinsight}
While Chapter 1 used z-scores (a threshold on statistics), this chapter moves
toward \textbf{probability estimation}. We want to know not just ``default or
not'' but ``how likely is default?'' This probability enables better decisions
and risk-based pricing.
\end{keyinsight}

\section{Introduction: Credit Risk Fundamentals}

\marginnote{%
\textbf{Credit Risk}\\[0.3em]
The risk that a borrower will fail to repay a loan according to agreed terms.
}

When a bank issues a loan, it faces \textbf{credit risk}---the possibility that
the borrower defaults. Unlike fraud detection where we catch bad actors, loan
default prediction is about assessing the \emph{likelihood} of future failure.

Key questions in credit risk:
\begin{itemize}
    \item What factors predict default? (Feature engineering)
    \item How do we combine factors into a score? (Scoring model)
    \item How do we evaluate the score's usefulness? (ROC-AUC)
\end{itemize}

\section{Data Model}

\subsection{Loan Application Structure}

\begin{lstlisting}[language=Rust]
pub struct LoanApplication {
    pub income: f64,           // Annual income
    pub loan_amount: f64,      // Requested loan amount
    pub interest_rate: f64,    // Interest rate offered
    pub dti: f64,              // Debt-to-income ratio
    pub grade: u8,             // Credit grade (A=1 to G=7)
    pub home_ownership: HomeOwnership,
    pub defaulted: bool,       // Target: did they default?
}

pub enum HomeOwnership {
    Rent,
    Own,
    Mortgage,
}
\end{lstlisting}

\marginnote{%
\textbf{Credit Grades}\\[0.3em]
A = lowest risk\\
G = highest risk\\[0.5em]
Each grade maps to different interest rates and approval criteria.
}

\section{Feature Engineering}

\subsection{Derived Features}

Raw data rarely enters a model directly. We engineer features that capture
meaningful relationships:

\marginnote{%
\textbf{Feature Formulas}\\[0.5em]
DTI: $\dfrac{\text{debt}}{\text{income}}$\\[0.8em]
LTI: $\dfrac{\text{loan}}{\text{income}}$\\[0.8em]
Grade Score: $\dfrac{g - 1}{6}$
}

\begin{lstlisting}[language=Rust]
pub struct FeatureExtractor;

impl FeatureExtractor {
    /// Debt-to-income ratio: total debt / income
    pub fn debt_to_income(income: f64, debt: f64) -> f64 {
        if income == 0.0 { return 0.0; }
        debt / income
    }

    /// Loan-to-income ratio: loan amount / income
    pub fn loan_to_income(loan_amount: f64, income: f64) -> f64 {
        if income == 0.0 { return 0.0; }
        loan_amount / income
    }

    /// Convert grade (1-7) to risk score (0.0-1.0)
    pub fn grade_to_score(grade: u8) -> f64 {
        ((grade.saturating_sub(1)) as f64) / 6.0
    }
}
\end{lstlisting}

\subsection{Categorical Encoding}

Numerical models require numerical inputs. We encode categorical variables
using \textbf{one-hot encoding}:

\begin{center}
\begin{tabular}{lccc}
\toprule
\textbf{Home Ownership} & \textbf{RENT} & \textbf{OWN} & \textbf{MORTGAGE} \\
\midrule
Rent     & 1 & 0 & 0 \\
Own      & 0 & 1 & 0 \\
Mortgage & 0 & 0 & 1 \\
\bottomrule
\end{tabular}
\end{center}

\begin{lstlisting}[language=Rust]
pub struct OneHotEncoder;

impl OneHotEncoder {
    /// Encode home ownership as [RENT, OWN, MORTGAGE]
    pub fn encode_home_ownership(h: &HomeOwnership) -> [f64; 3] {
        match h {
            HomeOwnership::Rent     => [1.0, 0.0, 0.0],
            HomeOwnership::Own      => [0.0, 1.0, 0.0],
            HomeOwnership::Mortgage => [0.0, 0.0, 1.0],
        }
    }
}
\end{lstlisting}

\subsection{Feature Normalization}

\marginnote{%
\textbf{Min-Max Scaling}\\[0.5em]
$x' = \dfrac{x - x_{\min}}{x_{\max} - x_{\min}}$\\[0.5em]
Scales features to $[0, 1]$.
}

Different features have different scales. Income might be in tens of thousands
while interest rates are single digits. We normalize to $[0, 1]$:

\begin{lstlisting}[language=Rust]
pub struct Normalizer {
    min: Vec<f64>,
    max: Vec<f64>,
}

impl Normalizer {
    pub fn fit(data: &[Vec<f64>]) -> Self {
        // Compute min/max per feature
    }

    pub fn transform(&self, features: &[f64]) -> Vec<f64> {
        features.iter().enumerate().map(|(i, &x)| {
            let range = self.max[i] - self.min[i];
            if range == 0.0 { 0.0 } else { (x - self.min[i]) / range }
        }).collect()
    }
}
\end{lstlisting}

\section{Linear Scoring Model}

\subsection{Weighted Sum}

The simplest scoring model is a weighted linear combination:

\[
\text{score} = \sum_{i=1}^{n} w_i \cdot x_i + b = \mathbf{w}^T \mathbf{x} + b
\]

\marginnote{%
\textbf{Linear Score}\\[0.5em]
$s = \mathbf{w}^T \mathbf{x} + b$\\[0.5em]
where $\mathbf{w}$ are weights, $b$ is bias.
}

\begin{lstlisting}[language=Rust]
pub struct LinearScorer {
    weights: Vec<f64>,
    bias: f64,
}

impl Scorer for LinearScorer {
    fn score(&self, features: &[f64]) -> f64 {
        let dot_product: f64 = self.weights.iter()
            .zip(features.iter())
            .map(|(w, f)| w * f)
            .sum();
        dot_product + self.bias
    }
}
\end{lstlisting}

\subsection{Probability via Sigmoid}

The raw score can be any real number. To convert it to a probability in $[0, 1]$,
we apply the \textbf{sigmoid function}:

\marginnote{%
\textbf{Sigmoid Function}\\[0.5em]
$\sigma(x) = \dfrac{1}{1 + e^{-x}}$\\[0.5em]
Maps $(-\infty, +\infty)$ to $(0, 1)$.
}

\[
\sigma(x) = \frac{1}{1 + e^{-x}}
\]

Properties of sigmoid:
\begin{itemize}
    \item $\sigma(0) = 0.5$ (neutral score = 50\% probability)
    \item $\sigma(x) \to 1$ as $x \to +\infty$
    \item $\sigma(x) \to 0$ as $x \to -\infty$
\end{itemize}

\begin{lstlisting}[language=Rust]
/// Hand-rolled sigmoid: 1 / (1 + exp(-x))
fn sigmoid(x: f64) -> f64 {
    1.0 / (1.0 + (-x).exp())
}

impl BinaryClassifier for LinearScorer {
    fn predict_proba(&self, features: &[f64]) -> f64 {
        sigmoid(self.score(features))
    }

    fn predict(&self, features: &[f64]) -> bool {
        self.predict_proba(features) > 0.5
    }
}
\end{lstlisting}

\section{ROC Curves and AUC}

\subsection{Why Not Just Accuracy?}

\marginnote{%
\textbf{Class Imbalance}\\[0.3em]
If 95\% of loans are repaid, a model that always predicts ``no default'' has
95\% accuracy but catches zero defaults!
}

Like fraud detection, loan default is imbalanced. Accuracy is misleading.
Instead, we use the \textbf{Receiver Operating Characteristic (ROC) curve}.

\subsection{ROC Curve Construction}

The ROC curve plots \textbf{True Positive Rate (TPR)} vs \textbf{False Positive
Rate (FPR)} at various classification thresholds:

\begin{itemize}
    \item \textbf{TPR} (Recall): $\frac{TP}{TP + FN}$ --- of actual defaults, how many did we catch?
    \item \textbf{FPR}: $\frac{FP}{FP + TN}$ --- of actual non-defaults, how many did we falsely flag?
\end{itemize}

\marginnote{%
\textbf{ROC Interpretation}\\[0.3em]
Top-left corner $(0, 1)$: perfect classifier.\\[0.2em]
Diagonal: random guessing.\\[0.2em]
Below diagonal: worse than random!
}

\begin{lstlisting}[language=Rust]
pub fn roc_curve(scores: &[f64], labels: &[bool]) -> Vec<(f64, f64)> {
    // Sort by score descending
    let mut pairs: Vec<(f64, bool)> = scores.iter()
        .copied()
        .zip(labels.iter().copied())
        .collect();
    pairs.sort_by(|a, b| b.0.partial_cmp(&a.0).unwrap());

    let total_pos = labels.iter().filter(|&&l| l).count();
    let total_neg = labels.len() - total_pos;

    let mut roc_points = vec![(0.0, 0.0)];
    let mut tp = 0;
    let mut fp = 0;

    for (_score, label) in pairs {
        if label { tp += 1; } else { fp += 1; }
        let fpr = fp as f64 / total_neg as f64;
        let tpr = tp as f64 / total_pos as f64;
        roc_points.push((fpr, tpr));
    }

    roc_points
}
\end{lstlisting}

\subsection{Area Under Curve (AUC)}

The \textbf{AUC} summarizes the ROC curve as a single number:

\marginnote{%
\textbf{AUC Interpretation}\\[0.3em]
$\text{AUC} = 1.0$: perfect\\
$\text{AUC} = 0.5$: random\\
$\text{AUC} < 0.5$: worse than random
}

\begin{itemize}
    \item $\text{AUC} = 1.0$: Perfect classifier (all defaults scored higher than non-defaults)
    \item $\text{AUC} = 0.5$: Random classifier (no predictive power)
    \item $\text{AUC} = 0.7$: Fair classifier (industry often requires $\geq 0.7$)
\end{itemize}

We compute AUC via trapezoidal integration:

% ROC curve figure in margin beside the auc function
\begin{marginfigure}[-1cm]
    \centering
    % ROC Curve Figure
% Usage: % ROC Curve Figure
% Usage: \input{figures/roc-curve}
\begin{tikzpicture}[scale=0.5]
    % Axes
    \draw[thick, ->] (0,0) -- (5.5,0) node[right] {\footnotesize FPR};
    \draw[thick, ->] (0,0) -- (0,5.5) node[above] {\footnotesize TPR};

    % Diagonal (random classifier)
    \draw[dashed, gray] (0,0) -- (5,5);

    % ROC curve (good classifier)
    \draw[thick, blue] plot[smooth] coordinates {
        (0,0) (0.2,1.5) (0.4,2.5) (0.7,3.5) (1.2,4.2) (2,4.6) (3,4.85) (5,5)
    };

    % Shaded AUC area
    \fill[blue!15] plot[smooth] coordinates {
        (0,0) (0.2,1.5) (0.4,2.5) (0.7,3.5) (1.2,4.2) (2,4.6) (3,4.85) (5,5)
    } -- (5,0) -- cycle;

    % Axis ticks
    \foreach \x in {0,1,...,5} {
        \draw (\x,0) -- (\x,-0.1);
    }
    \foreach \y in {0,1,...,5} {
        \draw (0,\y) -- (-0.1,\y);
    }

    % Labels
    \node at (0,-0.4) {\tiny 0};
    \node at (5,-0.4) {\tiny 1};
    \node at (-0.4,0) {\tiny 0};
    \node at (-0.4,5) {\tiny 1};

    % Legend
    \draw[thick, blue] (3.5,1) -- (4.2,1) node[right] {\tiny Classifier};
    \draw[dashed, gray] (3.5,0.3) -- (4.2,0.3) node[right] {\tiny Random};
\end{tikzpicture}

\begin{tikzpicture}[scale=0.5]
    % Axes
    \draw[thick, ->] (0,0) -- (5.5,0) node[right] {\footnotesize FPR};
    \draw[thick, ->] (0,0) -- (0,5.5) node[above] {\footnotesize TPR};

    % Diagonal (random classifier)
    \draw[dashed, gray] (0,0) -- (5,5);

    % ROC curve (good classifier)
    \draw[thick, blue] plot[smooth] coordinates {
        (0,0) (0.2,1.5) (0.4,2.5) (0.7,3.5) (1.2,4.2) (2,4.6) (3,4.85) (5,5)
    };

    % Shaded AUC area
    \fill[blue!15] plot[smooth] coordinates {
        (0,0) (0.2,1.5) (0.4,2.5) (0.7,3.5) (1.2,4.2) (2,4.6) (3,4.85) (5,5)
    } -- (5,0) -- cycle;

    % Axis ticks
    \foreach \x in {0,1,...,5} {
        \draw (\x,0) -- (\x,-0.1);
    }
    \foreach \y in {0,1,...,5} {
        \draw (0,\y) -- (-0.1,\y);
    }

    % Labels
    \node at (0,-0.4) {\tiny 0};
    \node at (5,-0.4) {\tiny 1};
    \node at (-0.4,0) {\tiny 0};
    \node at (-0.4,5) {\tiny 1};

    % Legend
    \draw[thick, blue] (3.5,1) -- (4.2,1) node[right] {\tiny Classifier};
    \draw[dashed, gray] (3.5,0.3) -- (4.2,0.3) node[right] {\tiny Random};
\end{tikzpicture}

    \caption{ROC curve: the blue line shows classifier performance, dashed line is random guessing. Area under curve (AUC) is shaded.}
    \label{fig:roc-curve}
\end{marginfigure}

\begin{lstlisting}[language=Rust]
pub fn auc(roc_points: &[(f64, f64)]) -> f64 {
    let mut area = 0.0;
    for i in 1..roc_points.len() {
        let (x0, y0) = roc_points[i - 1];
        let (x1, y1) = roc_points[i];
        // Trapezoidal rule: width * average height
        area += (x1 - x0) * (y0 + y1) / 2.0;
    }
    area
}
\end{lstlisting}

\begin{keyinsight}
AUC has a probabilistic interpretation: it equals the probability that a
randomly chosen positive (default) is ranked higher than a randomly chosen
negative (non-default).
\end{keyinsight}

\section{Trait-Based Architecture}

Following our library-first design, we define traits that enable composability:

\begin{lstlisting}[language=Rust]
/// Trait for models that produce probability predictions.
pub trait BinaryClassifier {
    fn predict(&self, features: &[f64]) -> bool;
    fn predict_proba(&self, features: &[f64]) -> f64;
}

/// Trait for models that produce raw scores.
pub trait Scorer {
    fn score(&self, features: &[f64]) -> f64;
}
\end{lstlisting}

\marginnote{%
\textbf{Future Benefits}\\[0.3em]
Any classifier implementing \texttt{BinaryClassifier} can be evaluated with our
\func{roc\_curve} and \func{auc} functions.
}

This means our ROC-AUC evaluation works with \emph{any} classifier, not just
\texttt{LinearScorer}. Future chapters will implement logistic regression,
decision trees, and ensemble methods---all reusing this evaluation code.

\section{Results and Analysis}

\subsection{Example Usage}

\begin{lstlisting}[language=bash]
$ cargo run -p qf-02-loan --example loan_analysis
Loading loan applications...
Loaded 10000 applications (850 defaults, 8.5%)

Feature extraction complete.
Training simple scorer (hardcoded weights)...

Evaluation Results:
  AUC: 0.724
  Optimal Threshold: 0.35
  At threshold 0.5:
    Precision: 0.42
    Recall: 0.58
    F1 Score: 0.49
\end{lstlisting}

\subsection{Interpreting the Results}

\begin{itemize}
    \item \textbf{AUC = 0.72}: The model has reasonable discriminative power. It correctly ranks defaults above non-defaults 72\% of the time.
    \item \textbf{Precision = 0.42}: Of predicted defaults, 42\% actually default. Many false alarms.
    \item \textbf{Recall = 0.58}: We catch 58\% of actual defaults. 42\% are missed.
\end{itemize}

\begin{warning}
We used hardcoded weights, not learned weights. Real credit scoring uses
logistic regression or gradient boosting to optimize weights from data. That
comes in later chapters.
\end{warning}

\section{Exercises}

\begin{enumerate}
    \item \textbf{Weight Sensitivity}: Change the weights in \texttt{LinearScorer}. How does AUC change? Which features matter most?

    \item \textbf{Threshold Selection}: Plot precision and recall at various thresholds. Find the threshold that maximizes F1 score.

    \item \textbf{Feature Importance}: Train separate single-feature scorers. Rank features by their individual AUC.

    \item \textbf{Cross-Validation}: Split data into 5 folds. Report mean and standard deviation of AUC across folds.

    \item \textbf{Calibration Plot}: Compare predicted probabilities to actual default rates. Is the model well-calibrated?
\end{enumerate}

\section{Key Takeaways}

\begin{itemize}
    \item \textbf{Feature engineering} transforms raw data into predictive signals
    \item \textbf{One-hot encoding} handles categorical variables
    \item \textbf{Min-max normalization} scales features to comparable ranges
    \item \textbf{Sigmoid} converts scores to probabilities
    \item \textbf{ROC-AUC} evaluates ranking quality, robust to class imbalance
    \item \textbf{Trait-based design} enables reusable evaluation code
\end{itemize}

\vspace{1em}
\noindent\textbf{Next chapter}: Market data fundamentals---OHLCV candlesticks, returns, and simple moving averages.


\chapter{Market Data: Stock Price Analysis}
\label{ch:stocks}

\begin{companioncode}{qf-03-stocks}
Code: \texttt{crates/qf-03-stocks/}\\
Dependencies: \crate{qf-common}\\
Run: \texttt{cargo run -p qf-03-stocks --example stock\_analysis}
\end{companioncode}

\section{Learning Objectives}

This chapter introduces the fundamental data structure of quantitative finance:
\textbf{OHLCV} (Open, High, Low, Close, Volume) market data. We learn:

\begin{itemize}
    \item OHLCV data structure and its significance
    \item Candlestick anatomy---body, shadows, and what they reveal
    \item Basic market statistics (range, volume, price change)
    \item Moving averages: SMA and EMA
    \item The \texttt{TimeSeries} trait for generic operations
\end{itemize}

\begin{keyinsight}
OHLCV data is the foundation of all quantitative analysis. Every strategy,
every backtest, every indicator starts here. Master this structure and you can
work with any market---stocks, forex, crypto, commodities.
\end{keyinsight}

\section{Introduction: What is OHLCV?}

\marginnote{%
\textbf{OHLCV}\\[0.3em]
Open, High, Low, Close, Volume---the five essential data points for any trading
period (day, hour, minute).
}

Every trading period (day, hour, minute) produces five key data points:

\begin{itemize}
    \item \textbf{Open}: First traded price of the period
    \item \textbf{High}: Highest price reached during the period
    \item \textbf{Low}: Lowest price reached during the period
    \item \textbf{Close}: Last traded price of the period
    \item \textbf{Volume}: Total number of shares/contracts traded
\end{itemize}

This data tells a story. The range between high and low shows volatility. The
relationship between open and close shows direction. Volume confirms conviction.

\section{Data Structure}

\subsection{Rust Implementation}

\begin{lstlisting}[language=Rust]
#[derive(Debug, Clone, Deserialize)]
pub struct Ohlcv {
    pub date: String,
    pub open: f64,
    pub high: f64,
    pub low: f64,
    pub close: f64,
    pub volume: u64,
    pub adj_close: Option<f64>,
}
\end{lstlisting}

\marginnote{%
\textbf{Adjusted Close}\\[0.3em]
Accounts for splits and dividends. If a stock splits 2:1, historical closes are
halved so charts remain continuous.
}

\subsection{Data Validation}

OHLCV data has natural constraints that we validate on load:

\begin{itemize}
    \item $\text{low} \leq \text{high}$ (always)
    \item $\text{low} \leq \text{open} \leq \text{high}$
    \item $\text{low} \leq \text{close} \leq \text{high}$
\end{itemize}

\begin{lstlisting}[language=Rust]
// Validate OHLCV relationships
if record.high < record.low {
    return Err(CommonError::ParseError(format!(
        "Row {}: high ({}) < low ({})",
        idx + 2, record.high, record.low
    )));
}
\end{lstlisting}

\section{Candlestick Anatomy}

\marginnote{%
\textbf{Candlestick Parts}\\[0.3em]
Body: $|close - open|$\\
Upper shadow: $high - \max(O, C)$\\
Lower shadow: $\min(O, C) - low$
}

A candlestick visualizes OHLCV data. Understanding its anatomy is essential:

\marginnote{%
\begin{tikzpicture}[scale=0.35]
\draw[thick, ->] (0,0) -- (6,0);
\draw[thick, ->] (0,0) -- (0,5);
% Bullish
\draw[thick] (1,1) -- (1,4);
\fill[green!60!black] (0.7,1.5) rectangle (1.3,3.5);
% Bearish
\draw[thick] (3,1.5) -- (3,4.5);
\fill[red!70!black] (2.7,2) rectangle (3.3,4);
% Doji
\draw[thick] (5,2) -- (5,4);
\draw[thick, line width=1.5pt] (4.7,3) -- (5.3,3);
\end{tikzpicture}\\[0.3em]
\textbf{Candlesticks}\\
Green: bullish\\
Red: bearish\\
Line: doji
}

\begin{center}
\begin{tabular}{ll}
\toprule
\textbf{Component} & \textbf{Formula} \\
\midrule
Daily Range & $high - low$ \\
Body & $close - open$ \\
Upper Shadow & $high - \max(open, close)$ \\
Lower Shadow & $\min(open, close) - low$ \\
Typical Price & $(high + low + close) / 3$ \\
\bottomrule
\end{tabular}
\end{center}

\subsection{Rust Implementation}

\begin{lstlisting}[language=Rust]
impl Ohlcv {
    /// Daily price range
    pub fn daily_range(&self) -> f64 {
        self.high - self.low
    }

    /// Candle body (positive = bullish)
    pub fn body(&self) -> f64 {
        self.close - self.open
    }

    /// Upper shadow length
    pub fn upper_shadow(&self) -> f64 {
        self.high - self.open.max(self.close)
    }

    /// Lower shadow length
    pub fn lower_shadow(&self) -> f64 {
        self.open.min(self.close) - self.low
    }

    /// True if close > open (green candle)
    pub fn is_bullish(&self) -> bool {
        self.close > self.open
    }

    /// Representative price: (H + L + C) / 3
    pub fn typical_price(&self) -> f64 {
        (self.high + self.low + self.close) / 3.0
    }
}
\end{lstlisting}

\section{Market Statistics}

\subsection{Basic Aggregations}

Given a series of OHLCV records, we compute summary statistics:

\begin{lstlisting}[language=Rust]
/// Average trading volume
pub fn average_volume(data: &[Ohlcv]) -> f64 {
    if data.is_empty() { return 0.0; }
    let sum: u64 = data.iter().map(|o| o.volume).sum();
    sum as f64 / data.len() as f64
}

/// Price change: (last - first) / first
pub fn price_change(data: &[Ohlcv]) -> f64 {
    if data.len() < 2 { return 0.0; }
    let first = data[0].close;
    let last = data[data.len() - 1].close;
    (last - first) / first
}

/// Highest high across all records
pub fn highest_high(data: &[Ohlcv]) -> f64 {
    data.iter().map(|o| o.high).fold(f64::NEG_INFINITY, f64::max)
}

/// Lowest low across all records
pub fn lowest_low(data: &[Ohlcv]) -> f64 {
    data.iter().map(|o| o.low).fold(f64::INFINITY, f64::min)
}
\end{lstlisting}

\marginnote{%
\textbf{Functional Style}\\[0.3em]
Note the use of \texttt{iter()}, \texttt{map()}, \texttt{fold()}---Rust's
functional combinators make statistical computations elegant.
}

\section{Moving Averages}

\subsection{Simple Moving Average (SMA)}

The SMA is the arithmetic mean of the last $n$ prices:

\marginnote{%
\textbf{SMA Formula}\\[0.5em]
$\text{SMA}_t = \dfrac{1}{n}\sum_{i=0}^{n-1} P_{t-i}$\\[0.5em]
where $n$ is the period.
}

\[
\text{SMA}_t = \frac{1}{n}\sum_{i=0}^{n-1} P_{t-i}
\]

\begin{lstlisting}[language=Rust]
pub fn sma(prices: &[f64], period: usize) -> Vec<f64> {
    if period == 0 || period > prices.len() {
        return Vec::new();
    }

    let mut result = Vec::with_capacity(prices.len() - period + 1);

    for i in 0..=(prices.len() - period) {
        let window = &prices[i..i + period];
        let sum: f64 = window.iter().sum();
        result.push(sum / period as f64);
    }

    result
}
\end{lstlisting}

\subsection{Exponential Moving Average (EMA)}

The EMA gives more weight to recent prices, making it more responsive:

\marginnote{%
\textbf{EMA Formula}\\[0.5em]
$\text{EMA}_t = \alpha \cdot P_t + (1-\alpha) \cdot \text{EMA}_{t-1}$\\[0.5em]
$\alpha = \dfrac{2}{n+1}$
}

\[
\text{EMA}_t = \alpha \cdot P_t + (1 - \alpha) \cdot \text{EMA}_{t-1}
\]

where $\alpha = \frac{2}{n+1}$ is the smoothing factor.

\begin{lstlisting}[language=Rust]
pub fn ema(prices: &[f64], period: usize) -> Vec<f64> {
    if period == 0 || prices.is_empty() {
        return Vec::new();
    }

    let alpha = 2.0 / (period + 1) as f64;
    let mut result = Vec::with_capacity(prices.len());

    // Initialize with first price
    result.push(prices[0]);

    // Recursive formula
    for i in 1..prices.len() {
        let prev_ema = result[i - 1];
        let current_ema = alpha * prices[i] + (1.0 - alpha) * prev_ema;
        result.push(current_ema);
    }

    result
}
\end{lstlisting}

\subsection{SMA vs EMA}

\begin{center}
\begin{tabular}{lcc}
\toprule
\textbf{Property} & \textbf{SMA} & \textbf{EMA} \\
\midrule
Weighting & Equal & Exponential \\
Responsiveness & Slower & Faster \\
Lag & More & Less \\
Smoothness & Smoother & Noisier \\
Computation & Simple & Recursive \\
\bottomrule
\end{tabular}
\end{center}

\begin{keyinsight}
SMA is a \emph{lagging} indicator---it tells you what already happened. EMA
reacts faster but can give false signals. Traders often use both: short-term
EMA for entries, long-term SMA for trend confirmation.
\end{keyinsight}

\section{TimeSeries Trait}

Following our trait-based architecture, we define a generic interface:

\begin{lstlisting}[language=Rust]
pub trait TimeSeries {
    fn len(&self) -> usize;
    fn is_empty(&self) -> bool;
    fn closes(&self) -> Vec<f64>;
    fn volumes(&self) -> Vec<u64>;
}

impl TimeSeries for Vec<Ohlcv> {
    fn len(&self) -> usize { self.len() }
    fn is_empty(&self) -> bool { self.is_empty() }

    fn closes(&self) -> Vec<f64> {
        self.iter().map(|o| o.close).collect()
    }

    fn volumes(&self) -> Vec<u64> {
        self.iter().map(|o| o.volume).collect()
    }
}
\end{lstlisting}

\marginnote{%
\textbf{Trait Benefits}\\[0.3em]
The \texttt{TimeSeries} trait allows functions to work with any data source
that can provide price/volume sequences.
}

This enables generic functions that work with any \texttt{TimeSeries}:

\begin{lstlisting}[language=Rust]
fn analyze<T: TimeSeries>(data: &T) -> Report {
    let closes = data.closes();
    let sma_20 = sma(&closes, 20);
    let sma_50 = sma(&closes, 50);
    // ...
}
\end{lstlisting}

\section{Golden and Death Crosses}

A classic trading signal using moving averages:

\begin{itemize}
    \item \textbf{Golden Cross}: Short-term MA crosses \emph{above} long-term MA
      (bullish signal)
    \item \textbf{Death Cross}: Short-term MA crosses \emph{below} long-term MA
      (bearish signal)
\end{itemize}

\begin{lstlisting}[language=Rust]
fn find_crossovers(short_ma: &[f64], long_ma: &[f64]) -> Vec<(usize, bool)> {
    let mut crossovers = Vec::new();

    for i in 1..short_ma.len().min(long_ma.len()) {
        let prev_above = short_ma[i-1] > long_ma[i-1];
        let curr_above = short_ma[i] > long_ma[i];

        if !prev_above && curr_above {
            crossovers.push((i, true));  // Golden cross
        } else if prev_above && !curr_above {
            crossovers.push((i, false)); // Death cross
        }
    }

    crossovers
}
\end{lstlisting}

\section{Results and Analysis}

\subsection{Example Output}

\begin{lstlisting}[language=bash]
$ cargo run -p qf-03-stocks --example stock_analysis
Loading OHLCV data from data/spy.csv...
Loaded 252 trading days

Statistics:
  Average Volume: 85,432,100
  Price Change: +12.5%
  Highest High: $485.32
  Lowest Low: $410.15
  Average Daily Range: $5.23

Moving Averages:
  SMA(20) latest: $472.15
  SMA(50) latest: $465.80
  EMA(20) latest: $474.22

Crossovers (last 12 months):
  Day 45: Golden Cross (bullish)
  Day 128: Death Cross (bearish)
  Day 201: Golden Cross (bullish)
\end{lstlisting}

\section{Exercises}

\begin{enumerate}
    \item \textbf{Candlestick Patterns}: Implement detection for ``doji''
      (open $\approx$ close) and ``hammer'' (small body, long lower shadow).

    \item \textbf{Volume Analysis}: Add \func{volume\_ma} that computes moving
      average of volume. Flag days where volume exceeds 2x the average.

    \item \textbf{Bollinger Bands}: Implement Bollinger Bands: SMA $\pm$ 2
      standard deviations. This requires computing rolling standard deviation.

    \item \textbf{MACD}: Implement the MACD indicator: EMA(12) - EMA(26), with
      a 9-period signal line.

    \item \textbf{Backtest}: Using SMA crossover signals, simulate a simple
      strategy: buy on golden cross, sell on death cross. Track P\&L.
\end{enumerate}

\section{Key Takeaways}

\begin{itemize}
    \item \textbf{OHLCV} is the atomic unit of market data
    \item \textbf{Candlesticks} encode price action: body shows direction,
      shadows show rejection
    \item \textbf{SMA} averages equally; \textbf{EMA} weights recent prices
    \item \textbf{Crossovers} generate trading signals (but beware of lag!)
    \item \textbf{TimeSeries trait} enables generic, reusable analysis code
\end{itemize}

\vspace{1em}
\noindent\textbf{Next chapter}: Returns and volatility---the foundation of risk
measurement and portfolio theory.


%----------------------------------------------------------------------------------------
% PART II: CORE CONCEPTS (Intermediate)
%----------------------------------------------------------------------------------------

\pagelayout{wide}
\addpart{Core Concepts}
\pagelayout{margin}

\chapter{Returns and Volatility}
\label{ch:returns}
% \chapter{Returns and Volatility}
\label{ch:returns}

\begin{companioncode}{qf-04-returns}
Code: \texttt{crates/qf-04-returns/}\\
Dependencies: \crate{qf-common}, \crate{qf-03-stocks}\\
Run: \texttt{cargo run -p qf-04-returns --example returns\_analysis}
\end{companioncode}

\section{Learning Objectives}

This chapter introduces the foundational vocabulary of quantitative finance.
After completing it you can:

\begin{itemize}
    \item Compute simple and logarithmic returns and know when to use each
    \item Measure volatility and annualise it for cross-frequency comparison
    \item Evaluate risk-adjusted performance with the Sharpe and Sortino ratios
    \item Quantify drawdowns and recovery time from a price series
    \item Use the \texttt{Returns} trait to abstract return computation across
      data sources
\end{itemize}

\begin{keyinsight}
Returns, not prices, are the currency of quant research. Prices are
non-stationary; returns are (approximately) stationary. Every quant model in
this book starts by transforming prices into returns.
\end{keyinsight}

\section{Simple vs Log Returns}

\subsection{Definitions}

\marginnote{%
\textbf{Simple return}\\[0.3em]
$r_t = \dfrac{P_t - P_{t-1}}{P_{t-1}}$\\[0.5em]
\textbf{Log return}\\[0.3em]
$r_t = \ln\!\left(\dfrac{P_t}{P_{t-1}}\right)$
}

The simple (arithmetic) return is the percentage change in price:
\[
r_t = \frac{P_t - P_{t-1}}{P_{t-1}}
\]
The log (continuously compounded) return is the natural log of the price ratio:
\[
r_t = \ln\!\left(\frac{P_t}{P_{t-1}}\right) = \ln(P_t) - \ln(P_{t-1})
\]

The two are related exactly by
\[
r^{\text{log}}_t = \ln\!\left(1 + r^{\text{simple}}_t\right),
\]
and for small returns ($|r| \ll 1$) they are nearly equal.

\subsection{When to use which}

\begin{center}
\begin{tabular}{lcc}
\toprule
\textbf{Property} & \textbf{Simple} & \textbf{Log} \\
\midrule
Interpretation & ``I made 10\%'' & Less intuitive \\
Additivity over time & No & Yes \\
Additivity over assets & Yes & No \\
Symmetry & Symmetric & Unbounded below \\
Compounding & Multiplicative & Additive \\
\bottomrule
\end{tabular}
\end{center}

\begin{keyinsight}
Log returns are additive across time: the $k$-period log return is the sum of
the $k$ single-period log returns. This makes them natural for time-series
modelling. Simple returns are additive across assets in a portfolio (weighted
by portfolio weights), which makes them natural for portfolio returns.
\end{keyinsight}

\subsection{Rust implementation}

\begin{lstlisting}[language=Rust]
/// Simple return: (P_t - P_{t-1}) / P_{t-1}
pub fn simple_returns(prices: &[f64]) -> Vec<f64> {
    if prices.len() < 2 { return Vec::new(); }
    let mut out = Vec::with_capacity(prices.len() - 1);
    for i in 1..prices.len() {
        let prev = prices[i - 1];
        out.push(if prev == 0.0 { 0.0 }
                 else { (prices[i] - prev) / prev });
    }
    out
}

/// Log return: ln(P_t / P_{t-1})
pub fn log_returns(prices: &[f64]) -> Vec<f64> {
    if prices.len() < 2 { return Vec::new(); }
    let mut out = Vec::with_capacity(prices.len() - 1);
    for i in 1..prices.len() {
        let (prev, curr) = (prices[i - 1], prices[i]);
        out.push(if prev > 0.0 && curr > 0.0 { (curr / prev).ln() }
                 else { 0.0 });
    }
    out
}
\end{lstlisting}

\marginnote{%
\textbf{Cumulative return}\\[0.3em]
$R_{0\to t} = \prod_{k=1}^{t}(1+r_k) - 1$
}

The cumulative return over a period is the running compounding of
single-period simple returns:
\[
R_{0 \to t} = \prod_{k=1}^{t}(1 + r_k) - 1.
\]

\begin{lstlisting}[language=Rust]
pub fn cumulative_returns(returns: &[f64]) -> Vec<f64> {
    let mut out = Vec::with_capacity(returns.len());
    let mut wealth = 1.0_f64;
    for &r in returns {
        wealth *= 1.0 + r;
        out.push(wealth - 1.0);
    }
    out
}
\end{lstlisting}

\section{Volatility}

\marginnote{%
\textbf{Sample std dev}\\[0.3em]
$\sigma = \sqrt{\dfrac{1}{n-1}\sum_{t=1}^{n}(r_t - \bar r)^2}$
}

Volatility is the standard deviation of returns. We use the \emph{sample}
estimator (denominator $n-1$), consistent with \func{qf\_common::compute\_stats}:
\[
\sigma = \sqrt{\frac{1}{n-1}\sum_{t=1}^{n}(r_t - \bar r)^2}.
\]

\subsection{Annualisation}

\marginnote{%
\textbf{Annualised vol}\\[0.3em]
$\sigma_{\text{ann}} = \sigma_{\text{period}} \cdot \sqrt{m}$\\[0.3em]
where $m$ is the number of periods per year (252 for daily).
}

To compare volatilities across frequencies (daily, weekly, monthly), we
annualise. Under the assumption that returns are i.i.d., variance scales
linearly with the horizon, so standard deviation scales with the square root:
\[
\sigma_{\text{annual}} = \sigma_{\text{period}} \cdot \sqrt{m}.
\]

\begin{lstlisting}[language=Rust]
pub fn volatility(returns: &[f64]) -> f64 {
    let n = returns.len();
    if n < 2 { return 0.0; }
    let mean: f64 = returns.iter().sum::<f64>() / n as f64;
    let ssd: f64 = returns.iter()
        .map(|&r| { let d = r - mean; d * d })
        .sum();
    (ssd / (n - 1) as f64).sqrt()
}

pub fn annualized_volatility(returns: &[f64], periods_per_year: f64) -> f64 {
    if periods_per_year <= 0.0 { return 0.0; }
    volatility(returns) * periods_per_year.sqrt()
}
\end{lstlisting}

\begin{warning}
The $\sqrt{m}$ rule assumes returns are independent and identically
distributed. For assets with autocorrelated or volatility-clustered returns
(see Chapter~\ref{ch:volatility}), annualisation by $\sqrt{m}$ overstates the
true annual volatility.
\end{warning}

\subsection{Rolling volatility}

A rolling window reveals how volatility changes over time---useful for
regime detection and risk monitoring:

\begin{lstlisting}[language=Rust]
pub fn rolling_volatility(returns: &[f64], window: usize) -> Vec<f64> {
    if window == 0 || window > returns.len() { return Vec::new(); }
    let mut out = Vec::with_capacity(returns.len() - window + 1);
    for i in window - 1..returns.len() {
        let slice = &returns[i + 1 - window..=i];
        out.push(volatility(slice));
    }
    out
}
\end{lstlisting}

\section{Risk-Adjusted Performance}

\subsection{Sharpe ratio}

\marginnote{%
\textbf{Sharpe ratio}\\[0.3em]
$S = \dfrac{\E[r] - r_f}{\sigma(r)}$
}

The Sharpe ratio measures excess return per unit of total risk:
\[
S = \frac{\bar r - r_f}{\sigma(r)}.
\]
To annualise, multiply by $\sqrt{m}$ (because both the mean excess return and
the standard deviation scale together under i.i.d.):
\[
S_{\text{annual}} = S_{\text{period}} \cdot \sqrt{m}.
\]

\begin{lstlisting}[language=Rust]
pub fn sharpe_ratio(returns: &[f64], risk_free_rate: f64) -> f64 {
    let n = returns.len();
    if n < 2 { return 0.0; }
    let vol = volatility(returns);
    if vol == 0.0 { return 0.0; }
    let mean: f64 = returns.iter().sum::<f64>() / n as f64;
    (mean - risk_free_rate) / vol
}

pub fn annualized_sharpe(returns: &[f64], rf: f64, periods_per_year: f64) -> f64 {
    if periods_per_year <= 0.0 { return 0.0; }
    sharpe_ratio(returns, rf) * periods_per_year.sqrt()
}
\end{lstlisting}

\subsection{Sortino ratio}

\marginnote{%
\textbf{Downside deviation}\\[0.3em]
$\sigma_D = \sqrt{\dfrac{1}{n}\sum_{r_t < r_f}(r_f - r_t)^2}$
}

The Sortino ratio replaces total volatility with \emph{downside} volatility
only. Upside volatility is good---it should not penalise the ratio:
\[
\text{Sortino} = \frac{\bar r - r_f}{\sigma_D}, \qquad
\sigma_D = \sqrt{\frac{1}{n}\sum_{r_t < r_f}(r_f - r_t)^2}.
\]

\begin{lstlisting}[language=Rust]
pub fn sortino_ratio(returns: &[f64], risk_free_rate: f64) -> f64 {
    let n = returns.len();
    if n < 2 { return 0.0; }
    let downside_sq: f64 = returns.iter()
        .map(|&r| {
            let short = risk_free_rate - r;
            if short > 0.0 { short * short } else { 0.0 }
        })
        .sum();
    if downside_sq == 0.0 { return 0.0; }
    let downside_dev = (downside_sq / n as f64).sqrt();
    let mean: f64 = returns.iter().sum::<f64>() / n as f64;
    (mean - risk_free_rate) / downside_dev
}
\end{lstlisting}

\begin{keyinsight}
Sortino $>$ Sharpe whenever the return distribution is right-skewed (more
upside than downside). For symmetric distributions the two ratios are equal in
expectation. Sortino is preferred for evaluating hedge-fund-style strategies
where upside volatility is desirable.
\end{keyinsight}

\section{Drawdown}

A drawdown is the percentage decline from a running peak:
\[
D_t = \frac{P_t - \max_{k \le t} P_k}{\max_{k \le t} P_k} \le 0.
\]

\marginnote{%
\textbf{Max drawdown}\\[0.3em]
$\text{MaxDD} = \min_t D_t$\\[0.3em]
The worst peak-to-trough decline over the period.
}

\begin{lstlisting}[language=Rust]
pub fn drawdown(prices: &[f64]) -> Vec<f64> {
    let mut out = Vec::with_capacity(prices.len());
    let mut running_max = f64::NEG_INFINITY;
    for &p in prices {
        if p > running_max { running_max = p; }
        out.push(if running_max == 0.0 { 0.0 }
                 else { (p - running_max) / running_max });
    }
    out
}

pub fn max_drawdown(prices: &[f64]) -> f64 {
    drawdown(prices).iter().copied().fold(0.0, f64::min)
}
\end{lstlisting}

The \texttt{DrawdownStats} struct summarises the worst episode:
\begin{itemize}
    \item \texttt{max\_drawdown}: the most negative drawdown value
    \item \texttt{max\_drawdown\_duration}: longest run of consecutive
      negative-drawdown observations
    \item \texttt{recovery\_time}: steps from the trough back to a new high
      (\texttt{None} if the series never recovered)
\end{itemize}

\begin{keyinsight}
Max drawdown measures \emph{path risk}, not just terminal risk. A strategy
that returns 10\% with a 50\% drawdown along the way is very different from one
that returns 10\% smoothly. Most investors abandon strategies with deep
drawdowns long before recovery---a risk that terminal-return statistics
completely miss.
\end{keyinsight}

\section{The Returns Trait}

\marginnote{%
\textbf{Trait benefits}\\[0.3em]
Define the capability (\emph{``can produce returns''}) once; implement for
slices, owned vectors, and OHLCV data.
}

Following the workspace-wide trait-first architecture rule, we abstract return
computation over price sources:

\begin{lstlisting}[language=Rust]
pub trait Returns {
    fn simple_returns(&self) -> Vec<f64>;
    fn log_returns(&self) -> Vec<f64>;
}

impl Returns for &[f64] { /* delegate to free functions */ }
impl Returns for Vec<f64> { /* delegate to the slice impl */ }
impl Returns for Vec<Ohlcv> {
    fn simple_returns(&self) -> Vec<f64> {
        let closes: Vec<f64> = self.iter().map(|o| o.close).collect();
        simple_returns(&closes)
    }
    // log_returns similarly uses closing prices
}
\end{lstlisting}

This lets generic code consume returns from any source that implements
\texttt{Returns}, keeping analytics decoupled from data loading.

\section{Example Output}

\begin{lstlisting}[language=bash]
$ cargo run -p qf-04-returns --example returns_analysis
Returns Analysis
================

Simple Returns: mean=0.0031, std=0.0215
Log Returns: mean=0.0029, std=0.0214
Volatility (daily): 0.0215
Volatility (annualized): 0.3417
Sharpe Ratio (annualized): 2.29
Sortino Ratio: 0.22
Max Drawdown: -7.50%
Drawdown duration: 6 steps, recovery: Some(6)

Trait check (Vec<f64>): first simple return = 0.0030
\end{lstlisting}

\section{Exercises}

\begin{enumerate}
    \item \textbf{Return conversion}: Write a function that converts a series
      of simple returns to log returns and vice versa. Verify the identity
      $r^{\text{log}} = \ln(1 + r^{\text{simple}})$ with a property test.

    \item \textbf{Annualisation sanity check}: For daily, weekly, and monthly
      returns of the same asset, verify that the annualised volatilities are
      approximately equal (they differ only by sampling noise).

    \item \textbf{Calmar ratio}: Implement the Calmar ratio: annualised return
      divided by the absolute value of the maximum drawdown. Compare it with
      the Sharpe ratio on a synthetic strategy.

    \item \textbf{Rolling Sharpe}: Implement a rolling Sharpe ratio (windowed
      mean and windowed std) and visualise how it changes over time. Where
      does it dip? Why?

    \item \textbf{Drawdown with recovery}: Extend \texttt{drawdown\_stats} to
      report the average drawdown (mean of negative drawdown values) and the
      proportion of time spent in drawdown.
\end{enumerate}

\section{Key Takeaways}

\begin{itemize}
    \item \textbf{Log returns} are time-additive; \textbf{simple returns} are
      portfolio-additive. Choose accordingly.
    \item \textbf{Volatility} uses sample std ($n-1$); annualise by multiplying
      by $\sqrt{m}$, valid under i.i.d.
    \item \textbf{Sharpe} penalises all volatility; \textbf{Sortino} penalises
      only downside.
    \item \textbf{Max drawdown} measures path risk---the deepest hole on the
      equity curve.
    \item \textbf{The \texttt{Returns} trait} keeps analytics decoupled from
      data sources, following the workspace trait-first rule.
\end{itemize}

\vspace{1em}
\noindent\textbf{Next chapter}: Your first backtest---an SMA crossover
strategy with P\&L accounting, transaction costs, and drawdown-based risk
evaluation.
\textit{Chapter content pending --- Phase 4 implementation.}

\chapter{Your First Backtest}
\label{ch:backtest-intro}
% \chapter{Backtesting Strategies}
\label{ch:backtest}

\begin{companioncode}{qf-05-backtest}
Code: \texttt{crates/qf-05-backtest/}\\
Dependencies: \crate{qf-common}, \crate{qf-03-stocks}, \crate{qf-04-returns}\\
Run: \texttt{cargo run -p qf-05-backtest --example backtest\_demo}
\end{companioncode}

\section{Learning Objectives}

This chapter introduces the mechanics of evaluating a trading strategy on
historical data. After completing it you can:

\begin{itemize}
    \item Explain what a backtest is --- and what it is not
    \item Generate trading signals from indicators with no look-ahead bias
    \item Account for transaction costs in P\&L
    \item Evaluate a strategy with the risk-adjusted metrics from
      \cref{ch:returns}
    \item Avoid the three classic backtest traps: overfitting, look-ahead
      bias, and survivorship bias
\end{itemize}

\begin{keyinsight}
A backtest is a \emph{simulation}, not a prediction. It tells you how a
strategy \emph{would have} performed, not how it \emph{will} perform. The number
one failure mode is overfitting to historical data: a strategy that looks
great in backtest and loses money in production.
\end{keyinsight}

\section{What Is Backtesting?}

Backtesting is the replay of a trading strategy over historical price data,
bar by bar, to estimate its hypothetical performance. The engine maintains
state (position, cash, equity) and lets the strategy decide, at each bar,
what action to take given the price history it has seen so far.

\marginnote{%
\textbf{Simulation vs prediction}\\[0.3em]
A backtest does not predict. It is a \emph{counterfactual}: ``if you had
traded this strategy with these rules in this market, here is what would have
happened.'' Any conclusion about the future requires additional assumptions
(market stationarity, capacity, regime persistence) that the backtest itself
cannot validate.
}

Three assumptions are implicit and worth naming:

\begin{enumerate}
    \item \textbf{Market stationarity}: the statistical properties of the
      historical sample are representative of the future. Often false.
    \item \textbf{Capacity}: the strategy can take the simulated sizes without
      moving the market. True for retail; false for institutional size.
    \item \textbf{Execution}: the simulated fills are achievable. Slippage,
      partial fills, and market impact are real and cost money.
\end{enumerate}

\section{Signals and Positions}

\marginnote{%
\textbf{Signal vs position}\\[0.3em]
A signal is an \emph{instruction} (what to do this bar); a position is
a \emph{state} (what we are currently holding). The engine applies the
signal to update the position. Decoupling them lets you replay the
same signal stream under different execution models.
}

A \emph{signal} is a discrete instruction for the current bar. We use three:

\begin{lstlisting}[language=Rust]
#[derive(Debug, Clone, Copy, PartialEq, Eq)]
pub enum Signal {
    Buy,  // Enter long
    Sell, // Exit long (long-only for this chapter)
    Hold, // Do nothing
}
\end{lstlisting}

A \emph{position} is what the engine is currently holding:

\begin{lstlisting}[language=Rust]
#[derive(Debug, Clone, Copy, PartialEq, Eq)]
pub enum Position {
    Long,  // Holding the asset
    Short, // Reserved for later phases
    Flat,  // No position
}
\end{lstlisting}

\begin{keyinsight}
Why separate \texttt{Signal} and \texttt{Position}? A signal is an
\emph{intention}; a position is a \emph{state}. The engine translates
intentions into state transitions, charging transaction costs on every
change. This separation lets the strategy stay declarative (``I want to be
long now'') and the engine stay responsible for the messy realities
(costs, fill assumptions, position limits).
\end{keyinsight}

\section{The Strategy Trait}

\marginnote{%
\textbf{No look-ahead}\\[0.3em]
\texttt{signal(data, index)} receives \texttt{data[..=index]}. The strategy
must never read past the current bar. The engine enforces this by passing a
slice and the current index --- implementations physically cannot access
\texttt{data[index+1..]} without unsafe code.
}

The extension point of the engine is the \texttt{Strategy} trait:

\begin{lstlisting}[language=Rust]
pub trait Strategy {
    fn signal(&self, data: &[Ohlcv], index: usize) -> Signal;
    fn name(&self) -> &str;
}
\end{lstlisting}

The engine is generic over \texttt{S: Strategy}, so new strategies drop in
without engine changes. The trait takes the full slice plus the current index
so the strategy can compute any indicator that fits in the observed history.

\section{SMA Crossover}

The canonical first strategy: go long when a short moving average crosses
above a long moving average (the \emph{golden cross}); exit when it crosses
below (the \emph{death cross}).

\marginnote{%
\textbf{Golden / death cross}\\[0.3em]
$\text{SMA}_S$ crosses above $\text{SMA}_L$ $\Rightarrow$ Buy\\[0.3em]
$\text{SMA}_S$ crosses below $\text{SMA}_L$ $\Rightarrow$ Sell\\[0.5em]
A crossover is a \emph{change in the relationship}, not a level. The previous
bar must satisfy the opposite inequality.
}

\begin{lstlisting}[language=Rust]
pub struct SmaCrossover {
    short_period: usize,
    long_period: usize,
}

impl Strategy for SmaCrossover {
    fn signal(&self, data: &[Ohlcv], index: usize) -> Signal {
        if index < self.long_period { return Signal::Hold; }

        let short_now = match sma_at(data, index, self.short_period) {
            Some(v) => v, None => return Signal::Hold,
        };
        let long_now = match sma_at(data, index, self.long_period) {
            Some(v) => v, None => return Signal::Hold,
        };
        let short_prev = match sma_at(data, index - 1, self.short_period) {
            Some(v) => v, None => return Signal::Hold,
        };
        let long_prev = match sma_at(data, index - 1, self.long_period) {
            Some(v) => v, None => return Signal::Hold,
        };

        let crossed_up   = short_prev <= long_prev  && short_now >  long_now;
        let crossed_down = short_prev >= long_prev  && short_now <  long_now;

        if crossed_up   { Signal::Buy  }
        else if crossed_down { Signal::Sell }
        else { Signal::Hold }
    }
    fn name(&self) -> &str { "SMA Crossover" }
}
\end{lstlisting}

\begin{keyinsight}
Crossover strategies are \emph{trend followers}: they only make money when
prices trend. In choppy, mean-reverting markets they lose, because each
crossover is a false signal that gets reversed at the next crossover --- with
transaction costs paid on both sides.
\end{keyinsight}

\section{P\&L with Transaction Costs}

\marginnote{%
\textbf{Transaction cost}\\[0.3em]
At entry: $c_{\text{entry}} = \text{notional} \cdot \tau$\\[0.3em]
At exit: $c_{\text{exit}} = \text{shares} \cdot P_{\text{exit}} \cdot \tau$\\[0.5em]
Net P\&L: $\text{shares} \cdot (P_{\text{exit}} - P_{\text{entry}}) - c_{\text{entry}} - c_{\text{exit}}$
}

Every state change costs money. The engine charges a proportional cost
$\tau$ (e.g.\ $\tau = 0.001$ for 10 basis points) on the traded notional:

\begin{lstlisting}[language=Rust]
// Enter long: spend all cash at the close, paying the cost first.
let notional = capital;
let cost = notional * config.transaction_cost;
shares = (notional - cost) / price;
total_costs += cost;

// Exit long: sell all shares at the close.
let notional = shares * price;
let cost = notional * config.transaction_cost;
total_costs += cost;
capital = notional - cost;
\end{lstlisting}

At each bar, equity is marked to market: \texttt{shares * close} when long,
cash when flat. The per-bar simple returns of the equity curve become the
\texttt{daily\_returns} used by the risk metrics from \cref{ch:returns}.

\section{Performance Evaluation}

\marginnote{%
\textbf{Why reuse \crate{qf-04-returns}?}\\[0.3em]
Sharpe, Sortino, and max drawdown are return-level statistics --- they
do not care how the returns were generated. Once the backtest produces
a daily return series, evaluation is identical to computing the same
metrics on a buy-and-hold benchmark. Code reuse follows the math.
}

\texttt{BacktestResult} delegates to \texttt{qf-04-returns} for the four
headline metrics:

\begin{lstlisting}[language=Rust]
impl BacktestResult {
    pub fn sharpe(&self, rf: f64) -> f64 {
        qf_04_returns::annualized_sharpe(&self.daily_returns, rf, 252.0)
    }
    pub fn sortino(&self, rf: f64) -> f64 {
        qf_04_returns::sortino_ratio(&self.daily_returns, rf)
    }
    pub fn max_drawdown(&self) -> f64 {
        qf_04_returns::max_drawdown(&self.equity_curve)
    }
    pub fn volatility(&self) -> f64 {
        qf_04_returns::annualized_volatility(&self.daily_returns, 252.0)
    }
}
\end{lstlisting}

\begin{keyinsight}
The risk metrics operate on the \emph{equity curve}, not the raw price
series. A strategy that goes flat for half the period has a different risk
profile from buy-and-hold even on the same underlying, because the equity
curve is flat whenever the position is flat. Volatility and drawdown reflect
this exposure pattern automatically.
\end{keyinsight}

\section{Buy-and-Hold Benchmark}

\marginnote{%
\textbf{Always compare to buy-and-hold}\\[0.3em]
An active strategy that returns $10\%$ looks good until you realise
the underlying asset returned $15\%$ over the same period. The
buy-and-hold benchmark is the floor any active strategy must beat
\emph{after} costs. If it does not, the activity added no value.
}

Every active strategy must be compared against the simplest alternative: buy
on the first bar and hold forever. \texttt{BuyAndHold} is a two-line strategy:

\begin{lstlisting}[language=Rust]
impl Strategy for BuyAndHold {
    fn signal(&self, _data: &[Ohlcv], index: usize) -> Signal {
        if index == 0 { Signal::Buy } else { Signal::Hold }
    }
    fn name(&self) -> &str { "Buy and Hold" }
}
\end{lstlisting}

If an active strategy cannot beat buy-and-hold after costs, the active
component is adding noise and cost, not value. The benchmark also lets you
decompose returns: alpha (skill) vs beta (market exposure).

\section{Common Pitfalls}

\marginnote{%
\textbf{The two cardinal sins}\\[0.3em]
\emph{Overfitting} and \emph{look-ahead bias} are the two ways a
backtest lies to you. Overfitting makes the past look good; look-ahead
bias makes the future knowable. Both produce strategies that fail
miserably in live trading. Treat any backtest Sharpe $> 2$ with deep
suspicion until both are ruled out.
}

\begin{description}
    \item[Overfitting] Tuning the SMA periods $(S, L)$ to maximise in-sample
      Sharpe almost always produces a strategy that fails out-of-sample. Use
      walk-forward validation (deferred to a later phase) and be suspicious of
      parameters tuned on the same data you evaluate on.
    \item[Look-ahead bias] Using information not available at decision time.
      Common sources: using the close of bar $t$ to decide at bar $t$ (only
      valid if the trade executes at the close), indicator warm-up that reads
      the future, corporate actions applied with hindsight. Our
      \texttt{signal(data, index)} interface rules out most of these by
      construction.
    \item[Survivorship bias] Backtesting only on companies that survived to the
      present. Companies that went bankrupt are absent, inflating historical
      returns. Use point-in-time universes when available.
\end{description}

\section{Example Output}

\begin{lstlisting}[language=bash]
$ cargo run -p qf-05-backtest --example backtest_demo
==============================
Backtest Results
==============================

Period: 2024-001 to 2024-252

--- Strategy: SMA Crossover (10/30) ---
  Total Return:   3.49%
  Sharpe (ann.):   2.55
  Sortino:         0.31
  Max Drawdown:    -0.41%
  Trades:          1

--- Strategy: Buy and Hold ---
  Total Return:    4.54%
  Sharpe (ann.):   1.70
  Max Drawdown:    -6.35%

Outperformance vs Buy & Hold: -1.04%
\end{lstlisting}

The SMA crossover underperforms buy-and-hold in this synthetic regime: the
cost of false signals in choppy sections outweighs the benefit of avoiding
the mid-year drawdown. This is the typical result for trend-following
strategies in noisy markets, and it is the right lesson to take from a first
backtest.

\section{Walk-Forward Analysis}
\label{sec:walk-forward}

\marginnote{%
\textbf{Walk-forward vs k-fold}\\[0.3em]
Unlike k-fold CV (which can use future data in some folds),
walk-forward strictly respects time ordering. Each window trains
on past data and tests on subsequent data --- mimicking real trading.
}

Walk-forward analysis addresses the overfitting trap by enforcing a strict
temporal separation between training and testing windows. Instead of tuning
parameters on the full dataset and evaluating on the same data, walk-forward
splits the time series into sequential windows:

\begin{enumerate}
    \item \textbf{In-sample window}: Historical data used for parameter tuning
    \item \textbf{Out-of-sample window}: Future data (not seen during tuning) for evaluation
    \item \textbf{Rolling forward}: Advance both windows and repeat
\end{enumerate}

\subsection{Walk-Forward Efficiency}

\marginnote{%
\textbf{Interpreting WFE}\\[0.3em]
$\text{WFE} = 1.0$: Strategy performs equally well in-sample and out-of-sample (ideal)\\[0.3em]
$\text{WFE} > 1.0$: Strategy performs \emph{better} out-of-sample (rare, check for bugs)\\[0.3em]
$\text{WFE} \approx 0.5$: Typical degradation for robust strategies\\[0.3em]
$\text{WFE} < 0.3$: Severe overfitting --- strategy likely fails in live trading
}

Walk-Forward Efficiency (WFE) quantifies how much performance degrades from
in-sample to out-of-sample:

\begin{equation}
\text{WFE} = \frac{\text{Sharpe}_{\text{out-of-sample}}}{\text{Sharpe}_{\text{in-sample}}}
\end{equation}

A WFE close to 1.0 indicates a robust strategy; values below 0.5 suggest
severe overfitting. The ratio measures how much ``alpha'' survived the
transition from tuning to live-like conditions.

\subsection{Implementation}

The \texttt{quant-backtest} crate provides walk-forward support via the
\texttt{WalkForward} cross-validator (see \cref{ch:afml} for
comprehensive AFML pipeline integration):

\begin{lstlisting}[language=Rust]
use quant_backtest::{WalkForward, WalkForwardConfig};
use quant_core::CrossValidator;

// Configure rolling 60-bar training, 20-bar testing, step 20 bars
let config = WalkForwardConfig::rolling(60, 20, 20);
let wf = WalkForward::new(config);

// Generate splits for 200 bars
let splits = wf.splits(0, 200);
// Returns: [ [0-60,60-80], [20-80,80-100], [40-100,100-120], ... ]
\end{lstlisting}

For each split: tune parameters on the in-sample window, freeze them, evaluate
on out-of-sample. Aggregate the out-of-sample Sharpe ratios and compute WFE.

\subsection{Visual Representation}

\begin{figure}[htbp]
\centering
\begin{tikzpicture}[>=stealth, thick]
  % Timeline
  \draw[->] (0,0) -- (12,0) node[right] {Time};

  % Window 1
  \draw[fill=blue!20] (0.5,0.3) rectangle (3.5,0.7) node[midway] {\small Train};
  \draw[fill=green!20] (3.5,0.3) rectangle (5,0.7) node[midway] {\small Test};

  % Window 2
  \draw[fill=blue!20] (2.5,1.0) rectangle (5.5,1.4) node[midway] {\small Train};
  \draw[fill=green!20] (5.5,1.0) rectangle (7,1.4) node[midway] {\small Test};

  % Window 3
  \draw[fill=blue!20] (4.5,1.7) rectangle (7.5,2.1) node[midway] {\small Train};
  \draw[fill=green!20] (7.5,1.7) rectangle (9,2.1) node[midway] {\small Test};

  % Labels
  \node[below] at (2,0) {\small $t_0$};
  \node[below] at (6,0) {\small $t_1$};
  \node[below] at (10,0) {\small $t_2$};

  % Legend
  \node[right] at (9.5,1.7) {\small Rolling windows};
  \draw[->] (9,1.5) -- (7,1.9);
\end{tikzpicture}
\caption{Walk-forward rolling windows: each training window slides forward,
  testing on the next unseen period. Unlike k-fold, every test window is
  strictly after its training window.}
\label{fig:walk-forward}
\end{figure}

\subsection{Anchored vs Rolling}

\marginnote{%
\textbf{Trade-off}\\[0.3em]
\emph{Anchored}: More data, but older data may be stale\\[0.3em]
\emph{Rolling}: Fresher data, but smaller sample size
}

\begin{itemize}
    \item \textbf{Rolling window}: Training window slides (fixed size, moves forward)
      \begin{itemize}
        \item Advantages: Uses only recent data (regime-aware)
        \item Disadvantages: Smaller training sets, higher variance
      \end{itemize}

    \item \textbf{Anchored window}: Training window expands (grows over time)
      \begin{itemize}
        \item Advantages: Larger training sets, lower variance
        \item Disadvantages: Older data may represent outdated regimes
      \end{itemize}
\end{itemize}

Choose rolling for high-frequency or regime-sensitive strategies; choose
anchored for slower, more stable strategies.

\begin{keyinsight}
Walk-forward analysis is the \emph{minimum viable safeguard} against
overfitting. It does not eliminate overfitting (you can still overfit
by cherry-picking the best split or tuning on WFE itself), but it forces
the strategy to prove it can perform on unseen data. A strategy with
$\text{WFE} > 0.7$ is far more credible than one evaluated only in-sample.
\end{keyinsight}

For advanced walk-forward workflows integrated with AFML labeling, purged
k-fold, and bet sizing, see \cref{ch:afml}.

\section{Exercises}

\begin{enumerate}
    \item \textbf{Cost sensitivity}: Run the 10/30 crossover on the synthetic
      series for $\tau \in \{0, 0.001, 0.0025, 0.005, 0.01\}$. Plot total
      return vs $\tau$. At what cost does the strategy stop being profitable?

    \item \textbf{Parameter sensitivity}: For $\tau = 0.001$, sweep
      $(S, L) \in \{5, 10, 20\} \times \{30, 50, 100\}$ with $S < L$. Which
      configuration maximises Sharpe? Is it the same one that maximises total
      return? Why or why not?

    \item \textbf{EMA crossover}: Implement an EMA crossover strategy by
      replacing SMA with \func{qf\_03\_stocks::ema}. Compare its turnover
      (number of trades) and Sharpe against the SMA crossover on the same
      data.

    \item \textbf{Position sizing}: Modify the engine to invest a fixed
      fraction $f \in (0, 1]$ of capital on each entry instead of 100\%.
      How does the equity curve change for $f = 0.5$? What is the trade-off?

    \item \textbf{Walk-forward (conceptual)}: Split the 252-bar series into
      an in-sample window (first 180 bars) and an out-of-sample window (last
      72 bars). Tune $(S, L)$ on the in-sample, freeze it, and evaluate on
      out-of-sample. How much does Sharpe degrade from in-sample to
      out-of-sample? This is the simplest walk-forward analysis.
\end{enumerate}

\section{Key Takeaways}

\begin{itemize}
    \item \textbf{Backtesting} is a simulation, not a prediction---it tells you
      how a strategy \emph{would have} performed
    \item \textbf{Signals vs positions}: signals are intentions, positions are
      state; the engine translates one to the other
    \item \textbf{Transaction costs} erode returns and must be modelled
      explicitly
    \item \textbf{SMA crossover} is a trend-following strategy that loses in
      choppy markets
    \item \textbf{Buy-and-hold} is the benchmark every active strategy must beat
    \item \textbf{Pitfalls}: overfitting, look-ahead bias, and survivorship bias
      are the classic backtest traps
\end{itemize}

\vspace{1em}
\noindent\textbf{Next chapter}: Foundations---moments, hand-rolled RNG, and
geometric Brownian motion simulation.
\textit{Chapter content pending --- Phase 5 implementation.}

%----------------------------------------------------------------------------------------
% PART III: QUANTITATIVE METHODS (Advanced)
%----------------------------------------------------------------------------------------

\pagelayout{wide}
\addpart{Quantitative Methods}
\pagelayout{margin}

\chapter{Foundations: Moments and Fat Tails}
\label{ch:moments}
% \section{Learning Objectives}

This chapter begins the advanced track. From here on, all math is hand-rolled —
no \texttt{rand}, \texttt{nalgebra}, or \texttt{statrs}. After completing it
you can:

\begin{itemize}
    \item Build a validated price series that rejects non-finite and
      non-positive values at construction time
    \item Compute the first four sample moments (mean, variance, skewness,
      excess kurtosis) from first principles
    \item Apply any aggregation over a sliding window with a generic
      \texttt{rolling} function
    \item Implement a deterministic PRNG (xorshift64\*) and a normal sampler
      (Box-Muller) without external crates
    \item Simulate geometric Brownian motion paths and understand the exact
      solution to the GBM SDE
    \item Detect fat tails via excess kurtosis and explain why Gaussian risk
      models underestimate tail risk
\end{itemize}

\begin{keyinsight}
The Gaussian distribution has excess kurtosis $0$. Real financial returns
have positive excess kurtosis: large losses and gains happen more often than
a normal model predicts. Any risk model that assumes Gaussian returns will
underestimate the probability of extreme events. The 2008 crisis was in
part a fat-tail crisis.
\end{keyinsight}

\section{A Validated Price Series}

\marginnote{%
\textbf{The invariant}\\[0.3em]
Every element of a \texttt{PriceSeries} is strictly positive and finite. The
constructor rejects the rest.
}

The foundation of every quant computation is a clean price series. We wrap
\texttt{Vec<f64>} in a newtype that enforces the invariant ``strictly positive
and finite'' at construction time:

\begin{lstlisting}[language=Rust]
#[derive(Debug, Clone, PartialEq)]
pub struct PriceSeries(Vec<f64>);

impl PriceSeries {
    pub fn new(prices: Vec<f64>) -> Result<Self, CoreError> {
        for &p in &prices {
            if !p.is_finite() || p <= 0.0 {
                return Err(CoreError::InvalidPrice);
            }
        }
        Ok(Self(prices))
    }
    pub fn as_slice(&self) -> &[f64] { &self.0 }
    pub fn len(&self) -> usize { self.0.len() }
}
\end{lstlisting}

Rejecting \texttt{NaN} and \texttt{+inf} at the boundary means every downstream
function can assume finite inputs and skip defensive checks. This is the
``parse, don't validate'' pattern applied to numeric data.

\section{Returns, Again}

Simple and log returns were introduced in \cref{ch:returns}. We re-implement
them here so \texttt{quant-core} is self-contained:

\begin{lstlisting}[language=Rust]
pub fn simple_returns(prices: &[f64]) -> Vec<f64> {
    if prices.len() < 2 { return Vec::new(); }
    let mut out = Vec::with_capacity(prices.len() - 1);
    for i in 1..prices.len() {
        let prev = prices[i - 1];
        out.push(if prev == 0.0 { 0.0 }
                 else { (prices[i] - prev) / prev });
    }
    out
}
\end{lstlisting}

\begin{keyinsight}
The identity $\ln(1 + r^{\text{simple}}) = r^{\text{log}}$ holds exactly. For
small returns the two are nearly equal; for large returns they diverge. This
identity is a cheap property test that catches implementation bugs in either
function.
\end{keyinsight}

\section{Moments}

\marginnote{%
\textbf{Central moment}\\[0.3em]
$m_k = \dfrac{1}{n}\sum_{i=1}^{n}(x_i - \bar x)^k$\\[0.5em]
\textbf{Skewness}\\[0.3em]
$g_1 = \dfrac{m_3}{m_2^{3/2}}$\\[0.3em]
\textbf{Excess kurtosis}\\[0.3em]
$g_2 = \dfrac{m_4}{m_2^2} - 3$
}

The $k$-th central moment is
\[
m_k = \frac{1}{n}\sum_{i=1}^{n}(x_i - \bar x)^k.
\]
Skewness and excess kurtosis are standardised ratios of these moments:

\begin{align*}
g_1 &= \frac{m_3}{m_2^{3/2}} \quad \text{(skewness)} \\
g_2 &= \frac{m_4}{m_2^2} - 3 \quad \text{(excess kurtosis)}
\end{align*}

The Gaussian distribution has $g_1 = 0$ and $g_2 = 0$. Positive skewness means
a longer right tail; negative skewness a longer left tail. Positive excess
kurtosis means heavier tails than Gaussian.

\begin{lstlisting}[language=Rust]
pub fn skewness(data: &[f64]) -> Result<f64, CoreError> {
    if data.len() < 3 { return Err(CoreError::InsufficientData { .. }); }
    let m = mean(data);
    let m2 = central_moment_2(data, m);
    let m3 = central_moment_3(data, m);
    Ok(m3 / m2.powf(1.5))
}

pub fn excess_kurtosis(data: &[f64]) -> Result<f64, CoreError> {
    if data.len() < 4 { return Err(CoreError::InsufficientData { .. }); }
    let m = mean(data);
    let m2 = central_moment_2(data, m);
    let m4 = central_moment_4(data, m);
    Ok(m4 / (m2 * m2) - 3.0)
}
\end{lstlisting}

\begin{keyinsight}
We use the sample variance ($n - 1$ denominator) for $\sigma^2$, consistent
with \func{qf\_common::compute\_stats}, but use population central moments
(denominator $n$) for $m_2, m_3, m_4$ in skewness and kurtosis. This is the
textbook definition; the difference between the two vanishes for large $n$.
\end{keyinsight}

The \texttt{Moments} trait abstracts over any data source that can produce
sample moments, so the same analytics work on slices, owned vectors, and
\texttt{PriceSeries}:

\begin{lstlisting}[language=Rust]
pub trait Moments {
    fn mean(&self) -> f64;
    fn variance(&self) -> Result<f64, CoreError>;
    fn std_dev(&self) -> Result<f64, CoreError>;
    fn skewness(&self) -> Result<f64, CoreError>;
    fn excess_kurtosis(&self) -> Result<f64, CoreError>;
}

impl Moments for &[f64] { /* delegate to free functions */ }
impl Moments for Vec<f64> { /* delegate to slice impl */ }
impl Moments for PriceSeries { /* delegate to slice impl */ }
\end{lstlisting}

\section{Rolling Windows}

\marginnote{%
\textbf{Rolling window}\\[0.3em]
For window $w$ on data of length $n$, the output has length $n - w + 1$ and
entry $i$ aggregates $\texttt{data}[i..i+w]$.
}

A rolling window applies an aggregation over each contiguous slice of length
$w$. The generic \texttt{rolling} function takes a closure, so any
aggregation (mean, std, max, custom) drops in:

\begin{lstlisting}[language=Rust]
pub fn rolling<F, T>(window: usize, data: &[f64], f: F)
    -> Result<Vec<T>, CoreError>
where F: Fn(&[f64]) -> T
{
    if window == 0 || window > data.len() {
        return Err(CoreError::InvalidWindow { window, len: data.len() });
    }
    let mut out = Vec::with_capacity(data.len() - window + 1);
    for i in 0..=(data.len() - window) {
        out.push(f(&data[i..i + window]));
    }
    Ok(out)
}

pub fn rolling_mean(window: usize, data: &[f64]) -> Result<Vec<f64>, CoreError> {
    rolling(window, data, mean)
}
\end{lstlisting}

\section{Hand-Rolled Simulation}

\subsection{xorshift64\*}

\marginnote{%
\textbf{xorshift64\*}\\[0.3em]
Three shifts and a multiply. Period $2^{64} - 1$. Fast, stateless, good
enough for Monte Carlo when cryptographic strength is not required.
}

We implement Vigna's xorshift64\* PRNG: three xor-shifts on a 64-bit state,
then a multiply by a constant. A seed of zero produces a degenerate stream,
so we replace it with a fixed default:

\begin{lstlisting}[language=Rust]
pub struct XorShift64 { state: u64 }

impl Rng for XorShift64 {
    fn next_u64(&mut self) -> u64 {
        let mut x = self.state;
        x ^= x >> 12;
        x ^= x << 25;
        x ^= x >> 27;
        self.state = x;
        x.wrapping_mul(0x2545_F491_4F6C_DD1D)
    }
    fn next_f64(&mut self) -> f64 {
        (self.next_u64() >> 11) as f64 / (1u64 << 53) as f64
    }
}
\end{lstlisting}

\subsection{Box-Muller Normal}

To sample $N(\mu, \sigma^2)$ without an external library, we use the
Box-Muller transform: two independent uniforms $u_1, u_2 \sim U(0,1)$ produce
one standard normal variate
\[
Z = \sqrt{-2 \ln u_1} \cos(2\pi u_2).
\]

\begin{lstlisting}[language=Rust]
pub fn box_muller_normal<R: Rng + ?Sized>(
    rng: &mut R, mean: f64, std: f64,
) -> f64 {
    let u1 = rng.next_f64().max(f64::EPSILON);
    let u2 = rng.next_f64();
    let r = (-2.0 * u1.ln()).sqrt();
    let theta = 2.0 * PI * u2;
    mean + std * (r * theta.cos())
}
\end{lstlisting}

The \texttt{Distribution} trait lets us plug in new distributions (exponential,
Student-$t$, jump-diffusion) without touching the simulator:

\begin{lstlisting}[language=Rust]
pub trait Distribution {
    fn sample<R: Rng + ?Sized>(&self, rng: &mut R) -> f64;
}

pub struct Normal { pub mean: f64, pub std: f64 }

impl Distribution for Normal {
    fn sample<R: Rng + ?Sized>(&self, rng: &mut R) -> f64 {
        box_muller_normal(rng, self.mean, self.std)
    }
}
\end{lstlisting}

\subsection{Geometric Brownian Motion}

\marginnote{%
\textbf{Exact GBM solution}\\[0.3em]
$S_{t+\Delta t} = S_t \exp\!\left(\left(\mu - \tfrac{1}{2}\sigma^2\right)\Delta t + \sigma\sqrt{\Delta t}\, Z\right)$\\[0.5em]
With $\sigma = 0$: $S_T = S_0 e^{\mu T}$, deterministic.
}

The stochastic differential equation
\[
dS_t = \mu S_t\, dt + \sigma S_t\, dW_t
\]
has the exact solution
\[
S_{t+\Delta t} = S_t \exp\!\left(\left(\mu - \tfrac{1}{2}\sigma^2\right)\Delta t + \sigma\sqrt{\Delta t}\, Z\right),
\]
where $Z \sim N(0, 1)$. We simulate one step at a time:

\begin{lstlisting}[language=Rust]
pub fn gbm_paths<R: Rng + ?Sized>(
    s0: f64, mu: f64, sigma: f64, t: f64,
    n_steps: usize, n_paths: usize, rng: &mut R,
) -> Vec<Vec<f64>> {
    let dt = t / n_steps as f64;
    let drift = (mu - 0.5 * sigma * sigma) * dt;
    let diffusion = sigma * dt.sqrt();
    let normal = Normal::standard();
    let mut paths = Vec::with_capacity(n_paths);
    for _ in 0..n_paths {
        let mut path = Vec::with_capacity(n_steps + 1);
        path.push(s0);
        let mut s = s0;
        for _ in 0..n_steps {
            let z = normal.sample(rng);
            s *= (drift + diffusion * z).exp();
            path.push(s);
        }
        paths.push(path);
    }
    paths
}
\end{lstlisting}

\begin{keyinsight}
With $\sigma = 0$, the $Z$ term drops out and $S_T = S_0 \exp(\mu T)$
deterministically. This is a useful sanity check: a zero-volatility GBM
should reproduce the risk-free growth exactly, with no RNG dependence.
\end{keyinsight}

\section{Fat Tails}

\begin{lstlisting}[language=bash]
$ cargo run -p quant-core --example fat_tails
1. Pure Gaussian (10k N(0,1) draws)
  skewness     = +0.0267
  excess kurt  = +0.0323

2. Fat tails (every 100th draw x10)
  skewness     = -0.3051
  excess kurt  = +74.1272

3. GBM terminal prices (10k paths, 252 steps)
  skewness     = +0.5957
  excess kurt  = +0.5959
\end{lstlisting}

The pure Gaussian sample has excess kurtosis near $0$, as expected. Scaling
every 100th draw by 10 injects rare large moves; the excess kurtosis jumps
from $\sim 0$ to $\sim 74$. GBM terminal prices are log-normally
distributed, so they have positive excess kurtosis (heavier right tail than
a Gaussian) and positive skewness --- a direct consequence of the
exponential in the exact solution.

\begin{keyinsight}
A Gaussian model fit to fat-tailed data will price a 4-sigma event as a
``1-in-15,000-years'' occurrence. The true frequency under fat tails may be
once a decade. This is why Value-at-Risk computed under Gaussian assumptions
systematically underestimates tail losses, and why option-implied volatilities
smile --- the market knows the tails are fatter than Gaussian allows.
\end{keyinsight}

\section{Exercises}

\begin{enumerate}
    \item \textbf{Sample vs population variance}: For
      $x \in \{1, 2, 3, 4, 5\}$, compute both the population variance
      (denominator $n$) and the sample variance (denominator $n - 1$). How
      large is the bias? When does it matter?

    \item \textbf{Skewness sign}: Construct a synthetic series with negative
      skew (long left tail) and one with positive skew. Verify
      \func{skewness} returns the correct sign. Which one is typical of
      equity returns?

    \item \textbf{RNG period}: The xorshift64\* generator has period
      $2^{64} - 1$. Write a test that verifies the first 1000 outputs from
      seed 42 are distinct, and reason about why ``distinct'' is necessary
      but not sufficient for a good PRNG.

    \item \textbf{GBM convergence}: Simulate 10\,000 GBM paths over one year
      (252 steps) with $S_0 = 100$, $\mu = 0.05$, $\sigma = 0.2$. Compute the
      mean terminal price. Compare with the analytical expectation
      $E[S_T] = S_0 e^{\mu T}$. How close is the Monte Carlo estimate?

    \item \textbf{Jump-diffusion}: Extend \func{gbm\_paths} with a Poisson
      jump term: with probability $\lambda \Delta t$ per step, multiply
      $S_t$ by $e^{J}$ where $J \sim N(\mu_J, \sigma_J^2)$. How does the
      excess kurtosis of terminal prices change as $\lambda$ grows?
\end{enumerate}
\textit{Chapter content pending --- Phase 6 implementation.}

\chapter{Time Series: Stationarity and Fractional Differentiation}
\label{ch:timeseries}
% \section{Learning Objectives}

Most financial time series are non-stationary: prices wander, volatilities
climb and fall, autocorrelations decay slowly. Naive regression on
non-stationary data produces spurious correlations and unstable forecasts.
After completing this chapter you can:

\begin{itemize}
    \item Fit an ordinary least squares regression by hand via Gaussian
      elimination with partial pivoting on the normal equations
    \item Compute standard errors, t-statistics, and $R^2$ without any
      external statistics crate
    \item Compute the sample autocorrelation function and read its structure
    \item Test for a unit root with the Augmented Dickey-Fuller test and
      interpret the MacKinnon critical value
    \item Apply López de Prado's fixed-width fractional differentiation to
      make a series stationary while preserving memory
    \item Find the minimum differencing order $d$ that achieves stationarity
      on real and synthetic price series
\end{itemize}

\begin{keyinsight}
Traditional differencing ($d = 1$) makes a price series stationary but erases
all long-range memory --- the resulting signal is essentially noise. Fractional
differentiation ($0 < d < 1$) achieves stationarity while keeping as much memory
as the data allow. The smallest such $d$ is the sweet spot for ML features that
predict future returns.
\end{keyinsight}

\section{Stationarity}

\marginnote{%
\textbf{Weak stationarity}\\[0.3em]
A process is weakly stationary if its mean is constant and its autocovariance
depends only on the lag, not on time. ``Stationary'' in this chapter means
weakly stationary.
}

A time series $\{y_t\}$ is \emph{weakly stationary} if its first two moments are
time-invariant: $E[y_t] = \mu$ and $\text{Cov}(y_t, y_{t-k}) = \gamma_k$ depend
only on the lag $k$. A random walk $y_t = y_{t-1} + \varepsilon_t$ fails this
definition because its variance grows without bound. Non-stationary data break
the assumptions behind OLS inference, forecast intervals, and most ML models.

The standard remedy is differencing: $\Delta y_t = y_t - y_{t-1}$. With
$d = 1$ the series becomes stationary but loses all memory of the distant
past. We need something in between.

\section{OLS via Gaussian Elimination}

\marginnote{%
\textbf{Normal equations}\\[0.3em]
$\hat\beta = (X'X)^{-1} X'y$. We never invert $X'X$ explicitly; we solve
$X'X \hat\beta = X'y$ by Gaussian elimination.
}

Given a design matrix $X$ (rows are observations, columns are features) and a
response $y$, the OLS estimator minimises $\|y - X\beta\|^2$. The
first-order condition gives the normal equations
\[
(X'X)\, \hat\beta = X'y.
\]
We solve this $k \times k$ linear system by Gaussian elimination with
partial pivoting: at each column we swap the row with the largest absolute
pivot into the diagonal position, then eliminate below.

\begin{lstlisting}[language=Rust]
pub fn ols(x: &[Vec<f64>], y: &[f64]) -> Result<OlsFit, TimeSeriesError> {
    let n = x.len();
    let k = x[0].len();
    // Normal equations: A = X'X, b = X'y.
    let mut a = vec![vec![0.0_f64; k]; k];
    let mut b = vec![0.0_f64; k];
    for (xi, yi) in x.iter().zip(y.iter()) {
        for i in 0..k {
            b[i] += xi[i] * yi;
            for j in 0..k {
                a[i][j] += xi[i] * xi[j];
            }
        }
    }
    let coeffs = gauss_solve(&mut a, &mut b)?;
    // residuals, SSE, SST, R^2, standard errors, t-stats ...
}
\end{lstlisting}

\begin{keyinsight}
Partial pivoting matters. Without it, a small diagonal entry relative to the
entries below causes catastrophic cancellation when we eliminate. Pivoting
costs nothing (a row swap per column) and lifts the worst-case error from
exponential in $k$ to roughly $k$ machine epsilons.
\end{keyinsight}

Standard errors require the diagonal of $(X'X)^{-1}$. Rather than inverting the
full matrix, we solve $X'X\, e_j = \text{unit}_j$ for each column $j$ and read
off $e_j[j] = (X'X)^{-1}_{jj}$. The standard error is
\[
\text{SE}(\hat\beta_j) = \sqrt{(X'X)^{-1}_{jj} \, s^2},
\quad s^2 = \frac{\text{SSE}}{n - k}.
\]
The t-statistic is $\hat\beta_j / \text{SE}(\hat\beta_j)$ and
$R^2 = 1 - \text{SSE}/\text{SST}$.

\section{Autocorrelation Function}

\marginnote{%
\textbf{Sample ACF}\\[0.3em]
$\hat\rho_k = \dfrac{\sum_{t=k+1}^{n}(y_t - \bar y)(y_{t-k} - \bar y)}
{\sum_{t=1}^{n}(y_t - \bar y)^2}$, with $\hat\rho_0 = 1$ by construction.
}

The autocorrelation at lag $k$ is the correlation between $y_t$ and
$y_{t-k}$. White noise has $\hat\rho_k \approx 0$ for $k > 0$; an AR(1) with
coefficient $\phi$ has $\hat\rho_k \approx \phi^k$; a random walk has
$\hat\rho_k \approx 1$ for all small $k$.

\begin{lstlisting}[language=Rust]
pub fn acf(data: &[f64], max_lag: usize)
    -> Result<Vec<f64>, TimeSeriesError>
{
    let n = data.len();
    let ybar: f64 = data.iter().sum::<f64>() / n as f64;
    let denom: f64 = data.iter()
        .map(|&v| { let d = v - ybar; d * d })
        .sum();
    let mut out = Vec::with_capacity(max_lag + 1);
    for k in 0..=max_lag {
        if k == 0 { out.push(1.0); continue; }
        let cov: f64 = (k..n)
            .map(|t| (data[t] - ybar) * (data[t - k] - ybar))
            .sum();
        out.push(cov / denom);
    }
    Ok(out)
}
\end{lstlisting}

\section{Augmented Dickey-Fuller Test}

\marginnote{%
\textbf{ADF regression}\\[0.3em]
$\Delta y_t = \alpha + \gamma\, y_{t-1} + \sum_{i=1}^{p} \delta_i\, \Delta y_{t-i} + \varepsilon_t$.\\[0.5em]
Reject $H_0: \gamma = 0$ (unit root) when the t-ratio on $\hat\gamma$ is below
the MacKinnon 5\% critical value $-2.86$.
}

The ADF test regresses the first difference on the lagged level and lagged
differences. The null hypothesis is a unit root ($\gamma = 0$); the alternative
is stationarity ($\gamma < 0$). Because the test statistic does not follow a
standard distribution under the null, we compare it to MacKinnon's simulated
critical values. For a constant-only specification at 5\% significance the
value is $-2.86$.

\begin{lstlisting}[language=Rust]
pub const MACKINNON_5PCT: f64 = -2.86;

pub fn adf_test(data: &[f64], lags: usize)
    -> Result<AdfResult, TimeSeriesError>
{
    let n = data.len();
    let dy: Vec<f64> = (1..n).map(|t| data[t] - data[t - 1]).collect();
    // Rows: [1, y_{t-1}, dy[t-1], dy[t-2], ..., dy[t-lags]]
    let mut x: Vec<Vec<f64>> = Vec::new();
    let mut y: Vec<f64> = Vec::new();
    for t in (lags + 1)..n {
        let mut row = Vec::with_capacity(lags + 2);
        row.push(1.0);
        row.push(data[t - 1]);
        for i in 1..=lags { row.push(dy[t - 1 - i]); }
        y.push(dy[t - 1]);
        x.push(row);
    }
    let fit = ols(&x, &y)?;
    let stat = fit.t_stats[1]; // gamma is index 1 (after intercept)
    Ok(AdfResult {
        statistic: stat,
        critical_value: MACKINNON_5PCT,
        lags,
        is_stationary: stat < MACKINNON_5PCT,
    })
}
\end{lstlisting}

\begin{keyinsight}
The lagged-level coefficient $\gamma$ is the key. If $\gamma = 0$, shocks to
$y_t$ are permanent (unit root). If $\gamma < 0$, shocks decay at rate
$1 + \gamma$ per period. The more negative $\gamma$, the faster the series
reverts to the mean --- and the more strongly stationary.
\end{keyinsight}

\section{Fractional Differentiation}

\marginnote{%
\textbf{Binomial series}\\[0.3em]
$(1 - B)^d = \sum_{k=0}^{\infty} w_k B^k$, where
$w_0 = 1$ and $w_k = -w_{k-1}\,\dfrac{d - k + 1}{k}$.\\[0.5em]
For $d = 1$: $w = [1, -1, 0, \ldots]$.\\
For $d \in (0, 1)$: infinitely many non-zero weights decaying in magnitude.
}

The backward shift $B$ satisfies $B y_t = y_{t-1}$. The first difference is
$(1 - B) y_t$. Fractional differentiation generalises this to
$(1 - B)^d$ for any real $d \in [0, 2]$. The binomial series gives the weights:

\begin{align*}
w_0 &= 1, \\
w_k &= -w_{k-1}\, \frac{d - k + 1}{k} \quad (k \geq 1).
\end{align*}

For $d = 1$ the weights are $[1, -1, 0, 0, \ldots]$ --- the ordinary first
difference. For $0 < d < 1$ infinitely many weights are non-zero, with
$|w_k|$ decaying roughly as $k^{-(1 + d)}$. We truncate when
$|w_k| < \text{threshold}$ (the fixed-width window of López de Prado, AFML
Ch.5):

\begin{lstlisting}[language=Rust]
pub fn ffd_weights(d: f64, threshold: f64)
    -> Result<Vec<f64>, TimeSeriesError>
{
    let mut weights = vec![1.0_f64];
    let mut w = 1.0_f64;
    let mut k = 1_usize;
    loop {
        w = -w * (d - (k as f64) + 1.0) / k as f64;
        if w.abs() < threshold { break; }
        weights.push(w);
        if k > 10_000 { break; }
        k += 1;
    }
    Ok(weights)
}
\end{lstlisting}

\begin{keyinsight}
The loop breaks \emph{before} pushing a weight below the threshold. This
makes $d = 0$ yield $[1.0]$ (identity, full memory) and $d = 1$ yield
$[1.0, -1.0]$ exactly (ordinary first difference, no memory). Every
intermediate $d$ keeps a finite prefix of the binomial series.
\end{keyinsight}

Applying the window is a weighted moving average over the last
$\texttt{weights.len()}$ observations:

\begin{lstlisting}[language=Rust]
pub fn frac_diff(data: &[f64], d: f64, threshold: f64)
    -> Result<Vec<f64>, TimeSeriesError>
{
    let weights = ffd_weights(d, threshold)?;
    let wlen = weights.len();
    let mut out = Vec::with_capacity(data.len() - wlen + 1);
    for t in (wlen - 1)..data.len() {
        let acc: f64 = weights.iter()
            .enumerate()
            .map(|(k, &w)| w * data[t - k])
            .sum();
        out.push(acc);
    }
    Ok(out)
}
\end{lstlisting}

\subsection{Finding the Minimum $d$}

Binary search over $d \in [0, 1]$ returns the smallest order for which the
ADF test rejects the unit root. We test the endpoints first: if $d = 0$
already passes, no differencing is needed; if $d = 1$ fails, we return $1.0$
conservatively.

\begin{lstlisting}[language=Rust]
pub fn find_min_d(data: &[f64], threshold: f64, tolerance: f64)
    -> Result<f64, TimeSeriesError>
{
    let passes = |d: f64| {
        let diff = frac_diff(data, d, threshold)?;
        Ok(adf_test(&diff, 1)?.is_stationary)
    };
    if passes(0.0)? { return Ok(0.0); }
    if !passes(1.0)? { return Ok(1.0); }
    let (mut lo, mut hi) = (0.0_f64, 1.0_f64);
    while hi - lo > tolerance {
        let mid = 0.5 * (lo + hi);
        if passes(mid)? { hi = mid; } else { lo = mid; }
    }
    Ok(hi)
}
\end{lstlisting}

\section{The Stationarity-Memory Tradeoff in Action}

\begin{lstlisting}[language=bash]
$ cargo run -p quant-timeseries --example stationarity
==============================
Stationarity Analysis
==============================

Random Walk (500 steps):
  ADF statistic: -1.1064
  Critical (5%):  -2.86
  Stationary:     NO

First Difference (d=1):
  ADF statistic: -15.8418
  Stationary:    YES
  ACF(1):        -0.0269 (memory destroyed)

Fractional Diff (d=0.4):
  ADF statistic: -4.8991
  Stationary:    YES
  ACF(1):         0.6878 (memory preserved)
==============================
\end{lstlisting}

The random walk has an ADF statistic of $-1.1$, well above the $-2.86$
critical value: we cannot reject the unit root. First differencing fixes
stationarity (statistic plunges to $-15.8$) but the ACF at lag 1 is
essentially zero --- all memory is gone. Fractional differentiation with
$d = 0.4$ also rejects the unit root, but the ACF at lag 1 is $0.69$:
most of the long-range information survives.

\begin{keyinsight}
This is the core finding of AFML Ch.5: there exists a $d^\star \in (0, 1)$
that simultaneously satisfies the stationarity requirement of statistical
inference \emph{and} the memory requirement of predictive ML features. Any
$d < d^\star$ is non-stationary; any $d > d^\star$ throws away more memory
than necessary.
\end{keyinsight}

\section{Exercises}

\begin{enumerate}
    \item \textbf{Partial autocorrelation}: Implement the PACF via the
      Yule-Walker equations. For an AR($p$) process, PACF should cut off after
      lag $p$. Verify on simulated AR(2) data with $\phi_1 = 0.5$,
      $\phi_2 = -0.3$.

    \item \textbf{KPSS test}: The KPSS test has the opposite null hypothesis
      to ADF --- $H_0$ is stationarity. Implement it and compare the
      conclusions on white noise, random walk, and fractionally differenced
      series. Where do the two tests disagree?

    \item \textbf{Real returns}: Generate a long GBM price path and find the
      minimum $d$ for stationarity. Repeat for different volatilities
      $\sigma \in \{0.1, 0.2, 0.4\}$. Does $\sigma$ affect $d^\star$?

    \item \textbf{Expanding-window frac\_diff}: The fixed-width window drops
      the first $K - 1$ observations. Implement an expanding-window variant
      that uses all available history for the early points (weights are
      truncated, not zero). What is the tradeoff versus fixed-width?

    \item \textbf{Lag selection for ADF}: The AIC and BIC criteria pick the
      number of lagged differences automatically. Implement AIC-based lag
      selection and check whether the stationarity conclusion changes for
      your test series.
\end{enumerate}
\textit{Chapter content pending --- Phase 7 implementation.}

\chapter{Volatility Models: EWMA, ARCH, GARCH}
\label{ch:volatility}
% \chapter{Volatility Models: EWMA, ARCH, and GARCH}
\label{ch:volatility}

\begin{companioncode}{quant-vol}
Code: \texttt{crates/quant-vol/}\\
Dependencies: \crate{quant-core}, \crate{thiserror}\\
Run: \texttt{cargo run -p quant-vol --example vol\_demo}
\end{companioncode}

\section{Learning Objectives}

Financial returns are not Gaussian --- their variance changes over time, and
large moves cluster. A constant-volatility model underestimates risk in
turbulent periods and overestimates it in calm ones. After completing this
chapter you can:

\begin{itemize}
    \item Describe the stylized facts of financial return volatility:
      clustering, fat tails, mean reversion, and slow decay of autocorrelation
      in squared returns
    \item Derive the exponentially weighted moving average (EWMA) recursion
      and justify the RiskMetrics $\lambda = 0.94$ choice for daily data
    \item Specify the ARCH($q$) model of Engle (1982) and fit it by Gaussian
      maximum likelihood with a hand-rolled optimiser
    \item Specify the GARCH($p, q$) model of Bollerslev (1986), state the
      covariance-stationarity condition, and interpret persistence
    \item Compute the long-run variance and the half-life of volatility shocks
    \item Fit GARCH(1,1) to simulated and clustered-volatility data and show
      that it beats constant-volatility models in log-likelihood
\end{itemize}

\begin{keyinsight}
Volatility clustering --- large moves follow large moves, calm follows calm
--- is the single most robust empirical regularity in financial returns. A
Gaussian return with constant variance cannot reproduce it. ARCH and GARCH
make the conditional variance a function of past squared returns, so a large
shock raises the variance of the next period. This one feature explains most
of the improvement of GARCH over Gaussian models for risk forecasting.
\end{keyinsight}

\section{Stylized Facts of Volatility}

\marginnote{%
\textbf{Volatility clustering}\\[0.3em]
$\Corr(r_t^2, r_{t-k}^2) > 0$ for many lags $k$, while $\Corr(r_t, r_{t-k})
\approx 0$. The ACF of returns is flat; the ACF of \emph{squared} returns
decays slowly.
}

Empirical daily returns exhibit four robust features that any realistic
volatility model must reproduce:

\begin{enumerate}
    \item \textbf{Volatility clustering.} The autocorrelation of $r_t^2$ is
      positive and decays slowly --- much more slowly than the autocorrelation
      of $r_t$, which is essentially zero after one or two lags.
    \item \textbf{Fat tails.} The unconditional distribution of returns has
      excess kurtosis: large moves occur more often than a Gaussian with the
      same variance would predict.
    \item \textbf{Mean reversion.} Volatility cannot grow without bound; it
      reverts toward a long-run level.
    \item \textbf{Asymmetry (leverage).} Negative returns tend to raise
      volatility more than positive returns of the same magnitude. Plain
      ARCH/GARCH do not capture this; GJR-GARCH and EGARCH do.
\end{enumerate}

\section{EWMA: Exponentially Weighted Moving Average}

\marginnote{%
\textbf{EWMA recursion}\\[0.3em]
$\sigma_t^2 = \lambda \sigma_{t-1}^2 + (1 - \lambda) r_{t-1}^2$\\[0.3em]
with $\sigma_0^2 = r_0^2$. RiskMetrics uses $\lambda = 0.94$ for daily data
and $\lambda = 0.97$ for monthly.
}

The simplest model that responds to clustering is the exponentially weighted
moving average:
\[
\sigma_t^2 = \lambda \sigma_{t-1}^2 + (1 - \lambda) r_{t-1}^2,
\qquad \lambda \in [0, 1].
\]
Each period's variance is a convex combination of the previous variance and
the previous squared return. The parameter $\lambda$ controls the memory: a
value close to $1$ makes the variance slow-moving; a value close to $0$
makes it track the last squared return.

\subsection{Boundary cases}

\begin{itemize}
    \item $\lambda = 0$: $\sigma_t^2 = r_{t-1}^2$ --- the variance is just the
      lagged squared return. No smoothing.
    \item $\lambda = 1$: $\sigma_t^2 = \sigma_{t-1}^2$ --- the variance is
      constant at its initial value $\sigma_0^2 = r_0^2$. No adaptation.
    \item $\lambda = 0.94$ (RiskMetrics daily): a single shock's influence
      decays geometrically. After $k$ periods without further shocks, the
      variance is multiplied by $\lambda^k$.
\end{itemize}

\subsection{Unrolling the recursion}

Substituting repeatedly:
\[
\sigma_t^2 = (1 - \lambda) \sum_{k=1}^{t} \lambda^{k-1} r_{t-k}^2
           + \lambda^t \sigma_0^2.
\]
The weights on past squared returns decay geometrically with rate $\lambda$.
For long series the $\lambda^t \sigma_0^2$ term vanishes and the weights sum
to $(1 - \lambda) / (1 - \lambda) = 1$, so $\sigma_t^2$ is a proper weighted
average of past squared returns.

\begin{lstlisting}[language=Rust]
pub fn ewma_vol(returns: &[f64], lambda: f64) -> Result<Vec<f64>, VolError> {
    if !(0.0..=1.0).contains(&lambda) {
        return Err(VolError::InvalidParam(...));
    }
    if returns.is_empty() {
        return Err(VolError::InsufficientData { required: 1, actual: 0 });
    }
    let mut sigma2 = vec![0.0_f64; returns.len()];
    sigma2[0] = returns[0] * returns[0];
    for t in 1..returns.len() {
        sigma2[t] = lambda * sigma2[t - 1]
                  + (1.0 - lambda) * returns[t - 1] * returns[t - 1];
    }
    Ok(sigma2)
}
\end{lstlisting}

\section{ARCH: Autoregressive Conditional Heteroskedasticity}

\marginnote{%
\textbf{ARCH($q$) of Engle (1982)}\\[0.3em]
$\sigma_t^2 = \omega + \sum_{i=1}^q \alpha_i r_{t-i}^2$\\[0.3em]
with $\omega > 0$, $\alpha_i \geq 0$, and $\sum \alpha_i < 1$ for
stationarity.
}

Engle's ARCH model makes the conditional variance an explicit linear function
of the past $q$ squared returns:
\[
\sigma_t^2 = \omega + \sum_{i=1}^q \alpha_i r_{t-i}^2.
\]
The coefficients are constrained to be non-negative (otherwise the variance
could go negative) and $\sum \alpha_i < 1$ ensures the unconditional variance
$\omega / (1 - \sum \alpha_i)$ is finite.

ARCH captures clustering directly: a large $r_{t-1}^2$ raises $\sigma_t^2$ by
$\alpha_1 r_{t-1}^2$, which in turn makes a large $r_t$ more likely, which
raises $\sigma_{t+1}^2$, and so on. The effect dies out geometrically when
shocks stop arriving.

\subsection{Gaussian log-likelihood}

Assuming $r_t \mid \mathcal{F}_{t-1} \sim \Normal(0, \sigma_t^2)$, the
conditional density is
\[
f(r_t \mid \mathcal{F}_{t-1}) = \frac{1}{\sqrt{2\pi\sigma_t^2}}
  \exp\!\left(-\frac{r_t^2}{2\sigma_t^2}\right),
\]
so the Gaussian log-likelihood is
\[
\mathcal{L}(\theta) = -\frac{1}{2} \sum_t
  \left[ \log(2\pi) + \log\sigma_t^2(\theta) + \frac{r_t^2}{\sigma_t^2(\theta)} \right].
\]
The MLE maximises $\mathcal{L}$ over $\theta = (\omega, \alpha_1, \ldots,
\alpha_q)$. Because $\sigma_t^2$ is a deterministic function of $\theta$ and
the data, the score and Hessian can be derived in closed form, but we will
not need them: a derivative-free optimiser is enough.

\begin{warning}
The Gaussian log-likelihood is \emph{not} bounded above by zero. When
$\sigma_t^2$ is small and $r_t^2 / \sigma_t^2$ is also small, each term
$-0.5[\log(2\pi) + \log\sigma_t^2 + r_t^2/\sigma_t^2]$ can be positive. The
only universal constraint is that $\mathcal{L}$ is finite when all
$\sigma_t^2 > 0$.
\end{warning}

\begin{lstlisting}[language=Rust]
#[derive(Debug, Clone)]
pub struct ArchModel {
    pub omega: f64,       // Constant term (omega > 0)
    pub alphas: Vec<f64>, // ARCH coefficients (alpha_i >= 0)
}

impl ArchModel {
    pub fn conditional_variances(&self, returns: &[f64]) -> Vec<f64> {
        let n = returns.len();
        let q = self.alphas.len();
        let mut sigma2 = vec![0.0_f64; n];
        for t in 0..n {
            let mut s = self.omega;
            for i in 0..q {
                if t > i {
                    s += self.alphas[i] * returns[t - 1 - i].powi(2);
                }
            }
            sigma2[t] = s.max(1e-300);
        }
        sigma2
    }

    pub fn log_likelihood(&self, returns: &[f64]) -> f64 {
        let sigma2 = self.conditional_variances(returns);
        let mut ll = 0.0_f64;
        for (t, &r) in returns.iter().enumerate() {
            let s2 = sigma2[t];
            ll -= 0.5 * (LN_2PI + s2.ln() + r * r / s2);
        }
        ll
    }
}
\end{lstlisting}

\section{GARCH: Generalised ARCH}

\marginnote{%
\textbf{GARCH($p, q$) of Bollerslev (1986)}\\[0.3em]
$\sigma_t^2 = \omega + \sum_{i=1}^q \alpha_i r_{t-i}^2 + \sum_{j=1}^p
\beta_j \sigma_{t-j}^2$\\[0.3em]
with $\omega > 0$, $\alpha_i, \beta_j \geq 0$, and $\sum\alpha_i + \sum\beta_j
< 1$ for covariance stationarity.
}

Bollerslev's GARCH adds lagged conditional variances to the ARCH recursion:
\[
\sigma_t^2 = \omega + \sum_{i=1}^q \alpha_i r_{t-i}^2
                   + \sum_{j=1}^p \beta_j \sigma_{t-j}^2.
\]
The GARCH terms $\beta_j \sigma_{t-j}^2$ propagate current volatility forward,
so shocks have a longer, smoother effect than in ARCH. In practice
GARCH(1,1) --- one ARCH term and one GARCH term --- is the workhorse: it fits
most daily return series well with only three parameters.

\begin{lstlisting}[language=Rust]
#[derive(Debug, Clone)]
pub struct GarchModel {
    pub omega: f64,       // Constant term (omega > 0)
    pub alphas: Vec<f64>, // ARCH coefficients, length q
    pub betas: Vec<f64>,  // GARCH coefficients, length p
}

impl GarchModel {
    pub fn conditional_variances(&self, returns: &[f64]) -> Vec<f64> {
        let n = returns.len();
        let (q, p) = (self.alphas.len(), self.betas.len());
        let warmup = p.max(q);
        let init_var = sample_variance(returns);
        let mut sigma2 = vec![init_var; n];
        for t in warmup..n {
            let mut s = self.omega;
            for i in 0..q {
                s += self.alphas[i] * returns[t - 1 - i].powi(2);
            }
            for j in 0..p {
                s += self.betas[j] * sigma2[t - 1 - j];
            }
            sigma2[t] = s.clamp(1e-300, 1e15);
        }
        sigma2
    }

    pub fn persistence(&self) -> f64 {
        self.alphas.iter().sum::<f64>() + self.betas.iter().sum::<f64>()
    }

    pub fn long_run_variance(&self) -> f64 {
        let pers = self.persistence();
        if pers < 1.0 { self.omega / (1.0 - pers) } else { f64::INFINITY }
    }

    pub fn half_life(&self) -> f64 {
        let pers = self.persistence();
        if pers > 0.0 && pers < 1.0 { 0.5_f64.ln() / pers.ln() }
        else { f64::INFINITY }
    }
}
\end{lstlisting}

\subsection{Persistence and stationarity}

\marginnote{%
\textbf{Persistence}\\[0.3em]
$\rho = \sum_i \alpha_i + \sum_j \beta_j$.\\[0.3em]
Covariance stationary iff $\rho < 1$.\\[0.3em]
Long-run variance: $\bar\sigma^2 = \omega / (1 - \rho)$.\\[0.3em]
Half-life of shocks: $\log(0.5) / \log\rho$.
}

The sum $\rho = \sum \alpha_i + \sum \beta_j$ is the \emph{persistence}. Taking
unconditional expectations of the GARCH recursion and solving for the fixed
point gives the long-run variance:
\[
\bar\sigma^2 = \frac{\omega}{1 - \rho}, \qquad \rho < 1.
\]
When $\rho \geq 1$ the model is non-stationary: shocks accumulate rather than
decay, and the unconditional variance is infinite. The IGARCH boundary
$\rho = 1$ is a random walk in variance and should be avoided for fitting.

\subsection{Half-life of shocks}

Starting from $\sigma_0^2 = \bar\sigma^2 + \delta$ (a perturbation $\delta$
above the long run), the expected deviation from $\bar\sigma^2$ after $k$
steps is $\rho^k \delta$. The half-life is the $k$ at which $\rho^k = 0.5$:
\[
k_{1/2} = \frac{\log 0.5}{\log \rho}.
\]
For a typical daily GARCH(1,1) with $\rho = 0.99$, the half-life is about
$69$ trading days --- roughly three months. High persistence is the norm for
daily financial data.

\subsection{Forecasting}

The multi-step forecast follows the AR(1)-in-variance recursion
\[
\E[\sigma_{t+h}^2 \mid \mathcal{F}_t] = \omega + \rho \,
  \E[\sigma_{t+h-1}^2 \mid \mathcal{F}_t],
\]
which reverts geometrically to $\bar\sigma^2$:
\[
\E[\sigma_{t+h}^2 \mid \mathcal{F}_t] = \bar\sigma^2 + \rho^h
  \bigl(\sigma_t^2 - \bar\sigma^2\bigr).
\]
The forecast is data-conditioned when the last fitted $\sigma_t^2$ is used as
the seed, or unconditional when seeded from $\bar\sigma^2$.

\section{Maximum Likelihood by Hand-Rolled Optimisation}

\marginnote{%
\textbf{Constrained MLE}\\[0.3em]
Optimise an unconstrained $\theta$ and map to the constrained space:
$\omega = e^{t_0}$, $\alpha_i = \mathrm{sigmoid}(t_i) \cdot 0.49/q$,
$\beta_j = \mathrm{sigmoid}(t_j) \cdot (0.998 - \sum\alpha)/p$. This
guarantees $\omega > 0$, $\alpha_i, \beta_j \geq 0$, and $\rho < 0.998$.
}

We fit both ARCH and GARCH by maximising the Gaussian log-likelihood with a
derivative-free optimiser. The constraint $\rho < 1$ is enforced not by
penalising the objective but by \emph{parameter transformation}: we optimise
an unconstrained vector $t$ and map it to the constrained space via
$\omega = e^{t_0}$, $\alpha_i = \mathrm{sigmoid}(t_i) \cdot 0.49 / q$, and
$\beta_j = \mathrm{sigmoid}(t_j) \cdot (0.998 - \sum\alpha) / p$. The
$0.998$ cap on persistence prevents the optimiser from wandering onto the
IGARCH boundary, where the likelihood becomes flat and numerically
ill-conditioned.

\subsection{Nelder-Mead simplex}

The Nelder-Mead algorithm maintains a simplex of $n+1$ points in
$n$-dimensional parameter space and cycles through four operations:
\begin{description}
    \item[Reflection] ($\alpha = 1$) Mirror the worst vertex through the
      centroid of the others.
    \item[Expansion] ($\gamma = 2$) If the reflection is better than the
      best, stretch further in the same direction.
    \item[Contraction] ($\rho = 0.5$) If the reflection is worse than the
      second-worst, contract toward the centroid.
    \item[Shrink] ($\sigma = 0.5$) If contraction fails, pull all vertices
      halfway toward the best.
\end{description}
Convergence is declared when both the simplex diameter and the spread of
function values fall below a tolerance.

\subsection{Coordinate ascent with multi-start}

\marginnote{%
\textbf{Why coordinate ascent?}\\[0.3em]
For 2--4 parameter MLE problems, Nelder-Mead can degenerate near the
constraint boundary. Coordinate ascent with step halving tracks the
$\alpha$--$\beta$ ridge cleanly, and multi-start from five perturbed
initial guesses escapes the occasional local plateau.
}

For the low-dimensional MLE problems here (ARCH(1) has 2 parameters,
GARCH(1,1) has 3), a simpler and more robust alternative is coordinate
ascent: cycle through each coordinate, try a positive and a negative step,
accept any improvement, and halve the step when no coordinate improves.
Convergence is declared when the step falls below the tolerance. We wrap
this in a \emph{multi-start} that runs coordinate ascent from the initial
guess plus five perturbations $[+0.1, -0.1, +0.3, -0.3, +0.5]$ and keeps the
best result.

\begin{lstlisting}[language=Rust]
/// Coordinate-wise hill climbing with adaptive step size.
pub fn coordinate_ascent<F>(f: F, x0: &[f64], step: f64, max_iter: usize, tol: f64)
    -> OptResult
where F: Fn(&[f64]) -> f64
{
    let mut x = x0.to_vec();
    let mut best_f = f(&x);
    let mut step = step;
    for _ in 0..max_iter {
        let mut improved = false;
        for i in 0..x.len() {
            let old = x[i];
            x[i] += step;  // try positive step
            if f(&x) > best_f { best_f = f(&x); improved = true; continue; }
            x[i] = old - step;  // try negative step
            if f(&x) > best_f { best_f = f(&x); improved = true; continue; }
            x[i] = old;  // restore
        }
        if !improved { step *= 0.5; if step < tol { break; } }
    }
    OptResult { x, f: best_f, .. }
}

/// Multi-start: run coordinate_ascent from x0 plus perturbations.
pub fn multistart<F>(f: F, x0: &[f64], step: f64, max_iter: usize, tol: f64)
    -> OptResult
where F: Fn(&[f64]) -> f64
{
    let mut best = coordinate_ascent(&f, x0, step, max_iter, tol);
    for &perturb in &[0.1, -0.1, 0.3, -0.3, 0.5] {
        let x_alt: Vec<f64> = x0.iter().map(|&v| v + perturb).collect();
        let result = coordinate_ascent(&f, &x_alt, step, max_iter, tol);
        if result.f > best.f { best = result; }
    }
    best
}
\end{lstlisting}

All three optimisers \emph{maximise} the objective --- the MLE closures pass
\texttt{model.log\_likelihood(returns)} directly, without negation.

\section{Implementation in Rust}

\marginnote{%
\textbf{Why three models in one crate?}\\[0.3em]
EWMA, ARCH, and GARCH form a progression: EWMA is one parameter
($\lambda$); ARCH(1) adds a lag coefficient ($\omega, \alpha$);
GARCH(1,1) adds the lagged-variance coefficient ($\beta$). Keeping
them in one crate lets you compare log-likelihoods and in-sample fit
on the same return series with the same MLE machinery.
}

\begin{lstlisting}[language=Rust]
use quant_vol::{ewma_vol, ArchModel, GarchModel};

let returns = vec![0.01, -0.02, 0.015, -0.01, 0.02, -0.005, 0.01, -0.03];

// EWMA (RiskMetrics lambda = 0.94).
let sigma2 = ewma_vol(&returns, 0.94).unwrap();

// ARCH(1) fit by Gaussian MLE.
let arch = ArchModel::fit(&returns, 1).unwrap();
let arch_ll = arch.log_likelihood(&returns);

// GARCH(1,1) fit by Gaussian MLE.
let garch = GarchModel::fit(&returns, 1, 1).unwrap();
println!("persistence = {:.4}", garch.persistence());
println!("long-run var = {:.6}", garch.long_run_variance());
println!("half-life = {:.1} periods", garch.half_life());
\end{lstlisting}

The constrained parameterisation for GARCH(1,1) is:
\begin{lstlisting}[language=Rust]
fn from_transformed(params: &[f64], p: usize, q: usize) -> Self {
    let omega = params[0].clamp(-30.0, 30.0).exp();
    let alpha_max = 0.49;
    let alphas: Vec<f64> = (0..q)
        .map(|i| sigmoid(params[1 + i].clamp(-30.0, 30.0)) * alpha_max / q as f64)
        .collect();
    let alpha_sum: f64 = alphas.iter().sum();
    let beta_max = (0.998 - alpha_sum).max(1e-6);
    let betas: Vec<f64> = (0..p)
        .map(|j| sigmoid(params[1 + q + j].clamp(-30.0, 30.0)) * beta_max / p as f64)
        .collect();
    GarchModel { omega, alphas, betas }
}
\end{lstlisting}
The initial guess targets $\omega \approx 0.05 \hat\sigma^2$,
$\alpha \approx 0.05$, $\beta \approx 0.90$, which are typical for daily
financial returns.

\section{Volatility Clustering in Action}

\begin{lstlisting}[language=bash]
$ cargo run -p quant-vol --example vol_clustering
Regimes: calm(1000) -> shock(500) -> calm(500)
  Calm sigma^2  = 0.000100
  Shock sigma^2 = 0.010000

Constant-volatility model:
  Sample variance = 0.002585
  Log-likelihood  = 3120.38

GARCH(1,1):
  omega = 0.000004
  alpha = 0.145513
  beta  = 0.849413
  Persistence = 0.9949
  Log-likelihood = 5072.78

  -> GARCH beats constant-vol by 1952.4 LL points
\end{lstlisting}

The constant-volatility model uses a single sample variance for every period,
so it over-estimates variance during the calm regime (producing too-small
densities for the small returns that actually occur) and under-estimates it
during the shock regime (producing too-small densities for the large returns
that actually occur). GARCH tracks the regime changes through its conditional
variance recursion and wins by nearly $2000$ log-likelihood points.

\begin{keyinsight}
The log-likelihood gain from GARCH over constant volatility is a direct
quantification of how much volatility clustering matters. A gain of $1952$
LL points across $2000$ observations is roughly $0.976$ nats per observation
--- a substantial amount of explained structure that a Gaussian-with-constant-
variance model leaves on the table.
\end{keyinsight}

\section{Parameter Recovery}

When the data are generated by a known GARCH(1,1), the MLE should recover the
true parameters up to sampling error. With $2000$ observations from
$\omega = 0.01$, $\alpha = 0.08$, $\beta = 0.90$ (true persistence $0.98$,
true long-run variance $0.5$):

\begin{lstlisting}[language=bash]
$ cargo run -p quant-vol --example vol_demo
DGP: GARCH(1,1) with omega=0.01, alpha=0.08, beta=0.90
True persistence:     0.9800
True long-run var:     0.500000

GARCH(1,1):
  omega = 0.0111
  alpha = 0.0910
  beta  = 0.8934
  Persistence = 0.9844  (true 0.9800)
  Long-run var = 0.7150  (true 0.500000)
  Half-life    = 44.1 periods
  Log-likelihood = -2298.62
\end{lstlisting}

The fitted persistence is within $0.005$ of the truth. The long-run variance
is harder to pin down because it divides $\omega$ by $1 - \rho$, so a small
error in $\rho$ near $1$ amplifies. This is a general feature of
high-persistence GARCH fits: $\omega$ and $\rho$ are well-identified, but
$\bar\sigma^2 = \omega / (1 - \rho)$ has a much wider confidence band.

\section{Exercises}

\begin{enumerate}
    \item \textbf{EGARCH}: Exponential GARCH models $\log \sigma_t^2$ as a
      linear function of past squared returns and their signs, which captures
      the leverage effect (negative returns raise volatility more than
      positive ones). Implement EGARCH(1,1) and compare its log-likelihood
      to plain GARCH(1,1) on a series with a visible leverage effect.

    \item \textbf{GJR-GARCH}: The Glosten-Jagannathan-Runkle model adds an
      indicator term that activates when the lagged return is negative:
      $\sigma_t^2 = \omega + \alpha r_{t-1}^2 + \gamma \mathbf{1}_{r_{t-1}<0}
      r_{t-1}^2 + \beta \sigma_{t-1}^2$. Implement it and test on equity
      index returns. Is $\gamma$ significantly positive?

    \item \textbf{Forecast evaluation}: Fit GARCH(1,1) on the first $80\%$ of
      a return series, forecast the conditional variance for the remaining
      $20\%$, and compare the forecast bands to the realised squared returns.
      Compute the mean squared error of the variance forecast and compare to
      the EWMA forecast.

    \item \textbf{Student-$t$ innovations}: Replace the Gaussian likelihood
      with a Student-$t$ likelihood with $\nu$ degrees of freedom (fit $\nu$
      jointly with $(\omega, \alpha, \beta)$). Does the $t$-GARCH fit
      outperform the Gaussian GARCH on a heavy-tailed series? What is the
      fitted $\nu$?

    \item \textbf{Long-memory FIGARCH}: The FIGARCH model replaces the
      geometric decay of shocks with a hyperbolic decay by fractionally
      differencing the variance recursion. Implement FIGARCH($1, d, 1$) and
      compare its half-life to plain GARCH(1,1) on a long daily series. How
      much longer is the effective memory?
\end{enumerate}

\section{Key Takeaways}

\begin{itemize}
    \item Financial returns show volatility clustering: squared returns are
      positively autocorrelated for many lags. Constant-volatility models
      cannot reproduce this.
    \item EWMA is the simplest fix: a one-parameter geometrically weighted
      average of past squared returns. RiskMetrics uses $\lambda = 0.94$ for
      daily data.
    \item ARCH($q$) makes the conditional variance a linear function of the
      past $q$ squared returns. It captures clustering but needs many lags to
      match the slow decay of the squared-return ACF.
    \item GARCH($p, q$) adds lagged conditional variances, giving shocks a
      long, smooth effect with only a few parameters. GARCH(1,1) is the
      standard workhorse.
    \item Persistence $\rho = \sum\alpha + \sum\beta$ must be $< 1$ for
      covariance stationarity. The long-run variance is
      $\omega / (1 - \rho)$ and the half-life of shocks is
      $\log 0.5 / \log \rho$.
    \item Fitting by Gaussian MLE with a hand-rolled optimiser is feasible
      because the likelihood is a deterministic function of the parameters
      and the data. Parameter transformation (exp, sigmoid) enforces
      positivity and stationarity without penalising the objective.
\end{itemize}

\vspace{1em}
\noindent\textbf{Next chapter}: Stochastic processes---Brownian motion, Poisson
processes, and Monte Carlo simulation for derivative pricing.
\textit{Chapter content pending --- Phase 8 implementation.}

\chapter{Stochastic Processes and Monte Carlo}
\label{ch:stochastic}
% \chapter{Stochastic Processes and Monte Carlo}
\label{ch:stochastic}

\begin{companioncode}{quant-stochastic}
Code: \texttt{crates/quant-stochastic/}\\
Dependencies: \crate{quant-core}, \crate{thiserror}\\
Run: \texttt{cargo run -p quant-stochastic --example mc\_pricing}
\end{companioncode}

\section{Learning Objectives}

Asset prices are often modelled as stochastic processes --- random functions
of time. The Brownian motion is the continuous-time limit of a random walk and
the building block of modern quantitative finance. From it we derive geometric
Brownian motion (the Black-Scholes asset model), the Poisson process (jump
models), and the Monte Carlo method (the generic pricing engine). After
completing this chapter you can:

\begin{itemize}
    \item Simulate standard Brownian motion from independent Gaussian
      increments and verify its quadratic variation property
    \item Simulate geometric Brownian motion via the exact closed-form
      solution to the SDE $dS_t = \mu S_t \, dt + \sigma S_t \, dW_t$
    \item Simulate a Poisson process via exponential interarrival times and
      build a Merton jump-diffusion on top of GBM
    \item Price a European option by Monte Carlo and report a standard error
    \item State the $1/\sqrt{N}$ convergence law and demonstrate it
      numerically
    \item Reduce estimator variance with antithetic variates and verify
      put-call parity on simulated prices
\end{itemize}

\begin{keyinsight}
Monte Carlo pricing converges to the analytical Black-Scholes price as
$N \to \infty$, with standard error decreasing as $1/\sqrt{N}$. This is the
law of large numbers in action: quadruple the paths, halve the error. No
other numerical method gives such a clean, dimension-independent error
bound --- which is why Monte Carlo is the workhorse of derivatives pricing
in high dimensions.
\end{keyinsight}

\section{Brownian Motion}

\marginnote{%
\textbf{Brownian motion}\\[0.3em]
$W_0 = 0$, $W_{t+\Delta t} - W_t \sim \Normal(0, \Delta t)$, independent
increments. Quadratic variation $[W]_t = t$.
}

A standard Brownian motion $\{W_t\}$ is the continuous-time limit of a random
walk with independent Gaussian increments:
\[
W_0 = 0, \qquad W_{t+\Delta t} = W_t + \sqrt{\Delta t} \, Z, \quad
Z \sim \Normal(0, 1).
\]
Two properties define it:
\begin{enumerate}
    \item \textbf{Independent increments.} $W_{t+s} - W_t$ is independent of
      $\{W_u : u \leq t\}$.
    \item \textbf{Quadratic variation.} The sum of squared increments
      $\sum (W_{t_k} - W_{t_{k-1}})^2$ converges in probability to $T$ as
      the partition refines --- the defining property of Brownian motion that
      distinguishes it from smooth functions.
\end{enumerate}

\subsection{Quadratic Variation}

For a Brownian motion sampled at spacing $\Delta t$, the realised quadratic
variation
\[
[W]_T^{(\Delta t)} = \sum_{k=1}^{n} (W_{k \Delta t} - W_{(k-1)\Delta t})^2
\]
converges in probability to $T = n \Delta t$ as $\Delta t \to 0$. This is the
empirical signature of Brownian motion: the path is continuous but
nowhere differentiable, and its roughness is exactly $t$.

\section{Geometric Brownian Motion}

\marginnote{%
\textbf{GBM SDE}\\[0.3em]
$dS_t = \mu S_t \, dt + \sigma S_t \, dW_t$\\[0.3em]
Closed form: $S_T = S_0 \exp\!\big((\mu - \tfrac{1}{2}\sigma^2) T + \sigma W_T\big)$
}

The Black-Scholes asset model is geometric Brownian motion, the exponential
of Brownian motion with drift:
\[
dS_t = \mu S_t \, dt + \sigma S_t \, dW_t.
\]
It\^o's lemma gives the exact solution
\[
S_T = S_0 \exp\!\left( \left(\mu - \tfrac{1}{2}\sigma^2\right) T + \sigma W_T \right).
\]
Two facts follow immediately:
\begin{itemize}
    \item $\log(S_T / S_0) \sim \Normal((\mu - \tfrac{1}{2}\sigma^2)T,\; \sigma^2 T)$.
    \item $\E[S_T] = S_0 \, e^{\mu T}$ (the log-normal mean).
\end{itemize}
The companion crate uses the closed form for each step, so the terminal
distribution is exact for any $\Delta t$ --- no Euler discretisation error.

\section{Poisson Process and Jump-Diffusion}

\marginnote{%
\textbf{Poisson process}\\[0.3em]
$N(t) \sim \text{Poisson}(\lambda t)$. Interarrival times $\sim
\text{Exp}(\lambda)$, mean $1/\lambda$.
}

The Poisson process counts rare events: $N(t)$ is the number of events in
$[0, t]$, with $\E[N(t)] = \lambda t$. Interarrival times are exponential
with rate $\lambda$. The inverse-CDF sampler draws
$X = -\ln(1 - U) / \lambda$ with $U \sim \text{Uniform}(0, 1)$.

\subsection{Merton Jump-Diffusion}

\marginnote{%
\textbf{Merton model}\\[0.3em]
$dS_t = \mu S_t \, dt + \sigma S_t \, dW_t + (J - 1) S_t \, dN_t$\\[0.3em]
Between jumps: GBM. At a jump: $S \to S \cdot J$.
}

Robert Merton's jump-diffusion superimposes Poisson jumps on GBM:
\[
dS_t = \mu S_t \, dt + \sigma S_t \, dW_t + (J - 1) S_t \, dN_t,
\]
where $N_t$ is a Poisson process with rate $\lambda$ and $J$ is the jump
multiplier. Between jumps the process follows GBM; at each event time the
price is multiplied by $J$. The companion crate uses $J = e^{\mu_J}$ (a
deterministic log-normal jump, with the randomness in the timing). With
$\lambda = 0$ there are no jumps and the process reduces exactly to GBM.

\section{The Monte Carlo Method}

\marginnote{%
\textbf{MC estimator}\\[0.3em]
$\hat{C} = e^{-rT} \frac{1}{N} \sum_{i=1}^N \max(S_T^{(i)} - K, 0)$\\[0.3em]
$\text{SE}(\hat{C}) = e^{-rT} \frac{\text{std}(\text{payoff})}{\sqrt{N}}$
}

Monte Carlo prices a derivative by simulating many paths and averaging the
discounted payoff. For a European call under risk-neutral GBM:
\[
\hat{C} = e^{-rT} \frac{1}{N} \sum_{i=1}^N \max(S_T^{(i)} - K, 0),
\qquad
S_T^{(i)} = S_0 \exp\!\left( \left(r - \tfrac{1}{2}\sigma^2\right) T + \sigma \sqrt{T} \, Z_i \right).
\]
Each path is a single normal draw --- the terminal $Z_i$ --- so the estimator
is minimal-variance. The standard error of the mean is
\[
\text{SE}(\hat{C}) = e^{-rT} \frac{s}{\sqrt{N}}, \qquad
s^2 = \frac{1}{N-1} \sum_i (\text{payoff}_i - \bar{\text{payoff}})^2.
\]

\subsection{The $1/\sqrt{N}$ Law}

The standard error shrinks as $1/\sqrt{N}$: quadruple the paths, halve the
error. This is slow --- to reduce the error by $10\times$ one needs $100\times$
the paths --- but it is \emph{dimension-independent}. Deterministic quadrature
suffers the curse of dimensionality; Monte Carlo does not. For a
50-dimensional exotic option, Monte Carlo's error is still $O(1/\sqrt{N})$.

\subsection{Antithetic Variates}

\marginnote{%
\textbf{Antithetic variates}\\[0.3em]
Pair $Z$ and $-Z$. For monotone payoffs the pair average has lower
variance than two independent samples.
}

Variance reduction lets us extract more precision per random draw.
\emph{Antithetic variates} pair each $Z$ with $-Z$:
\[
\hat{C}_{\text{anti}} = e^{-rT} \frac{1}{2N} \sum_{i=1}^N
\big[\max(S_T(Z_i) - K, 0) + \max(S_T(-Z_i) - K, 0)\big].
\]
The two payoffs are perfectly negatively correlated (the call payoff is
monotone increasing in $Z$), so the variance of the pair average is lower
than the variance of two independent samples. The cost is one extra payoff
evaluation per draw --- negligible versus the cost of the normal draw.

\section{Black-Scholes: The Analytical Benchmark}

\marginnote{%
\textbf{Black-Scholes call}\\[0.3em]
$C = S_0 \Phi(d_1) - K e^{-rT} \Phi(d_2)$\\[0.3em]
$d_1 = \frac{\ln(S_0/K) + (r + \frac{1}{2}\sigma^2)T}{\sigma\sqrt{T}},\;
d_2 = d_1 - \sigma\sqrt{T}$
}

The Black-Scholes formula gives the closed-form European call price:
\[
C = S_0 \Phi(d_1) - K e^{-rT} \Phi(d_2).
\]
The normal CDF $\Phi$ is computed via the Abramowitz-Stegun (1964) formula
7.1.26 rational approximation of $\text{erf}$, with $\Phi(x) = \tfrac{1}{2}(1 +
\text{erf}(x/\sqrt{2}))$. Maximum error $< 7.5 \times 10^{-8}$ --- far smaller
than Monte Carlo's statistical error for practical $N$.

\section{Implementation in Rust}

\subsection{Brownian Motion}

\begin{lstlisting}[language=Rust]
pub fn brownian_motion<R: Rng + ?Sized>(n: usize, dt: f64, rng: &mut R) -> Vec<f64> {
    let normal = Normal::standard();
    let sqrt_dt = dt.sqrt();
    let mut path = Vec::with_capacity(n + 1);
    path.push(0.0); // W_0 = 0
    let mut w = 0.0_f64;
    for _ in 0..n {
        let z = normal.sample(rng);
        w += sqrt_dt * z;
        path.push(w);
    }
    path
}
\end{lstlisting}

\subsection{Monte Carlo Call}

\begin{lstlisting}[language=Rust]
pub fn mc_call<R: Rng + ?Sized>(
    s0: f64, k: f64, r: f64, sigma: f64, t: f64,
    n_paths: usize, rng: &mut R,
) -> Result<McResult, StochError> {
    validate_mc_inputs(s0, k, r, sigma, t, n_paths)?;
    let normal = Normal::standard();
    let drift = (r - 0.5 * sigma * sigma) * t;
    let diffusion = sigma * t.sqrt();
    let discount = (-r * t).exp();
    let mut payoffs = Vec::with_capacity(n_paths);
    for _ in 0..n_paths {
        let z = normal.sample(rng);
        let s_t = s0 * (drift + diffusion * z).exp();
        payoffs.push((s_t - k).max(0.0));
    }
    Ok(reduce(&payoffs, discount))
}
\end{lstlisting}

\section{Results and Analysis}

\subsection{Monte Carlo Convergence}

Pricing an at-the-money call ($S_0 = K = 100$, $r = 0.05$, $\sigma = 0.2$,
$T = 1$) against the Black-Scholes benchmark $C_{\text{BS}} = 10.4506$:

\begin{center}
\begin{tabular}{rrrr}
\toprule
$N$ & MC price & BS price & SE \\
\midrule
100 & 8.9770 & 10.4506 & 1.4320 \\
1000 & 10.3204 & 10.4506 & 0.4641 \\
10000 & 10.5059 & 10.4506 & 0.1469 \\
100000 & 10.4194 & 10.4506 & 0.0465 \\
\bottomrule
\end{tabular}
\end{center}

The estimate converges to Black-Scholes and the standard error shrinks as
$1/\sqrt{N}$.

\subsection{The $1/\sqrt{N}$ Law}

\begin{center}
\begin{tabular}{lrl}
\toprule
$N$ & SE & ratio vs $N=10000$ \\
\midrule
10000 & 0.1479 & 1.00 \\
40000 & 0.0732 & 2.02$\times$ smaller \\
160000 & 0.0367 & 4.03$\times$ smaller \\
\bottomrule
\end{tabular}
\end{center}

Quadrupling paths halves the SE ($2.02\times$); 16$\times$ paths gives
$4.03\times$ smaller SE --- the $1/\sqrt{N}$ law made exact.

\subsection{Antithetic Variance Reduction}

With 50000 normal draws each:

\begin{center}
\begin{tabular}{lrrr}
\toprule
Method & price & SE & payoffs \\
\midrule
Plain MC & 10.4191 & 0.0656 & 50000 \\
Antithetic & 10.4962 & 0.0467 & 100000 \\
\bottomrule
\end{tabular}
\end{center}

Antithetic variates reduce the standard error by \textbf{28.8\%} per normal
draw. The pair average exploits the negative correlation between the $Z$ and
$-Z$ payoffs.

\subsection{Quadratic Variation}

\begin{center}
\begin{tabular}{rrrr}
\toprule
$n$ & QV & $T$ & rel.\ error \\
\midrule
100 & 0.346 & 0.397 & 12.8\% \\
1000 & 4.095 & 3.968 & 3.2\% \\
10000 & 40.509 & 39.683 & 2.1\% \\
100000 & 399.560 & 396.825 & 0.69\% \\
\bottomrule
\end{tabular}
\end{center}

The realised quadratic variation converges to $T$ as the partition refines
--- the defining property of Brownian motion.

\subsection{GBM Terminal Distribution}

With 100000 paths ($S_0 = 100$, $\mu = 0.05$, $\sigma = 0.2$, $T = 1$):
\[
\E[S_T] = 105.07 \quad \text{vs analytical } S_0 e^{\mu T} = 105.13.
\]
The 5th--95th percentile band is $[74.2, 143.2]$ --- the log-normal right
skew that produces fat tails in returns.

\begin{keyinsight}
The convergence table is the law of large numbers made visible. At
$N = 100$ the estimate is off by 14\%; at $N = 100000$ it is off by
0.03\%. The standard error shrinks by exactly $\sqrt{10}$ for each
$10\times$ increase in paths --- the $1/\sqrt{N}$ law, the single most
important rate in computational finance.
\end{keyinsight}

\section{Exercises}

\begin{enumerate}
    \item \textbf{Control variates.} Use the analytical GBM mean $\E[S_T] =
      S_0 e^{rT}$ as a control variate: $\hat{C}_{\text{cv}} = \hat{C} -
      \beta(\hat{M} - \E[S_T])$, where $\hat{M}$ is the MC estimate of
      $\E[S_T]$ and $\beta = \text{Cov}(\text{payoff}, S_T)/\text{Var}(S_T)$.
      How much does the SE drop versus plain MC?
    \item \textbf{Asian option.} Price an arithmetic-average Asian call
      $\max(\bar S - K, 0)$ with $\bar S = \frac{1}{n}\sum S_{t_k}$. This
      requires time-stepping (no closed form). Compare SE to the European
      call for the same $N$.
    \item \textbf{Euler vs exact.} Discretise GBM by Euler ($S_{t+\Delta t} =
      S_t + \mu S_t \Delta t + \sigma S_t \sqrt{\Delta t} Z$) and compare
      the terminal distribution to the exact solution at $\Delta t = 1/252$.
      How large is the discretisation bias?
    \item \textbf{Confidence intervals.} For $N = 10000$, compute the 95\%
      confidence interval $\hat{C} \pm 1.96 \cdot \text{SE}$ and verify it
      covers the Black-Scholes price. Repeat 1000 times and check the
      coverage --- is it close to 95\%?
\end{enumerate}

\section{Key Takeaways}

\begin{itemize}
    \item Brownian motion is the continuous-time random walk: independent
      Gaussian increments, quadratic variation $[W]_t = t$. It is the
      building block of continuous-time finance.
    \item Geometric Brownian motion is the exponential of Brownian motion
      with drift, $S_T = S_0 \exp((\mu - \tfrac{1}{2}\sigma^2)T + \sigma W_T)$,
      and the Black-Scholes asset model. The closed form gives exact terminal
      distributions for any $\Delta t$ --- no Euler bias.
    \item The Poisson process counts rare events via exponential interarrival
      times; the Merton jump-diffusion superimposes it on GBM and reduces to
      GBM when $\lambda = 0$.
    \item Monte Carlo prices any payoff by simulating many paths and averaging
      the discounted payoff. The standard error shrinks as $1/\sqrt{N}$ ---
      slow but dimension-independent, the workhorse for high-dimensional
      derivatives.
    \item Antithetic variates pair $Z$ with $-Z$; for monotone payoffs the
      pair average cuts the SE per normal draw by roughly 30\%.
    \item The Black-Scholes formula (normal CDF via Abramowitz-Stegun erf,
      error $< 7.5 \times 10^{-8}$) anchors the Monte Carlo estimator and
      verifies convergence to the truth.
\end{itemize}

\vspace{1em}
\noindent\textbf{Next chapter}: Options pricing --- Black-Scholes Greeks,
implied volatility, and finite-difference sensitivity analysis.
\textit{Chapter content pending --- Phase 9 implementation.}

%----------------------------------------------------------------------------------------
% PART IV: DERIVATIVES AND PORTFOLIO (Expert)
%----------------------------------------------------------------------------------------

\pagelayout{wide}
\addpart{Derivatives and Portfolio Theory}
\pagelayout{margin}

\chapter{Options Pricing: Black-Scholes and Beyond}
\label{ch:options}
% \chapter{Options Pricing: Black-Scholes and Beyond}
\label{ch:options}

\begin{companioncode}{quant-options}
Code: \texttt{crates/quant-options/}\\
Dependencies: \crate{quant-core}, \crate{quant-stochastic}, \crate{thiserror}\\
Run: \texttt{cargo run -p quant-options --example greeks}
\end{companioncode}

\section{Learning Objectives}

This chapter builds the full Black-Scholes options toolkit on top of the
closed-form pricing primitives from \cref{ch:stochastic}: the Greeks
(Delta, Gamma, Vega, Theta, Rho) analytically and by finite difference,
and implied volatility via Newton's method with a bisection fallback.
After completing it you can:

\begin{itemize}
    \item State the Black-Scholes PDE and the risk-neutral derivation
    \item Compute all five first- and second-order Greeks in closed form
    \item Recover the same Greeks by finite difference for any pricing
      model (not just Black-Scholes)
    \item Solve for implied volatility from a market price using Newton's
      method, and explain why a bisection fallback is needed for deep
      in/out-of-the-money options
    \item Reproduce a synthetic volatility smile and explain its shape
\end{itemize}

\begin{keyinsight}
The Greeks are the partial derivatives of the option price with respect to
its inputs. Delta-hedging (holding $-\Delta$ shares against a long call)
eliminates first-order price risk; Gamma controls the residual hedging
error. Implied volatility inverts the Black-Scholes formula to recover
the market's view of future variance --- the ``volatility smile'' is the
single most studied phenomenon in derivatives.
\end{keyinsight}

\section{The Black-Scholes PDE}

\marginnote{%
\textbf{Black-Scholes PDE}\\[0.3em]
$\dfrac{\partial V}{\partial t} + \dfrac{1}{2}\sigma^2 S^2
\dfrac{\partial^2 V}{\partial S^2} + r S \dfrac{\partial V}{\partial S}
- r V = 0$
}

The Black-Scholes model assumes the underlying follows geometric Brownian
motion (see \cref{ch:stochastic}):
\[
dS_t = \mu S_t \, dt + \sigma S_t \, dW_t.
\]
By holding a continuously rebalanced portfolio of the option and
$-\Delta$ shares of stock, the residual risk is eliminated to first order.
No-arbitrage then requires the portfolio to earn the risk-free rate,
which yields the Black-Scholes PDE:
\[
\frac{\partial V}{\partial t} + \frac{1}{2}\sigma^2 S^2
\frac{\partial^2 V}{\partial S^2} + r S \frac{\partial V}{\partial S}
- r V = 0.
\]
Solving this PDE with the terminal payoff $\max(S_T - K, 0)$ for a call
(or $\max(K - S_T, 0)$ for a put) gives the closed forms
\[
C = S_0 \Phi(d_1) - K e^{-rT} \Phi(d_2), \qquad
P = K e^{-rT} \Phi(-d_2) - S_0 \Phi(-d_1),
\]
with
\[
d_1 = \frac{\ln(S_0/K) + (r + \frac{1}{2}\sigma^2)T}{\sigma\sqrt{T}},
\quad d_2 = d_1 - \sigma\sqrt{T}.
\]
These are the formulas we reuse from \crate{quant-stochastic::blackscholes}.

\section{The Greeks as Partial Derivatives}

\marginnote{%
\textbf{Greeks}\\[0.3em]
$\Delta = \dfrac{\partial C}{\partial S}$ \quad (spot)\\
$\Gamma = \dfrac{\partial^2 C}{\partial S^2}$ \quad (curvature)\\
$\mathcal{V} = \dfrac{\partial C}{\partial \sigma}$ \quad (vol)\\
$\Theta = \dfrac{\partial C}{\partial t}$ \quad (time)\\
$\rho  = \dfrac{\partial C}{\partial r}$ \quad (rate)
}

The Greeks are the sensitivities of the option price to its inputs. They
are the building blocks of hedging and risk management:

\begin{description}
    \item[Delta ($\Delta$)] $\partial C / \partial S$. The hedge ratio:
      holding $-\Delta$ shares against a long option removes first-order
      spot risk. A call has $\Delta \in (0, 1)$; a put $\Delta \in (-1,
      0)$.
    \item[Gamma ($\Gamma$)] $\partial^2 C / \partial S^2$. The rate of
      change of Delta. High Gamma means the hedge must rebalance
      frequently --- Gamma is largest at-the-money and shrinks in the
      wings. Identical for calls and puts.
    \item[Vega ($\mathcal{V}$)] $\partial C / \partial \sigma$. The
      sensitivity to volatility. Identical for calls and puts.
    \item[Theta ($\Theta$)] $\partial C / \partial t$. The time decay.
      Long options lose value as expiry approaches, so $\Theta < 0$ for
      long calls and puts.
    \item[Rho ($\rho$)] $\partial C / \partial r$. The sensitivity to the
      risk-free rate. Calls gain with $r$ (positive $\rho$); puts lose
      (negative $\rho$).
\end{description}

\section{Analytical Greeks}

\marginnote{%
\textbf{Analytical Greeks (call)}\\[0.3em]
$\Delta = \Phi(d_1)$\\
$\Gamma = \dfrac{\phi(d_1)}{S_0 \sigma \sqrt{T}}$\\
$\mathcal{V} = S_0 \phi(d_1) \sqrt{T}$\\
$\Theta = -\dfrac{S_0 \phi(d_1) \sigma}{2\sqrt{T}} - r K e^{-rT} \Phi(d_2)$\\
$\rho = K T e^{-rT} \Phi(d_2)$
}

Closed forms for the five Greeks follow from differentiating the
Black-Scholes price. Let $\phi(x) = (2\pi)^{-1/2} e^{-x^2/2}$ be the
standard normal PDF (the derivative of $\Phi$). We implement it inline
in \crate{quant-options::greeks}:

\begin{lstlisting}[language=Rust]
/// Standard normal PDF phi(x) = (1/sqrt(2 pi)) exp(-x^2 / 2).
pub fn normal_pdf(x: f64) -> f64 {
    const INV_SQRT_2PI: f64 = 0.398_942_280_401_432_7;
    INV_SQRT_2PI * (-0.5 * x * x).exp()
}
\end{lstlisting}

The analytical Greeks then drop out:

\begin{lstlisting}[language=Rust]
pub fn delta(s0, k, r, sigma, t, is_call: bool) -> f64 {
    let d1_val = d1(s0, k, r, sigma, t);
    let n_d1 = normal_cdf(d1_val);
    if is_call { n_d1 } else { n_d1 - 1.0 }
}

pub fn gamma(s0, k, r, sigma, t) -> f64 {
    let d1_val = d1(s0, k, r, sigma, t);
    normal_pdf(d1_val) / (s0 * sigma * t.sqrt())
}

pub fn vega(s0, k, r, sigma, t) -> f64 {
    let d1_val = d1(s0, k, r, sigma, t);
    s0 * normal_pdf(d1_val) * t.sqrt()
}
\end{lstlisting}

Theta and Rho differ between call and put; the call versions are
\[
\Theta_c = -\frac{S_0 \phi(d_1) \sigma}{2\sqrt{T}} - r K e^{-rT} \Phi(d_2),
\qquad
\rho_c = K T e^{-rT} \Phi(d_2).
\]
The put versions follow from put-call parity:
$\Theta_p = \Theta_c + r K e^{-rT}$ and
$\rho_p = -K T e^{-rT} \Phi(-d_2) = \rho_c - K T e^{-rT}$.

\section{Numerical Greeks by Finite Difference}

\marginnote{%
\textbf{Central difference}\\[0.3em]
$f'(x) \approx \dfrac{f(x+h) - f(x-h)}{2 h}$ \\[0.3em]
Error: $O(h^2)$
}

Analytical Greeks require a closed form. For models without a closed form
(local volatility, Heston, Monte Carlo) the same sensitivities are
recovered by finite difference. The central difference
\[
f'(x) \approx \frac{f(x+h) - f(x-h)}{2h}
\]
is second-order accurate: the symmetric form cancels the $O(h)$
truncation term, leaving $O(h^2)$. For Gamma (a second derivative), the
central second difference
\[
f''(x) \approx \frac{f(x+h) - 2 f(x) + f(x-h)}{h^2}
\]
is likewise $O(h^2)$.

\begin{lstlisting}[language=Rust]
pub fn delta_fd(s0, k, r, sigma, t, is_call, h) -> f64 {
    let up = price(s0 + h, k, r, sigma, t, is_call);
    let down = price(s0 - h, k, r, sigma, t, is_call);
    (up - down) / (2.0 * h)
}

pub fn gamma_fd(s0, k, r, sigma, t, h) -> f64 {
    let up = price(s0 + h, k, r, sigma, t, true);
    let mid = price(s0, k, r, sigma, t, true);
    let down = price(s0 - h, k, r, sigma, t, true);
    (up - 2.0 * mid + down) / (h * h)
}
\end{lstlisting}

\begin{keyinsight}
Finite difference is a sanity check on the analytical Greeks. If the two
disagree by more than a few basis points, one of them is wrong. With
$h = 10^{-4}$ for spot bumps and $h = 10^{-3}$ for the second derivative,
the analytical and numerical Greeks agree to $10^{-4}$ or better --- the
truncation and round-off errors are well balanced.
\end{keyinsight}

Theta uses a \emph{forward} difference, not central, because we cannot
step to $t - h < 0$ for a short-dated option:
\[
\Theta \approx \frac{C(t - h) - C(t)}{h}.
\]
The sign is the ``desk convention'': Theta is the negative of
$\partial C / \partial T$ (time to maturity), since long options lose
value as calendar time advances.

\section{Implied Volatility}

\marginnote{%
\textbf{Newton's method}\\[0.3em]
$\sigma_{n+1} = \sigma_n - \dfrac{C(\sigma_n) - C_{\text{mkt}}}{\mathcal{V}(\sigma_n)}$
}

Given a market price $C_{\text{mkt}}$, the implied volatility is the
$\sigma$ such that $C(S, K, r, \sigma, T) = C_{\text{mkt}}$. Because the
Black-Scholes price is monotone increasing in $\sigma$ (Vega is positive),
there is a unique solution. Newton's method is the textbook choice
because Vega is the exact derivative:
\[
\sigma_{n+1} = \sigma_n - \frac{C(\sigma_n) - C_{\text{mkt}}}{\mathcal{V}(\sigma_n)}.
\]
Newton converges quadratically when the initial guess is reasonable and
the option is not deep in/out of the money. When the option is deep ITM
or OTM, Vega collapses toward zero (the option price is intrinsic and
barely moves with $\sigma$), and the Newton step becomes ill-conditioned.
A bisection fallback on $[\sigma_{\min}, \sigma_{\max}]$ is then robust:
bisection is linear but guaranteed.

\begin{lstlisting}[language=Rust]
let newton_step = if v.abs() > VEGA_FLOOR {
    sigma - diff / v
} else {
    f64::NAN  // hand off to bisection
};

let next = if newton_step.is_finite()
    && newton_step > lo
    && newton_step < hi
{
    newton_step
} else {
    0.5 * (lo + hi)  // bisection
};
\end{lstlisting}

The initial guess uses the Brenner-Subrahmanyam (1988) approximation
$\sigma_0 \approx \sqrt{2\pi / T} \cdot C / S_0$, which is exact for
at-the-money options and bracketed by bisection if it diverges.

\section{Put-Call Parity and Implied Vol}

\marginnote{%
\textbf{Put-call parity}\\[0.3em]
$C - P = S_0 - K e^{-rT}$
}

Put-call parity says $C - P = S_0 - K e^{-rT}$. The implied vol that
reproduces a call price must therefore equal the implied vol that
reproduces the corresponding put price. The solver always inverts the
call formula; for a put we translate the put price into the equivalent
call price via $C = P + S_0 - K e^{-rT}$ and proceed. The IV is invariant
across calls and puts of the same strike and maturity.

\section{Results and Analysis}

\subsection{Greeks for an ATM Call}

Pricing an at-the-money call ($S_0 = K = 100$, $r = 0.05$, $\sigma = 0.2$,
$T = 1$) gives the call price $C = 10.4506$ (matching the Black-Scholes
benchmark from \cref{ch:stochastic}) and the Greeks:

\begin{center}
\begin{tabular}{lrrr}
\toprule
Greek & Analytical & Finite diff. & $|$diff$|$ \\
\midrule
$\Delta_c$ &  0.636831 &  0.636836 & $5.24 \times 10^{-6}$ \\
$\Delta_p$ & -0.363169 & -0.363164 & $5.24 \times 10^{-6}$ \\
$\Gamma$   &  0.018762 &  0.018763 & $1.34 \times 10^{-6}$ \\
$\mathcal{V}$ & 37.524035 & 37.523872 & $1.63 \times 10^{-4}$ \\
$\Theta_c$ & -6.414028 & -6.414142 & $1.15 \times 10^{-4}$ \\
$\Theta_p$ & -1.657881 & -1.657983 & $1.03 \times 10^{-4}$ \\
$\rho_c$   & 53.232483 & --- & --- \\
$\rho_p$   & -41.890459 & --- & --- \\
\bottomrule
\end{tabular}
\end{center}

The analytical and finite-difference Greeks agree to $10^{-4}$ or better
with $h = 10^{-4}$ for spot/vol bumps and $h = 10^{-3}$ for Gamma. Theta
matches to $10^{-4}$ with a forward difference in $t$.

\subsection{Implied Volatility Recovery}

Recovering $\sigma$ from the ATM call price $C = 10.4506$:

\begin{lstlisting}[language=bash]
$ cargo run -p quant-options --example implied_vol
Recovery (ATM call):
  Market price  = 10.450575
  Implied vol   = 0.20000000
  True vol      = 0.20000000
  |iv - true|   = 3.54e-14
\end{lstlisting}

The implied vol recovers the true $\sigma = 0.2$ to $\sim 10^{-14}$ ---
far below the typical bid-ask spread. The same IV is recovered from the
put price, as put-call parity demands.

\subsection{Volatility Smile}

A synthetic smile is generated by pricing calls under
$\sigma(K) = \sigma_0 + a \cdot (\ln(K/S_0))^2$ with $a = 0.5$, so the
ATM vol is $\sigma_0$ and the wings have higher vol (the classic equity
smile). Inverting the BS formula at each strike recovers the smile:

\begin{center}
\begin{tabular}{rrrr}
\toprule
$K$ & moneyness $K/S_0$ & market price & IV \\
\midrule
 70 & 0.700 & 33.981835 & 0.263609 \\
 80 & 0.800 & 24.965341 & 0.224897 \\
 90 & 0.900 & 16.851425 & 0.205550 \\
 95 & 0.950 & 13.390174 & 0.201316 \\
100 & 1.000 & 10.450575 & 0.200000 \\
105 & 1.050 &  8.068580 & 0.201190 \\
110 & 1.100 &  6.219914 & 0.204542 \\
120 & 1.200 &  3.823180 & 0.216621 \\
130 & 1.300 &  2.578807 & 0.234418 \\
\bottomrule
\end{tabular}
\end{center}

The IV is symmetric around the ATM strike in log-moneyness, exactly as
the model prescribes. In real markets the smile is not symmetric: equity
smiles are skewed (puts bid up by crash protection), FX smiles are more
symmetric, and rates smiles can be monotone. Fitting these is the job of
local-volatility, stochastic-volatility, and SVI parameterisations.

\begin{keyinsight}
The smile is the market's fingerprint. If the Black-Scholes model were
literally true, IV would be a flat line at $\sigma$. Real markets quote
options at many strikes with different IVs, and the shape of that curve
encodes the market's view of tail risk. Modelling the smile is the bridge
between the Black-Scholes world and the post-1987 derivatives industry.
\end{keyinsight}

\section{Exercises}

\begin{enumerate}
    \item \textbf{Black-Scholes with dividends.} Extend the BS formula to
      a continuous dividend yield $q$: replace $S_0$ with $S_0 e^{-qT}$
      everywhere, and re-derive the Greeks. Verify Delta becomes
      $e^{-qT} \Phi(d_1)$ for the call.

    \item \textbf{Local volatility.} Implement a simple local-volatility
      pricer where $\sigma(S, t)$ is a deterministic function of spot and
      time (e.g. $\sigma(S, t) = \sigma_0 + \alpha (S - S_0)^2$). Price an
      option by explicit Euler on a spot-time grid and compare the smile
      to the IV surface.

    \item \textbf{Heston model.} Implement the Heston stochastic-volatility
      SDE $dv_t = \kappa(\theta - v_t) dt + \xi \sqrt{v_t} dW^{(v)}_t$
      with Monte Carlo. Generate call prices for a strip of strikes and
      invert to recover the smile. Compare to the equity smile shape.

    \item \textbf{Vega-clipped bisection.} For an extreme deep-ITM call
      (e.g. $S_0 = 1000$, $K = 10$) show that the BS price equals the
      intrinsic value to machine precision, so IV is genuinely
      unrecoverable. Modify the solver to return \texttt{Ok(sigma\_min)}
      with a warning flag rather than attempting to recover.

    \item \textbf{Smile fitting.} Download a real option chain (e.g.
      SPX) and fit a quadratic-in-log-moneyness smile
      $\sigma(K) = \beta_0 + \beta_1 \ln(K/S_0) + \beta_2 (\ln(K/S_0))^2$
      via least squares. Report the residual and the skew coefficient
      $\beta_1$.

    \item \textbf{Newton convergence rate.} Instrument \func{implied\_vol}
      to print $|C(\sigma_n) - C_{\text{mkt}}|$ at each iteration. Verify
      quadratic convergence for ATM options (the error roughly squares
      each step) and linear convergence once bisection takes over.
\end{enumerate}

\section{Key Takeaways}

\begin{itemize}
    \item \textbf{Black-Scholes PDE} prices any payoff via continuous
      delta-hedging; the closed-form call and put follow from solving
      the PDE with the terminal payoff.
    \item \textbf{Greeks} are the partial derivatives of the price.
      Delta is the hedge ratio; Gamma the curvature; Vega the vol
      sensitivity; Theta the time decay; Rho the rate sensitivity. Gamma
      and Vega are identical for calls and puts.
    \item \textbf{Finite difference} recovers the Greeks for any pricing
      model, not just Black-Scholes. Central differences are $O(h^2)$;
      forward differences are needed for Theta (cannot step to $t < 0$).
    \item \textbf{Implied volatility} inverts the BS formula via Newton's
      method, with a bisection fallback when Vega collapses (deep ITM/OTM).
    \item \textbf{Put-call parity} implies the IV from a call equals the
      IV from the corresponding put --- a basic arbitrage check that
      real-market quotes must satisfy.
    \item \textbf{The volatility smile} is the market-implied IV curve
      across strikes; modelling it is the bridge from Black-Scholes to
      post-1987 derivatives practice.
\end{itemize}

\vspace{1em}
\noindent\textbf{Next chapter}: Portfolio optimization --- Markowitz
mean-variance frontier, the efficient frontier, and the capital asset
pricing model.
\textit{Chapter content pending --- Phase 10 implementation.}

\chapter{Portfolio Optimization: Markowitz and Beyond}
\label{ch:portfolio}
% \chapter{Portfolio Optimization: Markowitz and Beyond}
\label{ch:portfolio}

\begin{companioncode}{quant-portfolio}
Code: \texttt{crates/quant-portfolio/}\\
Dependencies: \crate{quant-core}, \crate{quant-timeseries}, \crate{thiserror}\\
Run: \texttt{cargo run -p quant-portfolio --example frontier}
\end{companioncode}

\section{Learning Objectives}

This chapter builds the Markowitz mean-variance framework on top of the
returns and covariance machinery from \cref{ch:returns} and \cref{ch:timeseries}:
the efficient frontier, the global minimum-variance portfolio, the tangency
portfolio, the capital market line, two-fund separation, and CAPM. After
completing it you can:

\begin{itemize}
    \item State the Markowitz mean-variance optimisation problem and its
      Lagrangian dual
    \item Compute the two-asset efficient frontier in closed form and
      identify the minimum-varance weight
    \item Solve the $n$-asset minimum-variance and target-return portfolios
      via the inverse of the covariance matrix
    \item Derive the tangency portfolio (maximum Sharpe) in closed form and
      draw the capital market line
    \item Apply two-fund separation: any efficient portfolio is a mix of the
      risk-free asset and the tangency portfolio
    \item Estimate CAPM beta and Jensen's alpha from a return series and
      interpret them on the security market line
    \item Compute historical Value-at-Risk and Conditional VaR
\end{itemize}

\begin{keyinsight}
Diversification is the only free lunch in finance. The efficient frontier
is the set of portfolios that maximise expected return for a given variance
(or minimise variance for a given return). The tangency portfolio is the
point on the frontier that maximises the Sharpe ratio; once a risk-free
asset exists, every efficient portfolio is a mix of the risk-free asset
and the tangency portfolio --- the celebrated \emph{two-fund separation}
theorem.
\end{keyinsight}

\section{The Markowitz Problem}

\marginnote{%
\textbf{Markowitz in one sentence}\\[0.3em]
For each target return, pick the portfolio of minimum variance that
achieves it. Sweep the target to trace the efficient frontier. The
recipe is 70 years old and still the starting point of every
portfolio construction problem.
}

Consider $n$ risky assets with expected return vector $\mu$ and covariance
matrix $\Sigma$. A portfolio is a vector of weights $w \in \mathbb{R}^n$
with budget constraint $1^\top w = 1$. The portfolio return and variance are
\[
\mu_p = w^\top \mu, \qquad \sigma_p^2 = w^\top \Sigma w.
\]
The Markowitz problem is to find, for each target return $\mu_\star$, the
portfolio that achieves it at minimum variance:
\[
\min_w \; w^\top \Sigma w
\quad \text{s.t.} \quad
1^\top w = 1, \;\; \mu^\top w = \mu_\star.
\]
The set of solutions, swept over $\mu_\star$, is the \emph{efficient
frontier}. Without a risk-free asset the frontier is a hyperbola in
$(\sigma_p, \mu_p)$ space; with a risk-free asset the relevant efficient set
becomes the straight line from $(0, r_f)$ through the tangency portfolio.

The \func{Portfolio} struct carries the three inputs (weights, expected
returns, covariance) together and exposes the three headline statistics:

\begin{lstlisting}[language=Rust]
pub struct Portfolio {
    pub weights: Vec<f64>,
    pub expected_returns: Vec<f64>,
    pub covariance: Vec<Vec<f64>>,
}

impl Portfolio {
    pub fn expected_return(&self) -> f64 {
        self.weights.iter()
            .zip(self.expected_returns.iter())
            .map(|(w, m)| w * m).sum()
    }
    pub fn variance(&self) -> f64 {
        quadratic_form(&self.weights, &self.covariance)
    }
    pub fn volatility(&self) -> f64 { self.variance().sqrt() }
    pub fn sharpe(&self, rf: f64) -> f64 {
        let vol = self.volatility();
        if vol == 0.0 { 0.0 } else { (self.expected_return() - rf) / vol }
    }
}
\end{lstlisting}

The quadratic form $w^\top \Sigma w$ is a one-liner once a matrix-vector
product is available; the Sharpe ratio guards against a zero-volatility
degenerate portfolio (all cash).

\section{Two-Asset Closed-Form Frontier}

\marginnote{%
\textbf{Two-asset frontier}\\[0.3em]
$\mu_p = w \mu_A + (1{-}w) \mu_B$\\[0.2em]
$\sigma_p^2 = w^2 \sigma_A^2 + (1{-}w)^2 \sigma_B^2 + 2 w (1{-}w) \rho \sigma_A \sigma_B$\\[0.3em]
\textbf{Min-variance weight}\\[0.2em]
$w_A^\star = \dfrac{\sigma_B^2 - \rho \sigma_A \sigma_B}{\sigma_A^2 + \sigma_B^2 - 2 \rho \sigma_A \sigma_B}$
}

For two assets the frontier is a single-parameter curve indexed by the
weight $w$ on asset A. The minimum-variance weight has a closed form that
falls out of setting $d\sigma_p^2 / dw = 0$:
\[
w_A^\star = \frac{\sigma_B^2 - \rho \sigma_A \sigma_B}
                 {\sigma_A^2 + \sigma_B^2 - 2 \rho \sigma_A \sigma_B}.
\]
Three boundary cases are worth knowing by heart:
\begin{itemize}
    \item $\rho = 0$: $w_A^\star = \sigma_B^2 / (\sigma_A^2 + \sigma_B^2)$
      --- the inverse-variance weights.
    \item $\rho = 1$ and $\sigma_A = \sigma_B$: any split has the same
      variance; the formula is degenerate ($0/0$), so we fall back to
      $w_A^\star = 1/2$.
    \item $\rho = -1$: a perfect hedge exists; the minimum variance is
      zero (a risk-free portfolio can be synthesised from two risky
      assets).
\end{itemize}

\begin{lstlisting}[language=Rust]
pub fn two_asset_min_variance_weight(
    var_a: f64, var_b: f64, cov_ab: f64,
) -> f64 {
    let denom = var_a + var_b - 2.0 * cov_ab;
    if denom.abs() < 1e-15 { return 0.5; }  // degenerate
    (var_b - cov_ab) / denom
}

pub fn two_asset_frontier_point(
    w: f64, mu_a: f64, mu_b: f64,
    var_a: f64, var_b: f64, cov_ab: f64,
) -> FrontierPoint {
    let mu_p = w * mu_a + (1.0 - w) * mu_b;
    let var_p = w * w * var_a
        + (1.0 - w) * (1.0 - w) * var_b
        + 2.0 * w * (1.0 - w) * cov_ab;
    FrontierPoint { expected_return: mu_p, volatility: var_p.sqrt() }
}
\end{lstlisting}

The degenerate case (equal variances, $\rho = 1$) returns $w = 1/2$ rather
than panicking on a $0/0$ --- any split is minimum-variance when the assets
are perfectly correlated and equally volatile.

\section{The $n$-Asset Lagrangian}

\marginnote{%
\textbf{Why closed form?}\\[0.3em]
With only equality constraints ($1^\top w = 1$, $\mu^\top w =
\mu_\star$) the KKT system is linear: $2\Sigma w = \lambda_1 1 +
\lambda_2 \mu$. Inverting $\Sigma$ once gives every frontier
portfolio as a linear combination of $\Sigma^{-1} 1$ and
$\Sigma^{-1} \mu$. Inequality constraints (no short sales, box
constraints) destroy this closed form and require a real QP solver.
}

For $n$ assets the Markowitz problem is a constrained quadratic program.
With no inequality constraints (short sales allowed), the KKT conditions
reduce to a linear system and everything has a closed form in terms of
$\Sigma^{-1}$.

\subsection{Global Minimum-Variance Portfolio}

Minimising $w^\top \Sigma w$ subject only to $1^\top w = 1$ has the
Lagrangian $\mathcal{L} = w^\top \Sigma w - \lambda (1^\top w - 1)$, whose
stationarity gives $\Sigma w = \lambda 1$ and the budget constraint
recovers $\lambda$. The closed form is
\[
w^\star_{\text{GMV}} = \frac{\Sigma^{-1} 1}{1^\top \Sigma^{-1} 1}.
\]

\subsection{Efficient Frontier at a Target Return}

Adding the return constraint $\mu^\top w = \mu_\star$ produces a second
Lagrange multiplier. Define the standard scalars
\[
a = 1^\top \Sigma^{-1} 1, \quad
b = 1^\top \Sigma^{-1} \mu, \quad
c = \mu^\top \Sigma^{-1} \mu, \quad
d = a c - b^2.
\]
Then the target-return efficient portfolio is
\[
w^\star(\mu_\star) = \frac{1}{d} \, \Sigma^{-1}
\bigl[ (c - b \mu_\star) 1 + (a \mu_\star - b) \mu \bigr],
\]
and the minimum variance achievable for that target is
\[
\sigma_p^2(\mu_\star) = \frac{a \mu_\star^2 - 2 b \mu_\star + c}{d}.
\]
This is a parabola in $(\mu_\star, \sigma_p^2)$ space, opening upward; its
tip is the global minimum-variance portfolio at $\mu_\star = b/a$ with
variance $1/a$.

\begin{lstlisting}[language=Rust]
pub fn min_variance_portfolio(
    expected_returns: &[f64],
    covariance: &[Vec<f64>],
) -> Result<Vec<f64>, PortfolioError> {
    let n = expected_returns.len();
    let sigma_inv = inverse(covariance)?;
    let ones = vec![1.0_f64; n];
    let sigma_inv_ones = matvec(&sigma_inv, &ones);
    let denom: f64 = sigma_inv_ones.iter().sum();
    Ok(sigma_inv_ones.iter().map(|v| v / denom).collect())
}

pub fn efficient_frontier_point(
    expected_returns: &[f64],
    covariance: &[Vec<f64>],
    mu_target: f64,
) -> Result<Vec<f64>, PortfolioError> {
    let sigma_inv = inverse(covariance)?;
    let ones = vec![1.0_f64; expected_returns.len()];
    let sigma_inv_ones = matvec(&sigma_inv, &ones);
    let sigma_inv_mu = matvec(&sigma_inv, expected_returns);
    let a: f64 = sigma_inv_ones.iter().sum();
    let b: f64 = sigma_inv_ones.iter()
        .zip(expected_returns.iter()).map(|(x, m)| x * m).sum();
    let c: f64 = expected_returns.iter()
        .zip(sigma_inv_mu.iter()).map(|(m, x)| m * x).sum();
    let d = a * c - b * b;
    let coef_ones = (c - b * mu_target) / d;
    let coef_mu = (a * mu_target - b) / d;
    Ok((0..expected_returns.len())
        .map(|i| coef_ones * sigma_inv_ones[i] + coef_mu * sigma_inv_mu[i])
        .collect())
}
\end{lstlisting}

Both functions hinge on the small Gaussian-elimination \func{inverse}
in \crate{quant-portfolio}'s \func{linalg} module --- the only linear
algebra primitive needed for the entire Markowitz framework. The budget
constraint is implicit in the closed form: the numerator of
\func{min\_variance\_portfolio} is $\Sigma^{-1} 1$ (unnormalised), and
dividing by $1^\top \Sigma^{-1} 1$ enforces $\sum w_i = 1$.

\section{The Tangency Portfolio}

\marginnote{%
\textbf{Tangency portfolio}\\[0.3em]
$w^\star_{\text{tan}} = \dfrac{\Sigma^{-1} (\mu - r_f 1)}{1^\top \Sigma^{-1} (\mu - r_f 1)}$\\[0.3em]
\textbf{Sharpe}\\[0.2em]
$\text{Sharpe} = \dfrac{\mu_p - r_f}{\sigma_p}$
}

With a risk-free rate $r_f$, every efficient portfolio is a mix of the
risk-free asset and a single risky portfolio: the \emph{tangency
portfolio}, the point on the efficient frontier that maximises the Sharpe
ratio. The closed form is
\[
w^\star_{\text{tan}} = \frac{\Sigma^{-1} (\mu - r_f 1)}{1^\top \Sigma^{-1} (\mu - r_f 1)}.
\]
The numerator is the direction that maximises the excess-return-to-variance
trade-off; the denominator normalises the budget to one. If all expected
returns equal $r_f$ the tangency is undefined (the excess return vector is
zero) and we surface an error.

\begin{lstlisting}[language=Rust]
let excess: Vec<f64> = expected_returns.iter().map(|m| m - rf).collect();
let sigma_inv = inverse(covariance)?;
let sigma_inv_excess = matvec(&sigma_inv, &excess);
let denom: f64 = sigma_inv_excess.iter().sum();
let weights: Vec<f64> = sigma_inv_excess.iter().map(|v| v / denom).collect();
\end{lstlisting}

\section{The Capital Market Line and Two-Fund Separation}

\marginnote{%
\textbf{Capital market line}\\[0.3em]
$\mu_p = r_f + \text{Sharpe}_{\text{tan}} \cdot \sigma_p$\\[0.3em]
\textbf{Two-fund separation}\\[0.2em]
$y = \sigma_p / \sigma_{\text{tan}}$\\
$w = y \, w_{\text{tan}} + (1{-}y) \, r_f$
}

The capital market line (CML) is the straight line in $(\sigma, \mu)$
space from $(0, r_f)$ through the tangency point:
\[
\mu_p = r_f + \text{Sharpe}_{\text{tan}} \cdot \sigma_p.
\]
Every point on the CML is achievable by mixing the risk-free asset with
the tangency portfolio. The fraction invested in the tangency portfolio
to achieve a target volatility $\sigma_p$ is $y = \sigma_p / \sigma_{\tan}$;
the rest, $1 - y$, is invested at $r_f$. This is the
\emph{two-fund separation} theorem: investors with different risk
preferences differ only in $y$, not in the composition of the risky
fund. The tangency portfolio is the market portfolio under CAPM
assumptions.

\begin{marginfigure}
    \centering
    \vspace{-3cm}
    % Markowitz Efficient Frontier with Tangency Portfolio and CML
% Usage: % Markowitz Efficient Frontier with Tangency Portfolio and CML
% Usage: \input{figures/efficient-frontier}
% Coordinates: scale_x = 20 (sigma in [0, 0.30] -> 0..6 cm)
%              scale_y = 50 (mu    in [0, 0.12] -> 0..6 cm)
\begin{tikzpicture}[scale=0.5]
    % Axes (clipped to chart area: x in [0, 6.5], y in [0, 6.5])
    \draw[thick, ->] (0,0) -- (6.8,0) node[right] {\footnotesize $\sigma_p$};
    \draw[thick, ->] (0,0) -- (0,6.8) node[above] {\footnotesize $\mu_p$};

    % Chart border (light, optional) to visualise the plot area
    % \draw[gray!30, very thin] (0,0) rectangle (6.5,6.5);

    % Tick labels on x-axis (volatility)
    \foreach \x/\xlab in {0/0.00, 1/0.05, 2/0.10, 3/0.15, 4/0.20, 5/0.25, 6/0.30} {
        \draw (\x,0) -- (\x,-0.1);
        \node[below] at (\x,-0.1) {\tiny \xlab};
    }
    % Tick labels on y-axis (expected return)
    \foreach \y/\ylab in {0/0.00, 1/0.02, 2/0.04, 3/0.06, 4/0.08, 5/0.10, 6/0.12} {
        \draw (0,\y) -- (-0.1,\y);
        \node[left] at (-0.1,\y) {\tiny \ylab};
    }

    % Risk-free rate point on the y-axis
    \fill[black] (0,1.0) circle (2pt);
    \node[left] at (-0.1,1.0) {\tiny $r_f$};

    % Efficient frontier (upper branch of the risky-assets hyperbola).
    % Starts at GMV (w = w_GMV) and rises through the tangency portfolio
    % to asset A and beyond (short B to leverage A).
    \draw[thick, blue] plot[smooth] coordinates {
        (3.33,4.23) (3.42,4.50) (3.53,4.65) (3.65,4.75)
        (4.00,5.00) (4.95,5.50)
    };

    % Inefficient lower branch (w < w_GMV): same variance as the upper
    % branch, but lower return. Runs from GMV down-right to asset B
    % (all-B portfolio) and beyond (short A to buy more B).
    \draw[thick, blue!40, dashed] plot[smooth] coordinates {
        (3.33,4.23) (3.39,4.00) (3.61,3.75) (3.94,3.50)
        (4.37,3.25) (4.87,3.00) (5.42,2.75) (6.00,2.50)
    };

    % Capital market line: from (0, rf) through tangency, clipped at y=6.5.
    % Slope in scaled coords: (4.65-1.0)/(3.53-0) = 1.034.
    % At y=6.5: x = (6.5-1.0)/1.034 = 5.318.
    \draw[thick, red] (0,1.0) -- (5.32, 6.5);
    \node[red, above, rotate=58] at (2.6, 3.7) {\footnotesize CML};

    % Asset points
    \fill[black] (4.00,5.00) circle (2pt);
    \node[above right] at (4.00,5.00) {\tiny Asset A};
    \fill[black] (6.00,2.50) circle (2pt);
    \node[below right] at (6.00,2.50) {\tiny Asset B};

    % Global minimum-variance portfolio
    \fill[black] (3.33,4.23) circle (2pt);
    \node[above left] at (3.33,4.23) {\tiny GMV};

    % Tangency portfolio
    \fill[red!70!black] (3.53,4.65) circle (2.5pt);
    \node[above right] at (3.53,4.65) {\tiny Tangency};

    % Legend in the lower-right empty area (below the inefficient branch).
    \node[draw, fill=white, inner sep=3pt, anchor=north east]
        at (6.6, 0.4) {%
        \begin{tabular}{@{}ll@{}}
            \textcolor{blue}{\rule[0.5ex]{8pt}{1pt}} & \tiny Efficient frontier \\
            \textcolor{blue!40}{\rule[0.5ex]{8pt}{1pt}} & \tiny Inefficient branch \\
            \textcolor{red}{\rule[0.5ex]{8pt}{1pt}} & \tiny Capital market line \\
        \end{tabular}};
\end{tikzpicture}
% Coordinates: scale_x = 20 (sigma in [0, 0.30] -> 0..6 cm)
%              scale_y = 50 (mu    in [0, 0.12] -> 0..6 cm)
\begin{tikzpicture}[scale=0.5]
    % Axes (clipped to chart area: x in [0, 6.5], y in [0, 6.5])
    \draw[thick, ->] (0,0) -- (6.8,0) node[right] {\footnotesize $\sigma_p$};
    \draw[thick, ->] (0,0) -- (0,6.8) node[above] {\footnotesize $\mu_p$};

    % Chart border (light, optional) to visualise the plot area
    % \draw[gray!30, very thin] (0,0) rectangle (6.5,6.5);

    % Tick labels on x-axis (volatility)
    \foreach \x/\xlab in {0/0.00, 1/0.05, 2/0.10, 3/0.15, 4/0.20, 5/0.25, 6/0.30} {
        \draw (\x,0) -- (\x,-0.1);
        \node[below] at (\x,-0.1) {\tiny \xlab};
    }
    % Tick labels on y-axis (expected return)
    \foreach \y/\ylab in {0/0.00, 1/0.02, 2/0.04, 3/0.06, 4/0.08, 5/0.10, 6/0.12} {
        \draw (0,\y) -- (-0.1,\y);
        \node[left] at (-0.1,\y) {\tiny \ylab};
    }

    % Risk-free rate point on the y-axis
    \fill[black] (0,1.0) circle (2pt);
    \node[left] at (-0.1,1.0) {\tiny $r_f$};

    % Efficient frontier (upper branch of the risky-assets hyperbola).
    % Starts at GMV (w = w_GMV) and rises through the tangency portfolio
    % to asset A and beyond (short B to leverage A).
    \draw[thick, blue] plot[smooth] coordinates {
        (3.33,4.23) (3.42,4.50) (3.53,4.65) (3.65,4.75)
        (4.00,5.00) (4.95,5.50)
    };

    % Inefficient lower branch (w < w_GMV): same variance as the upper
    % branch, but lower return. Runs from GMV down-right to asset B
    % (all-B portfolio) and beyond (short A to buy more B).
    \draw[thick, blue!40, dashed] plot[smooth] coordinates {
        (3.33,4.23) (3.39,4.00) (3.61,3.75) (3.94,3.50)
        (4.37,3.25) (4.87,3.00) (5.42,2.75) (6.00,2.50)
    };

    % Capital market line: from (0, rf) through tangency, clipped at y=6.5.
    % Slope in scaled coords: (4.65-1.0)/(3.53-0) = 1.034.
    % At y=6.5: x = (6.5-1.0)/1.034 = 5.318.
    \draw[thick, red] (0,1.0) -- (5.32, 6.5);
    \node[red, above, rotate=58] at (2.6, 3.7) {\footnotesize CML};

    % Asset points
    \fill[black] (4.00,5.00) circle (2pt);
    \node[above right] at (4.00,5.00) {\tiny Asset A};
    \fill[black] (6.00,2.50) circle (2pt);
    \node[below right] at (6.00,2.50) {\tiny Asset B};

    % Global minimum-variance portfolio
    \fill[black] (3.33,4.23) circle (2pt);
    \node[above left] at (3.33,4.23) {\tiny GMV};

    % Tangency portfolio
    \fill[red!70!black] (3.53,4.65) circle (2.5pt);
    \node[above right] at (3.53,4.65) {\tiny Tangency};

    % Legend in the lower-right empty area (below the inefficient branch).
    \node[draw, fill=white, inner sep=3pt, anchor=north east]
        at (6.6, 0.4) {%
        \begin{tabular}{@{}ll@{}}
            \textcolor{blue}{\rule[0.5ex]{8pt}{1pt}} & \tiny Efficient frontier \\
            \textcolor{blue!40}{\rule[0.5ex]{8pt}{1pt}} & \tiny Inefficient branch \\
            \textcolor{red}{\rule[0.5ex]{8pt}{1pt}} & \tiny Capital market line \\
        \end{tabular}};
\end{tikzpicture}
    \caption{Markowitz efficient frontier for two uncorrelated risky assets.}
    \label{fig:efficient-frontier}
\end{marginfigure}

\section{CAPM: Beta, Alpha, and the SML}

\marginnote{%
\textbf{Beta as a slope}\\[0.3em]
$\beta_i = \text{Cov}(R_i, R_m) / \text{Var}(R_m)$ is the slope of the
linear regression of $R_i$ on $R_m$. It is \emph{not} correlation ---
two assets can have $\beta = 1$ with very different correlations if
their idiosyncratic vol differs. Beta is a regression coefficient, so
it depends on the variance of the regressor.
}

The Capital Asset Pricing Model pins down the relationship between an
asset's expected excess return and its comovement with the market:
\[
\mathbb{E}[R_i] - r_f = \beta_i (\mathbb{E}[R_m] - r_f),
\qquad
\beta_i = \frac{\text{Cov}(R_i, R_m)}{\text{Var}(R_m)}.
\]
The \emph{security market line} (SML) is this linear relationship in
$(\beta, \mathbb{E}[R])$ space. An asset lying above the SML has positive
\emph{Jensen's alpha}:
\[
\alpha_i = \bar{R}_i - \bigl(r_f + \beta_i (\bar{R}_m - r_f)\bigr),
\]
which under CAPM should be zero in equilibrium. Real-world alpha is
nonzero because of risk-factor exposures CAPM does not capture (size,
value, momentum --- the Fama--French factors of \cref{ch:timeseries}).
Beta is a regression coefficient and can be estimated by OLS; here we
compute it directly from the covariance definition.

\begin{lstlisting}[language=Rust]
pub fn beta(asset: &[f64], market: &[f64]) -> Result<f64, PortfolioError> {
    let n = asset.len();
    let mean_a: f64 = asset.iter().sum::<f64>() / n as f64;
    let mean_m: f64 = market.iter().sum::<f64>() / n as f64;
    let mut cov = 0.0_f64;
    let mut var_m = 0.0_f64;
    for (a, m) in asset.iter().zip(market.iter()) {
        cov += (a - mean_a) * (m - mean_m);
        var_m += (m - mean_m) * (m - mean_m);
    }
    if var_m == 0.0 {
        return Err(PortfolioError::InvalidParam(
            "market variance is zero".into()));
    }
    Ok(cov / var_m)
}

pub fn alpha(asset: &[f64], market: &[f64], rf: f64) -> Result<f64, PortfolioError> {
    let b = beta(asset, market)?;
    let mean_a: f64 = asset.iter().sum::<f64>() / asset.len() as f64;
    let mean_m: f64 = market.iter().sum::<f64>() / market.len() as f64;
    Ok(mean_a - (rf + b * (mean_m - rf)))
}
\end{lstlisting}

The $(n - 1)$ denominators in sample covariance and sample variance
cancel, so the ratio is identical whether we use population or sample
moments. Beta is a pure ratio of comovement to market variance; alpha is
the gap between realised mean return and the SML prediction.

\section{Risk Metrics: Historical VaR and CVaR}

Value-at-Risk (VaR) at confidence $\alpha$ is the loss that is exceeded
with probability $1 - \alpha$. The historical estimator uses the empirical
quantile of the return series:
\[
\text{VaR}_\alpha = -Q_{1-\alpha}(R),
\]
where $Q_p$ is the $p$-th empirical quantile with linear interpolation
between adjacent order statistics. Conditional VaR (CVaR), also called
Expected Shortfall, is the average loss in the tail beyond the VaR level:
\[
\text{CVaR}_\alpha = -\mathbb{E}[R \,|\, R \le Q_{1-\alpha}(R)].
\]
CVaR is a coherent risk measure (sub-additive, convex) whereas VaR is
not. Both are non-parametric here: no Gaussian or Student-$t$ assumption
is imposed, so the tail shape is whatever the data says it is.

\begin{lstlisting}[language=Rust]
pub fn historical_var(returns: &[f64], confidence: f64) -> Result<f64, PortfolioError> {
    let p = 1.0 - confidence;
    let quantile = empirical_quantile(returns, p);
    Ok(-quantile)
}

pub fn historical_cvar(returns: &[f64], confidence: f64) -> Result<f64, PortfolioError> {
    let p = 1.0 - confidence;
    let quantile = empirical_quantile(returns, p);
    let tail: Vec<f64> = returns.iter()
        .filter(|r| **r <= quantile + 1e-12).copied().collect();
    let mean_tail: f64 = tail.iter().sum::<f64>() / tail.len() as f64;
    Ok(-mean_tail)
}

fn empirical_quantile(returns: &[f64], p: f64) -> f64 {
    let mut sorted: Vec<f64> = returns.to_vec();
    sorted.sort_by(|a, b| a.partial_cmp(b).unwrap_or(std::cmp::Ordering::Equal));
    let n = sorted.len();
    let pos = p * (n as f64 - 1.0);
    let lower = pos.floor() as usize;
    let upper = pos.ceil() as usize;
    if lower == upper { return sorted[lower]; }
    let frac = pos - lower as f64;
    sorted[lower] * (1.0 - frac) + sorted[upper] * frac
}
\end{lstlisting}

The quantile uses linear interpolation between adjacent order
statistics (the same convention as NumPy's default). The VaR is the
\emph{negative} of the low quantile (a positive loss); CVaR is the
negative of the mean of the tail below that quantile. The small
$10^{-12}$ slack in the filter guards against floating-point ties at
the boundary.

\section{Results and Analysis}

\subsection{Two-Asset Universe}

Throughout the chapter we use two uncorrelated risky assets:
$A$ with $\mu_A = 0.10$, $\sigma_A = 0.20$; $B$ with $\mu_B = 0.05$,
$\sigma_B = 0.30$; $\rho = 0$. The risk-free rate is $r_f = 0.02$.

\begin{center}
\begin{tabular}{lrr}
\toprule
Portfolio & $w_A$ & $w_B$ \\
\midrule
Asset A only            & 1.0000 & 0.0000 \\
Asset B only            & 0.0000 & 1.0000 \\
Global min-variance     & 0.6923 & 0.3077 \\
Tangency (max Sharpe)   & 0.8571 & 0.1429 \\
\bottomrule
\end{tabular}
\end{center}

The global minimum-variance portfolio places weight inversely
proportional to variance: $w_A = (1/0.04) / (1/0.04 + 1/0.09) = 0.6923$.
The tangency portfolio tilts further toward asset A because of its
higher expected return.

\subsection{Tangency Portfolio and CML}

\begin{center}
\begin{tabular}{lr}
\toprule
Quantity & Value \\
\midrule
$\mu_{\text{tan}}$     & 0.092857 \\
$\sigma_{\text{tan}}$  & 0.176705 \\
Sharpe$_{\text{tan}}$  & 0.412311 \\
\bottomrule
\end{tabular}
\end{center}

The capital market line $\mu = 0.02 + 0.4123 \cdot \sigma$ dominates the
risky-asset-only frontier everywhere above the tangency point. To achieve
$\sigma_p = 0.25$, for instance, an investor puts $y = 0.25 / 0.1767 \approx
1.415$ in the tangency portfolio (borrowing at $r_f$) and earns
$\mu_p = 0.02 + 0.4123 \cdot 0.25 = 0.1231$ --- better than any risky-only
portfolio at the same volatility.

\subsection{N-Asset Efficient Frontier (Lagrangian)}

Sweeping the target return from $0.05$ to $0.15$ recovers the parabolic
frontier. Up to $\mu_\star = 0.10$ the frontier is long-only; above that,
the Lagrangian short-sells asset B to finance more of asset A:

\begin{center}
\begin{tabular}{rrrrr}
\toprule
target & $w_A$ & $w_B$ & $\mu_p$ & $\sigma_p$ \\
\midrule
0.050 &  0.000 &  1.000 & 0.0500 & 0.3000 \\
0.070 &  0.400 &  0.600 & 0.0700 & 0.1970 \\
0.090 &  0.800 &  0.200 & 0.0900 & 0.1709 \\
0.100 &  1.000 &  0.000 & 0.1000 & 0.2000 \\
0.110 &  1.200 & $-$0.200 & 0.1100 & 0.2474 \\
0.130 &  1.600 & $-$0.600 & 0.1300 & 0.3672 \\
0.150 &  2.000 & $-$1.000 & 0.1500 & 0.5000 \\
\bottomrule
\end{tabular}
\end{center}

The minimum-variance target (the tip of the parabola) is at
$\mu_\star = b/a = 0.0846$ with $\sigma_{\text{GMV}} = 0.1664$ --- matching
the global minimum-variance portfolio from the closed form.

\subsection{CAPM Regression}

A synthetic asset is generated as $R_i = r_f + \beta (R_m - r_f) + \varepsilon$
with true $\beta = 1.2$, $r_f = 0.02$, market mean $0.08$, and Gaussian
noise of standard deviation $0.02$. Regressing 60 observations:

\begin{center}
\begin{tabular}{lr}
\toprule
Estimate & Value \\
\midrule
$\hat{\beta}$ (true 1.20) & 1.175155 \\
$\hat{\alpha}$            & $-$0.000196 \\
SML-predicted $\bar{R}_i$ & 0.091882 \\
Realised $\bar{R}_i$       & 0.091686 \\
\bottomrule
\end{tabular}
\end{center}

Jensen's alpha is essentially zero ($\sim 10^{-4}$), as expected for a
data-generating process that lies exactly on the SML. The beta estimate
is biased by about $0.025$ due to finite-sample noise; with 60
observations the standard error of $\hat{\beta}$ is roughly
$\sigma_\varepsilon / (\sigma_m \sqrt{n}) \approx 0.02 / (0.04 \cdot
7.75) \approx 0.065$, so the bias is well within one standard error.

\begin{keyinsight}
CAPM is a one-factor model: expected excess return is proportional to
beta. In real markets, beta alone explains only a fraction of the cross
section of expected returns --- the residual alpha is the door through
which Fama--French size/value, momentum, and other multi-factor models
walk. The next chapter extends CAPM to a multi-factor world via PCA and
Fama--French.
\end{keyinsight}

\section{Exercises}

\begin{enumerate}
    \item \textbf{Correlated assets.} Reproduce the two-asset frontier with
      $\rho = 0.5$, $\rho = -0.5$, and $\rho = -1$. Verify that $\rho = -1$
      admits a zero-variance portfolio and find its weight. Plot the
      three frontiers on the same axes.

    \item \textbf{Three-asset frontier.} Add a third asset with
      $\mu_C = 0.08$, $\sigma_C = 0.15$ and correlations $\rho_{AB} = 0$,
      $\rho_{AC} = 0.3$, $\rho_{BC} = -0.2$. Compute the global
      minimum-variance and tangency portfolios and compare the Sharpe to
      the two-asset case.

    \item \textbf{Long-only constraint.} The Lagrangian solution allows
      short sales. Implement a projected-gradient descent that enforces
      $w_i \ge 0$ and $1^\top w = 1$ (project onto the probability simplex
      after each gradient step). Compare the constrained and unconstrained
      frontiers for the two-asset universe when the target is very high
      ($\mu_\star = 0.15$).

    \item \textbf{Black-Litterman.} Implement the Black-Litterman reverse
      optimisation: start from the market-portfolio weights, assume a
      prior $\mu \propto \Sigma w_{\text{mkt}}$, and update with a view
      (e.g. ``asset A will return 0.15'') using Bayesian shrinkage.
      Compare the posterior tangency to the prior tangency.

    \item \textbf{Rolling CAPM.} Download a year of daily returns for an
      equity and a market index. Compute rolling 60-day betas and plot
      them. Identify regime changes (e.g. a stock that becomes more
      correlated with the market during a crash).

    \item \textbf{CVaR optimisation.} Replace variance with CVaR at the
      95\% level as the objective and minimise over a long-only simplex
      via grid search for the two-asset universe. Compare the optimal
      weights to the Markowitz minimum-variance weights and discuss when
      the two differ.
\end{enumerate}

\section{Key Takeaways}

\begin{itemize}
    \item \textbf{Markowitz} frames portfolio choice as mean-variance
      optimisation. The efficient frontier is the set of portfolios that
      are Pareto-optimal in $(\sigma_p, \mu_p)$.
    \item \textbf{Two-asset frontier} has a closed form indexed by a single
      weight; the minimum-variance weight is inverse-variance when
      $\rho = 0$.
    \item \textbf{$n$-asset Lagrangian} turns the constrained quadratic
      program into a linear system in $\Sigma^{-1}$. Both the global
      minimum-variance and the target-return efficient portfolios have
      closed forms.
    \item \textbf{Tangency portfolio} maximises the Sharpe ratio; its
      weights are $\Sigma^{-1} (\mu - r_f 1)$ normalised to budget one.
    \item \textbf{Two-fund separation} says every efficient portfolio is a
      mix of the risk-free asset and the tangency portfolio. Investors
      with different risk preferences differ only in $y$.
    \item \textbf{CAPM} pins the expected excess return to
      $\beta (\mathbb{E}[R_m] - r_f)$. Real-world alpha captures
      risk-factor exposures CAPM omits.
    \item \textbf{Historical VaR / CVaR} are non-parametric tail risk
      summaries; CVaR is coherent (sub-additive), VaR is not.
\end{itemize}

\vspace{1em}
\noindent\textbf{Next chapter}: Factor models --- principal component
analysis of returns, the Fama--French three-factor model, and the bridge
from CAPM's single-factor world to multi-factor asset pricing.
\textit{Chapter content pending --- Phase 11 implementation.}

\chapter{Factor Models and PCA}
\label{ch:factors}
% \chapter{Factor Attribution: PCA and Fama-French}
\label{ch:factors}

\begin{companioncode}{quant-factors}
Code: \texttt{crates/quant-factors/}\\
Dependencies: \crate{quant-core}, \crate{quant-timeseries}, \crate{quant-portfolio}, \crate{thiserror}\\
Run: \texttt{cargo run -p quant-factors --example pca\_demo}\\
Run: \texttt{cargo run -p quant-factors --example fama\_french}
\end{companioncode}

\section{Learning Objectives}

This chapter builds the factor attribution framework on top of the
covariance and CAPM machinery from \cref{ch:timeseries} and
\cref{ch:portfolio}: principal component analysis, the Fama-French
3-factor model, and risk decomposition. After completing it you can:

\begin{itemize}
    \item Explain why a single-factor model (CAPM) is insufficient and how
      multi-factor models capture additional systematic risk
    \item Compute the dominant eigenvalue and eigenvector of a symmetric
      matrix via the power method
    \item Deflate a matrix to find subsequent eigenpairs and recover the
      full spectrum
    \item Perform PCA on a returns matrix: centre, form the covariance,
      decompose, and project
    \item Quantify the explained variance ratio and the reconstruction
      error as a function of the number of retained components
    \item Estimate Fama-French 3-factor exposures (alpha, beta\_mkt,
      beta\_smb, beta\_hml) via OLS regression
    \item Decompose portfolio variance into systematic and idiosyncratic
      components
\end{itemize}

\begin{keyinsight}
CAPM says the market is the only systematic risk factor. Real returns
have multiple drivers: size (SMB), value (HML), momentum, quality. PCA
extracts the dominant statistical factors from the covariance matrix
without economic labels; the Fama-French model names three specific
factors with economic interpretation. The risk decomposition
$\sigma_p^2 = \text{systematic} + \text{idiosyncratic}$ tells you how
much of your portfolio's risk comes from factors you can hedge versus
asset-specific noise you must diversify away.
\end{keyinsight}

\section{Why Multi-Factor?}

\marginnote{%
\textbf{CAPM's blind spot}\\[0.3em]
CAPM says the market is the only systematic risk factor. Empirically,
small stocks and value stocks earn higher average returns than CAPM
predicts --- they carry systematic risk that the market beta does not
absorb. Multi-factor models name those additional risks so they can be
measured, hedged, and priced.
}

The CAPM of \cref{ch:portfolio} is a single-factor model: the expected
return of an asset is linear in its beta to the market portfolio alone.
Empirically, CAPM betas explain a modest fraction of the
cross-sectional variation in stock returns. Fama and French (1992)
showed that two additional factors --- \emph{size} (Small Minus Big,
SMB) and \emph{value} (High Minus Low, HML) --- capture systematic risk
that the market beta misses. Small stocks and value stocks earn higher
average returns than CAPM predicts, suggesting they carry additional
systematic risk that the market factor does not absorb.

\begin{marginfigure}
    \centering
    \vspace{-2cm}
    % Scree plot: explained variance per principal component.
% Usage: % Scree plot: explained variance per principal component.
% Usage: \input{figures/scree-plot}
% Scale: x = 1 cm per PC (3 PCs), y = 6 cm per unit EVR (0..1.0 -> 0..6 cm)
\begin{tikzpicture}[scale=0.5]
    % Axes
    \draw[thick, ->] (0,0) -- (4.5,0) node[right] {\footnotesize PC};
    \draw[thick, ->] (0,0) -- (0,6.8) node[above] {\footnotesize EVR};

    % Y-axis ticks
    \foreach \y/\ylab in {0/0.0, 1/0.2, 2/0.4, 3/0.6, 4/0.8, 5/1.0} {
        \draw (0,\y) -- (-0.1,\y);
        \node[left] at (-0.1,\y) {\tiny \ylab};
    }
    % X-axis labels
    \foreach \x/\xlab in {1/PC1, 2/PC2, 3/PC3} {
        \node[below] at (\x, -0.1) {\tiny \xlab};
    }

    % Bars: EVR = 0.9608, 0.0389, 0.0003 (scaled by 6)
    % PC1
    \fill[blue!70] (0.6, 0) rectangle (1.4, 5.765);
    \node[above] at (1.0, 5.765) {\tiny 0.961};
    % PC2
    \fill[blue!70] (1.6, 0) rectangle (2.4, 0.233);
    \node[above] at (2.0, 0.233) {\tiny 0.039};
    % PC3 (tiny, but still drawn)
    \fill[blue!70] (2.6, 0) rectangle (3.4, 0.002);
    \node[above] at (3.0, 0.15) {\tiny 0.0003};

    % Cumulative line (optional, shows the "knee")
    % Cumulative: 0.9608, 0.9997, 1.0000 (scaled by 6)
    \draw[thick, red!70!black, mark=*, mark size=1.5pt] plot coordinates {
        (1.0, 5.765) (2.0, 5.998) (3.0, 6.000)
    };
    \node[red!70!black, right] at (3.2, 6.0) {\tiny Cum.};

    % Grid lines (light)
    \foreach \y in {1,2,3,4,5} {
        \draw[gray!20, very thin] (0,\y) -- (4.0,\y);
    }
\end{tikzpicture}
% Scale: x = 1 cm per PC (3 PCs), y = 6 cm per unit EVR (0..1.0 -> 0..6 cm)
\begin{tikzpicture}[scale=0.5]
    % Axes
    \draw[thick, ->] (0,0) -- (4.5,0) node[right] {\footnotesize PC};
    \draw[thick, ->] (0,0) -- (0,6.8) node[above] {\footnotesize EVR};

    % Y-axis ticks
    \foreach \y/\ylab in {0/0.0, 1/0.2, 2/0.4, 3/0.6, 4/0.8, 5/1.0} {
        \draw (0,\y) -- (-0.1,\y);
        \node[left] at (-0.1,\y) {\tiny \ylab};
    }
    % X-axis labels
    \foreach \x/\xlab in {1/PC1, 2/PC2, 3/PC3} {
        \node[below] at (\x, -0.1) {\tiny \xlab};
    }

    % Bars: EVR = 0.9608, 0.0389, 0.0003 (scaled by 6)
    % PC1
    \fill[blue!70] (0.6, 0) rectangle (1.4, 5.765);
    \node[above] at (1.0, 5.765) {\tiny 0.961};
    % PC2
    \fill[blue!70] (1.6, 0) rectangle (2.4, 0.233);
    \node[above] at (2.0, 0.233) {\tiny 0.039};
    % PC3 (tiny, but still drawn)
    \fill[blue!70] (2.6, 0) rectangle (3.4, 0.002);
    \node[above] at (3.0, 0.15) {\tiny 0.0003};

    % Cumulative line (optional, shows the "knee")
    % Cumulative: 0.9608, 0.9997, 1.0000 (scaled by 6)
    \draw[thick, red!70!black, mark=*, mark size=1.5pt] plot coordinates {
        (1.0, 5.765) (2.0, 5.998) (3.0, 6.000)
    };
    \node[red!70!black, right] at (3.2, 6.0) {\tiny Cum.};

    % Grid lines (light)
    \foreach \y in {1,2,3,4,5} {
        \draw[gray!20, very thin] (0,\y) -- (4.0,\y);
    }
\end{tikzpicture}
    \caption{Scree plot: explained variance per principal component.
    The first PC captures 96\% of the total variance; the second adds
    3.9\%; the third adds only 0.3\%.}
    \label{fig:scree-plot}
\end{marginfigure}

A \emph{factor model} expresses each asset's excess return as a linear
combination of $K$ factor returns plus an idiosyncratic residual:
\[
  R_i - R_f = \alpha_i + \sum_{k=1}^{K} \beta_{ik} F_k + \varepsilon_i
\]
where $F_k$ is the return of factor $k$, $\beta_{ik}$ is the loading of
asset $i$ on factor $k$, and $\varepsilon_i$ is the asset-specific
residual. The factors are assumed to be uncorrelated with the
residuals, and the residuals are assumed uncorrelated across assets.

\section{Eigenvalue Decomposition: The Power Method}

The spectral theorem guarantees that every real symmetric matrix $A$
has an orthonormal basis of eigenvectors with real eigenvalues:
$A = V \Lambda V^\top$. The \emph{power method} finds the dominant
eigenpair $(\lambda_1, v_1)$ by repeated matrix-vector multiplication:
\[
  v^{(k+1)} = \frac{A v^{(k)}}{\|A v^{(k)}\|} \;\to\; v_1, \qquad
  \lambda_1 = v_1^\top A v_1 \quad \text{(Rayleigh quotient)}
\]
Starting from a random vector, each iteration amplifies the component
along $v_1$ by a factor $|\lambda_1 / \lambda_2|$ relative to the next
eigenvalue, so convergence is geometric with rate $|\lambda_2/\lambda_1|$.

\marginnote{%
\textbf{Power method intuition}\\[0.3em]
Repeated multiplication by $A$ amplifies the dominant eigendirection
the most. After $k$ steps, $v^{(k)} \propto \lambda_1^k v_1 +
\lambda_2^k v_2 + \ldots$ The ratio $|\lambda_2/\lambda_1|$ controls
convergence: well-separated spectra converge in a handful of
iterations; clustered eigenvalues are slow.
}

A sign-flip artefact arises when $\lambda_1 < 0$ (or when the initial
vector is exactly anti-aligned with the previous iterate): $v^{(k+1)}$
points opposite to $v^{(k)}$ while still converging in direction. We
fix this by aligning each iterate with its predecessor before testing
convergence, which makes the residual $\|v^{(k+1)} - v^{(k)}\|$
monotone non-increasing.

\begin{lstlisting}[language=Rust, caption={Power method for the dominant eigenpair.}]
pub fn power_method(
    matrix: &[Vec<f64>],
    max_iter: usize,
    tol: f64,
) -> Result<(f64, Vec<f64>), FactorError> {
    let n = matrix.len();
    let mut v = vec![1.0 / (n as f64).sqrt(); n];
    let mut lambda_prev = f64::INFINITY;
    for _ in 0..max_iter {
        let w = matvec_sym(matrix, &v);
        let norm = l2_norm(&w);
        let mut v_new: Vec<f64> = w.iter().map(|&x| x / norm).collect();
        // Align sign with the previous iterate to avoid spurious flips.
        let dot: f64 = v_new.iter().zip(v.iter()).map(|(a, b)| a * b).sum();
        if dot < 0.0 { for x in &mut v_new { *x = -*x; } }
        let av = matvec_sym(matrix, &v_new);
        let lambda: f64 = v_new.iter().zip(av.iter()).map(|(a, b)| a * b).sum();
        let delta = l2_norm(&v_new.iter().zip(v.iter())
            .map(|(a, b)| a - b).collect::<Vec<_>>());
        v = v_new;
        if (lambda - lambda_prev).abs() < tol && delta < tol * 10.0 {
            lambda_prev = lambda;
            break;
        }
        lambda_prev = lambda;
    }
    Ok((lambda_prev, v))
}
\end{lstlisting}

\section{Deflation and the Full Spectrum}

\marginnote{%
\textbf{Deflation trade-off}\\[0.3em]
Deflation is exact in infinite precision: $A' = A - \lambda v v^\top$
removes the eigendirection $v$ cleanly. In floating point, each
deflation step accumulates rounding error, so the smaller eigenvalues
are less accurate. For a well-conditioned $3 \times 3$ covariance,
the top two eigenpairs reach $10^{-6}$; the third may carry
$10^{-3}$ relative error. Use a QR/SVD solver if you need all
eigenpairs to full precision.
}

Once the dominant eigenpair is found, \emph{deflation} removes it from
the matrix:
\[
  A' = A - \lambda_1 \, v_1 v_1^\top
\]
For a symmetric matrix, $v_1 v_1^\top$ is the projector onto the
$e_1$ direction, so $A' v_1 = 0$ and the spectrum of $A'$ is
$\{0, \lambda_2, \lambda_3, \ldots\}$. Repeating the power method on
$A'$ yields $(\lambda_2, v_2)$, and so on.

\begin{lstlisting}[language=Rust, caption={Deflation and the top-$k$ eigenpairs.}]
pub fn deflate(
    matrix: &[Vec<f64>],
    eigenvalue: f64,
    eigenvector: &[f64],
) -> Vec<Vec<f64>> {
    let n = matrix.len();
    let mut out = vec![vec![0.0; n]; n];
    for i in 0..n {
        for j in 0..n {
            out[i][j] = matrix[i][j] - eigenvalue * eigenvector[i] * eigenvector[j];
        }
    }
    out
}

pub fn top_k_eigen(matrix: &[Vec<f64>], k: usize)
    -> Result<(Vec<f64>, Vec<Vec<f64>>), FactorError>
{
    let mut current = matrix.to_vec();
    let mut eigenvalues = Vec::with_capacity(k);
    let mut eigenvectors = Vec::with_capacity(k);
    for _ in 0..k {
        let (lambda, v) = power_method(&current, 0, 0.0)?;
        eigenvalues.push(lambda);
        eigenvectors.push(v.clone());
        current = deflate(&current, lambda, &v);
    }
    Ok((eigenvalues, eigenvectors))
}
\end{lstlisting}

The accumulated floating-point error in deflation limits the accuracy
of the smaller eigenvalues. For a well-conditioned matrix with
well-separated eigenvalues, the top two or three eigenpairs are
accurate to $10^{-6}$; the last eigenpair of a near-singular matrix
may carry relative error of $10^{-3}$ or worse.

\section{PCA: Covariance Eigendecomposition}

\marginnote{%
\textbf{Why centre?}\\[0.3em]
Subtracting the per-column mean before forming $\Sigma$ removes the
``level'' component. Without centring, the first PC would just pick up
the mean return direction; with centring, the first PC captures the
direction of maximum \emph{variance} around that mean --- which is
what we want for risk analysis.
}

\marginnote{%
\textbf{EVR denominator}\\[0.3em]
Use $\text{tr}(\Sigma) = \sum_{j=1}^{N} \lambda_j$, the \emph{full}
covariance trace --- not $\sum_{j=1}^{k} \lambda_j$ over retained
components only. With the full trace, $\sum_i \text{EVR}_i < 1$ when
$k < N$, which is the correct statement: retained components explain
\emph{this fraction} of total variance.
}

Given a returns matrix $X$ ($T$ observations $\times$ $N$ assets), PCA
centres the data by subtracting the per-column mean, forms the sample
covariance $\Sigma = \frac{1}{T-1} \tilde{X}^\top \tilde{X}$, and
decomposes it as $\Sigma = V \Lambda V^\top$. The principal components
are the columns of $V$ (the eigenvectors), and the explained variance
ratio of component $i$ is
\[
  \text{EVR}_i = \frac{\lambda_i}{\text{tr}(\Sigma)}
  = \frac{\lambda_i}{\sum_{j=1}^{N} \lambda_j}
\]
where $\text{tr}(\Sigma) = \sum_j \lambda_j$ is the total variance.
The cumulative explained variance after $k$ components is
$\sum_{i=1}^{k} \text{EVR}_i$.

\begin{lstlisting}[language=Rust, caption={PCA fit: centre, covariance, eigendecompose.}]
pub fn pca(returns: &[Vec<f64>], n_components: usize)
    -> Result<PcaResult, FactorError>
{
    let (t, n) = (returns.len(), returns[0].len());
    let mean: Vec<f64> = (0..n)
        .map(|j| returns.iter().map(|r| r[j]).sum::<f64>() / t as f64)
        .collect();
    let mut cov = vec![vec![0.0; n]; n];
    for row in returns {
        for i in 0..n {
            let di = row[i] - mean[i];
            for j in 0..n { cov[i][j] += di * (row[j] - mean[j]); }
        }
    }
    let denom = (t - 1) as f64;
    for row in &mut cov { for cell in row.iter_mut() { *cell /= denom; } }
    let (eigenvalues, eigenvectors) = top_k_eigen(&cov, n_components)?;
    let trace: f64 = (0..n).map(|i| cov[i][i]).sum();
    let explained_variance_ratio: Vec<f64> =
        eigenvalues.iter().map(|&l| l / trace).collect();
    Ok(PcaResult { eigenvalues, eigenvectors, explained_variance_ratio,
        cumulative_variance: /* ... */, mean })
}
\end{lstlisting}

\begin{lstlisting}[language=Rust, caption={Project and reconstruct.}]
pub fn pca_transform(returns: &[Vec<f64>], eigenvectors: &[Vec<f64>],
                    mean: &[f64]) -> Vec<Vec<f64>> {
    returns.iter().map(|row| {
        let centred: Vec<f64> = row.iter().zip(mean.iter())
            .map(|(x, m)| x - m).collect();
        eigenvectors.iter().map(|pc|
            pc.iter().zip(centred.iter()).map(|(a, b)| a * b).sum()
        ).collect()
    }).collect()
}

pub fn pca_reconstruct(scores: &[Vec<f64>], eigenvectors: &[Vec<f64>],
                       mean: &[f64]) -> Vec<Vec<f64>> {
    scores.iter().map(|row| {
        let mut out = mean.to_vec();
        for (k, score) in row.iter().enumerate() {
            for (out_j, ev) in out.iter_mut().zip(eigenvectors[k].iter()) {
                *out_j += score * ev;
            }
        }
        out
    }).collect()
}
\end{lstlisting}

\subsection{Demo Results}

Running \texttt{pca\_demo} on 12 synthetic observations of 3 assets
yields the following decomposition:

\begin{center}
\begin{tabular}{lrrr}
\toprule
PC & Eigenvalue & EVR & Cumulative \\
\midrule
1 & $1.96 \times 10^{-4}$ & 0.9608 & 0.9608 \\
2 & $7.94 \times 10^{-6}$ & 0.0389 & 0.9997 \\
3 & $7.00 \times 10^{-8}$ & 0.0003 & 1.0000 \\
\bottomrule
\end{tabular}
\end{center}

The top eigenvector $v_1 \approx (0.892, 0.452, -0.022)$ shows that the
first PC is essentially a long-asset-1, long-asset-2 position ---
the dominant source of co-movement. The reconstruction error (SSE)
decreases from $8.8 \times 10^{-5}$ (1 component) to $7.2 \times 10^{-7}$
(2 components) to $\approx 0$ (3 components), confirming the
$1/\sqrt{k}$ intuition.

\section{Fama-French 3-Factor Regression}

The Fama-French model regresses asset excess returns on three factors:
\[
  R_i - R_f = \alpha_i + \beta_{\text{Mkt}} (R_m - R_f)
  + \beta_{\text{SMB}} \cdot \text{SMB}
  + \beta_{\text{HML}} \cdot \text{HML}
  + \varepsilon_i
\]
The design matrix is $X = [1, \text{Mkt-Rf}, \text{SMB}, \text{HML}]$
and the OLS fit reuses the Gaussian-elimination solver from
\crate{quant-timeseries}.

\marginnote{%
\textbf{Alpha = skill or noise?}\\[0.3em]
The intercept $\alpha_i$ is the average return left over after the
three factors explain what they can. In-sample, any asset has a
non-zero alpha by construction; the question is whether it is
statistically different from zero out-of-sample. In a well-specified
factor model, alphas should be small and mostly insignificant.
}

\marginnote{%
\textbf{Partial vs simple regression}\\[0.3em]
$\beta_{\text{Mkt}}$ from FF3 is a \emph{partial} coefficient: it
measures market exposure \emph{after} controlling for SMB and HML. It
differs from the CAPM beta (a simple regression on the market alone)
whenever the factors are correlated. Do not expect FF3 $\beta_{\text{Mkt}}$
to equal CAPM $\beta$.
}

\begin{lstlisting}[language=Rust, caption={FF3 regression via the OLS machinery.}]
pub fn ff3_regression(asset_excess: &[f64], factors: &[Vec<f64>])
    -> Result<FF3Exposure, FactorError>
{
    let x: Vec<Vec<f64>> = factors.iter()
        .map(|row| vec![1.0, row[0], row[1], row[2]])
        .collect();
    let fit = ols(&x, asset_excess)
        .map_err(|e| FactorError::Singular(e.to_string()))?;
    let n = asset_excess.len() as f64;
    let dof = (n - 4.0).max(1.0);
    let residual_var = fit.residuals.iter()
        .map(|e| e * e).sum::<f64>() / dof;
    Ok(FF3Exposure {
        alpha: fit.coeffs[0],
        beta_mkt: fit.coeffs[1],
        beta_smb: fit.coeffs[2],
        beta_hml: fit.coeffs[3],
        r_squared: fit.r_squared,
        residual_var,
    })
}
\end{lstlisting}

\subsection{Demo Results}

Running \texttt{fama\_french} on 100 simulated observations with the
DGP $\alpha = 0.001$, $\beta_{\text{Mkt}} = 1.2$,
$\beta_{\text{SMB}} = 0.4$, $\beta_{\text{HML}} = -0.3$ plus noise:

\begin{center}
\begin{tabular}{lrr}
\toprule
Parameter & True & Estimated \\
\midrule
$\alpha$ & 0.001 & 0.001003 \\
$\beta_{\text{Mkt}}$ & 1.2 & 1.2008 \\
$\beta_{\text{SMB}}$ & 0.4 & 0.4016 \\
$\beta_{\text{HML}}$ & $-0.3$ & $-0.2977$ \\
$R^2$ (FF3) & --- & 0.9893 \\
$R^2$ (CAPM) & --- & 0.9391 \\
\bottomrule
\end{tabular}
\end{center}

The 3-factor model's $R^2$ (98.9\%) beats the single-factor CAPM's
$R^2$ (93.9\%) by 5.3 percentage points --- the incremental explanatory
power of SMB and HML beyond the market.

\section{Risk Attribution}

\marginnote{%
\textbf{Systematic vs idiosyncratic}\\[0.3em]
Systematic risk comes from factors you can hedge (market, size, value)
--- it is shared across assets and does not diversify away.
Idiosyncratic risk is asset-specific noise; it does diversify away as
$N$ grows. A well-diversified portfolio has $\sigma_p^2$ dominated by
the systematic term.
}

Given portfolio weights $w$, factor loadings $B$ ($N \times K$), factor
covariance $\Sigma_F$ ($K \times K$), and per-asset idiosyncratic
variances $D$ (diagonal), the portfolio variance decomposes as
\[
  \sigma_p^2 = \underbrace{w^\top B \, \Sigma_F \, B^\top w}_{\text{systematic}}
  + \underbrace{w^\top D \, w}_{\text{idiosyncratic}}
\]
The portfolio's factor exposure is $f_p = B^\top w$ (a $K$-vector);
the systematic variance is $f_p^\top \Sigma_F f_p$; the per-factor
contribution is the diagonal approximation
$\text{contrib}_k = f_{p,k}^2 \, \Sigma_{F,kk}$.

\begin{lstlisting}[language=Rust, caption={Risk attribution.}]
pub fn risk_attribution(
    weights: &[f64],
    factor_loadings: &[Vec<f64>],
    factor_covariance: &[Vec<f64>],
    residual_variances: &[f64],
) -> Result<RiskAttribution, FactorError> {
    let n = weights.len();
    let k = factor_loadings[0].len();
    // Portfolio factor exposures: f_p[k] = sum_i w_i * beta_ik.
    let mut f_p = vec![0.0; k];
    for (i, w_i) in weights.iter().enumerate() {
        for (kk, &beta) in factor_loadings[i].iter().enumerate() {
            f_p[kk] += w_i * beta;
        }
    }
    // Systematic: f_p' Sigma_F f_p.
    let mut systematic = 0.0;
    for i in 0..k {
        for j in 0..k {
            systematic += f_p[i] * factor_covariance[i][j] * f_p[j];
        }
    }
    // Idiosyncratic: sum_i w_i^2 * D_i.
    let idiosyncratic: f64 = weights.iter()
        .zip(residual_variances.iter())
        .map(|(w, d)| w * w * d).sum();
    let factor_contributions: Vec<f64> = (0..k)
        .map(|kk| f_p[kk] * f_p[kk] * factor_covariance[kk][kk])
        .collect();
    Ok(RiskAttribution {
        total_variance: systematic + idiosyncratic,
        systematic_variance: systematic,
        idiosyncratic_variance: idiosyncratic,
        factor_contributions,
    })
}
\end{lstlisting}

For the FF3 demo, the risk attribution on the single-asset portfolio
($w = 1$) gives:

\begin{center}
\begin{tabular}{lr}
\toprule
Component & Variance \\
\midrule
Total & $8.94 \times 10^{-5}$ \\
Systematic & $8.84 \times 10^{-5}$ (98.9\%) \\
Idiosyncratic & $9.85 \times 10^{-7}$ (1.1\%) \\
\midrule
Mkt contribution & $8.41 \times 10^{-5}$ \\
SMB contribution & $3.40 \times 10^{-6}$ \\
HML contribution & $8.54 \times 10^{-7}$ \\
\bottomrule
\end{tabular}
\end{center}

The market factor dominates the systematic variance (95.1\% of
total), consistent with the large market beta ($\beta_{\text{Mkt}} = 1.2$)
and the large market variance.

\section{Exercises}

\begin{enumerate}
    \item Add a 4th and 5th factor (momentum, quality) to the FF
      regression and compare $R^2$ to the 3-factor model.
    \item Implement the Jacobi eigenvalue algorithm as an alternative
      to the power method and compare convergence on a $5 \times 5$
      covariance matrix.
    \item Run PCA on a real equity returns matrix (e.g., 10 stocks,
      252 days) and interpret the first three principal components
      economically.
    \item Compute the rolling 60-day beta of a stock to the market
      factor and plot the time series.
    \item Construct a long-short portfolio that is neutral to the
      market ($\beta_{\text{Mkt}} = 0$) but exposed to SMB, and
      decompose its risk into systematic and idiosyncratic parts.
\end{enumerate}

\section{Key Takeaways}

\begin{itemize}
    \item The power method finds the dominant eigenpair in
      $O(n^2)$ per iteration; deflation extends it to the full spectrum
      at the cost of accumulated floating-point error.
    \item PCA centres the data, decomposes the covariance, and ranks
      components by explained variance; the reconstruction error
      decreases as more components are retained.
    \item The Fama-French 3-factor model adds size (SMB) and value
      (HML) to the market factor; the $R^2$ improvement over CAPM
      quantifies the incremental explanatory power.
    \item Risk decomposition splits portfolio variance into
      systematic (factor-driven) and idiosyncratic (asset-specific)
      parts; diversification eliminates the idiosyncratic part.
    \item The power method's sign-flip artefact is handled by aligning
      each iterate with its predecessor before testing convergence.
\end{itemize}
\textit{Chapter content pending --- Phase 12 implementation.}

%----------------------------------------------------------------------------------------
% PART V: ADVANCED TOPICS (Expert)
%----------------------------------------------------------------------------------------

\pagelayout{wide}
\addpart{Advanced Topics}
\pagelayout{margin}

\chapter{Market Microstructure}
\label{ch:microstructure}
% \chapter{Market Microstructure}
\label{ch:microstructure}

\begin{companioncode}{quant-microstructure}
Code: \texttt{crates/quant-microstructure/}\\
Dependencies: \crate{quant-core}, \crate{thiserror}\\
Run: \texttt{cargo run -p quant-microstructure --example lob\_demo}\\
Run: \texttt{cargo run -p quant-microstructure --example ofi\_demo}
\end{companioncode}

\section{Learning Objectives}

This chapter moves from the daily/return horizon of portfolio theory
into the intraday world of the limit order book (LOB). After completing
it you can:

\begin{itemize}
    \item Represent a limit order book with price-time priority using
      integer tick prices and a \texttt{BTreeMap} keyed by price level
    \item Implement the three core operations --- add, cancel, market
      order --- and reason about their complexity
    \item Compute the order flow imbalance (OFI) series from a sequence
      of book snapshots
    \item Compute volume-weighted average price (VWAP) from fills and
      trade imbalance from signed trades
    \item Apply the square-root (Almgren-Chriss) and linear market
      impact models to estimate the cost of executing a market order
    \item Combine half-spread and impact into a single execution-cost
      estimate in basis points
\end{itemize}

\begin{keyinsight}
Microstructure is where the frictionless world of Markowitz meets
reality. The bid-ask spread is the price of immediacy: takers pay it,
makers earn it. The square-root impact law says that moving size
$Q$ through the market costs $\sigma\sqrt{Q/V}$ in fraction of price
--- a non-linear function of participation rate. Together, spread
plus impact turn the efficient frontier into a \emph{reachable}
frontier net of trading costs.
\end{keyinsight}

\section{Why Microstructure Matters}

The portfolio theory of \cref{ch:portfolio} assumes frictionless
trading: you can rebalance to the tangency portfolio at no cost. In
reality every trade pays the half-spread and pushes the price against
itself. The combined cost is small for a retail order of 100 shares,
but it dominates the economics of executing a $\$100$M rebalance.

Market microstructure studies the mechanism by which trades occur and
how that mechanism affects price formation. Three quantities capture
most of the relevant friction:

\begin{enumerate}
    \item \textbf{The bid-ask spread} --- the difference between the
      best (lowest) ask and the best (highest) bid. A market taker
      buys at the ask and sells at the bid, paying the full spread;
      the maker on the other side earns half of it on average.
    \item \textbf{Market impact} --- the adverse price move caused by
      your own order. A large buy lifts the ask; subsequent fills
      happen at worse prices. Empirically the impact scales as the
      square root of the order size relative to daily volume.
    \item \textbf{Order flow imbalance (OFI)} --- a signed measure of
      buying versus selling pressure at the top of the book. Positive
      OFI predicts positive short-horizon returns.
\end{enumerate}

\section{The Limit Order Book}

A limit order book aggregates all outstanding resting orders. Each
order has a side (bid or ask), a price, a quantity, and a timestamp.
Orders at the same price execute in \emph{price-time priority}: the
best price first, and at that price, the earliest timestamp first
(FIFO).

We represent prices as integer ticks (\texttt{u64}) to avoid
floating-point comparison issues. A tick size of 1 means one unit of
price (e.g.\ one cent); a tick size of 10 means prices must be
multiples of 10. The book stores bids and asks in two
\texttt{BTreeMap<u64, Vec<Order>>} keyed by price. For bids, the best
(highest) price is the last key; for asks, the best (lowest) price is
the first key. A \texttt{HashMap<order\_id, (side, price)>} index
gives $O(1)$ cancel lookup.

\vspace{-2cm}
\marginnote{%
\textbf{Why integer ticks?}\\[0.3em]
Float comparison is non-associative and rounding-dependent. The book
must compare prices \emph{exactly}: an order either matches the tick
grid or it does not. \texttt{u64} ticks make that comparison trivial
and unambiguous.
}

\marginnote{%
\textbf{BTreeMap per side}\\[0.3em]
A \texttt{BTreeMap<u64, \_>} gives ordered price levels with
$O(\log L)$ insert and best-price lookup. The price \emph{is} the key,
so the book never has to sort --- ordering is structural. A
\texttt{HashMap} would not preserve order; a sorted \texttt{Vec} would
shift data on every insert.
}

\begin{marginfigure}
    \centering
    \vspace{-1cm}
    % Limit order book depth diagram.
% Horizontal bars: each price level is one row, bar length = quantity.
% Asks (sellers) above in red, bids (buyers) below in blue. The spread
% is the gap between best bid (100) and best ask (101).

\begin{tikzpicture}[
    scale=0.85,
    every node/.style={font=\scriptsize},
    bidbar/.style={fill=blue!45, draw=blue!70!black},
    askbar/.style={fill=red!45,  draw=red!70!black},
]

% Axes
\draw[->, thick, >=stealth] (0,0) -- (5.2,0) node[below right=-1pt] {Qty};
\draw[->, thick, >=stealth] (0,0) -- (0,6.6) node[above] {Price};

% Asks (descending from best ask at 101, drawn above the mid line)
%   y-position chosen so each level sits on its own row
\node[left] at (0, 6.10) {103};
\filldraw[askbar] (0, 5.95) rectangle (1.0, 6.25);
\node[right] at (1.05, 6.10) {20};

\node[left] at (0, 5.40) {102};
\filldraw[askbar] (0, 5.25) rectangle (2.5, 5.55);
\node[right] at (2.55, 5.40) {50};

\node[left] at (0, 4.70) {101};
\filldraw[askbar] (0, 4.55) rectangle (1.5, 4.85);
\node[right] at (1.55, 4.70) {30};

% Spread annotation between best ask (101) and best bid (100)
\draw[<->, dashed, gray!70] (3.8, 4.40) -- (3.8, 4.10);
\node[right] at (3.8, 4.25) {spread = 1};

% Bids (descending from best bid at 100, drawn below the mid line)
\node[left] at (0, 4.10) {100};
\filldraw[bidbar] (0, 3.95) rectangle (2.0, 4.25);
\node[right] at (2.05, 4.10) {40};

\node[left] at (0, 3.40) {99};
\filldraw[bidbar] (0, 3.25) rectangle (3.0, 3.55);
\node[right] at (3.05, 3.40) {60};

\node[left] at (0, 2.70) {98};
\filldraw[bidbar] (0, 2.55) rectangle (1.5, 2.85);
\node[right] at (1.55, 2.70) {30};

% Mid price dashed line
\draw[dashed, gray!50] (0, 4.40) -- (4.7, 4.40);
\node[right, gray!60!black] at (4.7, 4.40) {mid = 100.5};

% Side labels
\node[red!70!black,  font=\bfseries] at (4.4, 5.85) {Ask};
\node[blue!70!black, font=\bfseries] at (4.4, 3.40) {Bid};

\end{tikzpicture}
    \caption{Resting depth at three price levels on each side. Best
    bid at 100 (40 shares), best ask at 101 (30 shares), spread = 1,
    mid = 100.5.}
    \label{fig:lob-depth}
\end{marginfigure}

\subsection{Core Types}

The basic types are intentionally small and \texttt{Copy} where
possible so that the book can pass them by value without ceremony.

\begin{lstlisting}[caption={Core microstructure types}]
pub enum Side { Bid, Ask }

pub struct Order {
    pub id: u64,
    pub side: Side,
    pub price: u64,
    pub quantity: u64,
    pub timestamp: u64,
}

pub struct Level {
    pub price: u64,
    pub quantity: u64,
    pub order_count: usize,
}

pub struct Fill {
    pub price: u64,
    pub quantity: u64,
    pub maker_order_id: u64,
}
\end{lstlisting}

\subsection{Adding and Cancelling Orders}

\marginnote{%
\textbf{Price-time priority}\\[0.3em]
Orders at the same price execute in FIFO order: the earliest
timestamp fills first. At each price level we store orders in a
\texttt{Vec} in insertion order, and market orders consume the front
of the vector. Time is secondary to price --- a worse-priced order
never jumps ahead of a better one, regardless of age.
}

\func{add\_order} validates that the price is a multiple of the tick
size, that the quantity is positive, and that the order id is unique.
It then inserts the order at the back of the vector for its price
level and records the id in the index.

\func{cancel\_order} looks up the (side, price) pair in the index,
finds the order in the level's vector, removes it, and prunes the
level if it becomes empty. Both operations are $O(\log L + K)$ where
$L$ is the number of price levels and $K$ is the number of orders at
that level; the index makes cancel $O(1)$ to locate the level.

\begin{lstlisting}[caption={Adding a limit order}]
pub fn add_order(&mut self, order: Order) -> Result<(), MicroError> {
    if self.tick_size > 0 && order.price % self.tick_size != 0 {
        return Err(MicroError::InvalidOrder(format!(
            "price {} is not a multiple of tick size {}",
            order.price, self.tick_size
        )));
    }
    if order.quantity == 0 {
        return Err(MicroError::InvalidOrder("quantity must be positive".into()));
    }
    if self.index.contains_key(&order.id) {
        return Err(MicroError::InvalidOrder(format!(
            "duplicate order id {}", order.id
        )));
    }
    self.index.insert(order.id, (order.side, order.price));
    let book = match order.side {
        Side::Bid => &mut self.bids,
        Side::Ask => &mut self.asks,
    };
    book.entry(order.price).or_insert_with(Vec::new).push(order);
    Ok(())
}
\end{lstlisting}

\subsection{Market Orders and Price-Time Priority}

A market order has no price: it takes liquidity from the opposite side
until filled or until that side is empty. A market buy walks the asks
in ascending price order; a market sell walks the bids in descending
price order. At each price level, fills happen in FIFO order --- the
front of the vector is consumed first, exactly matching price-time
priority.

\begin{lstlisting}[caption={Market-order execution walks the opposite side}]
pub fn market_order(&mut self, side: Side, quantity: u64) -> Vec<Fill> {
    let mut fills = Vec::new();
    let mut remaining = quantity;
    let prices: Vec<u64> = match side {
        Side::Bid => self.asks.keys().copied().collect(),         // buy -> asks ascending
        Side::Ask => self.bids.keys().rev().copied().collect(),   // sell -> bids descending
    };
    for price in prices {
        if remaining == 0 { break; }
        let book = match side {
            Side::Bid => &mut self.asks,
            Side::Ask => &mut self.bids,
        };
        if let Some(orders) = book.get_mut(&price) {
            while remaining > 0 && !orders.is_empty() {
                let front = &mut orders[0];
                let fill_qty = remaining.min(front.quantity);
                let maker_id = front.id;
                fills.push(Fill { price, quantity: fill_qty, maker_order_id: maker_id });
                remaining -= fill_qty;
                front.quantity -= fill_qty;
                if front.quantity == 0 {
                    let removed = orders.remove(0);
                    self.index.remove(&removed.id);
                }
            }
            if orders.is_empty() { book.remove(&price); }
        }
    }
    fills
}
\end{lstlisting}

\subsection{Top-of-Book Accessors}

The book exposes \func{best\_bid}, \func{best\_ask}, \func{mid\_price}
and \func{spread}. For a \texttt{BTreeMap}, the best bid is
\texttt{iter().next\_back()} (highest key) and the best ask is
\texttt{iter().next()} (lowest key). Depth at the top $L$ levels is
the first $L$ keys from each side, aggregated.

\section{Order Flow Imbalance}

\marginnote{%
\textbf{OFI sign convention}\\[0.3em]
Positive OFI = buy-side pressure (bids growing relative to asks).
Negative OFI = sell-side pressure. The signal predicts short-horizon
returns: a rising bid queue means buyers are more eager than sellers,
and the mid price tends to drift up over the next few seconds to
minutes.
}

Cont, Kukanov and Stoikov (2014) introduced OFI as a signed measure of
buying versus selling pressure at the top of the book. Given a sequence
of snapshots, the OFI at transition $t$ is
\[
    \mathrm{OFI}_t = \Delta q^{bid}_t - \Delta q^{ask}_t,
\]
where $\Delta q^{bid}_t$ is the change in available quantity at the
best bid between $t-1$ and $t$, and similarly for the ask. Positive
OFI indicates buy-side pressure (bids growing relative to asks);
negative OFI indicates sell-side pressure.

\begin{lstlisting}[caption={Order flow imbalance from book snapshots}]
pub fn order_flow_imbalance(snapshots: &[BookSnapshot]) -> Vec<f64> {
    if snapshots.len() < 2 { return Vec::new(); }
    let mut result = Vec::with_capacity(snapshots.len() - 1);
    for i in 1..snapshots.len() {
        let prev_bid = snapshots[i - 1].best_bid.as_ref().map(|l| l.quantity as f64).unwrap_or(0.0);
        let prev_ask = snapshots[i - 1].best_ask.as_ref().map(|l| l.quantity as f64).unwrap_or(0.0);
        let curr_bid = snapshots[i].best_bid.as_ref().map(|l| l.quantity as f64).unwrap_or(0.0);
        let curr_ask = snapshots[i].best_ask.as_ref().map(|l| l.quantity as f64).unwrap_or(0.0);
        result.push((curr_bid - prev_bid) - (curr_ask - prev_ask));
    }
    result
}
\end{lstlisting}

\subsection{VWAP and Trade Imbalance}

Two complementary measures aggregate the trade tape. VWAP is the
volume-weighted average fill price over a window:
\[
    \mathrm{VWAP} = \frac{\sum_i p_i q_i}{\sum_i q_i}.
\]
Trade imbalance normalises the signed volume into $[-1, 1]$:
\[
    \mathrm{TI} = \frac{V^{buy} - V^{sell}}{V^{buy} + V^{sell}}.
\]
A value near $+1$ means all trades were buy-initiated; near $-1$ all
sell-initiated; near $0$ balanced flow.

\section{Market Impact Models}

\marginnote{%
\textbf{Square-root law}\\[0.3em]
The most robust empirical regularity in market impact: across equity
markets, futures, and FX, impact scales as $\sigma\sqrt{Q/V}$. The
exponent is not $1$ (linear) and not $2/3$ (nonlinear propagation) ---
it is consistently close to $1/2$. Doubling order size multiplies
impact by $\sqrt{2}\approx 1.414$, \emph{not} by $2$.
}

The square-root impact law (Almgren and Chriss, 2000) states that the
impact of an order is proportional to the volatility times the square
root of the participation rate $Q / V$:
\[
    \text{Impact} = \sigma \sqrt{\frac{Q}{V}},
\]
where $\sigma$ is the daily volatility, $Q$ the order size, and $V$
the average daily volume. The result is a fraction of price (e.g.\
$0.001 = 10$ bps). This is an empirical regularity observed across
equity markets with remarkable consistency.

A simpler linear model replaces the square root with a linear term:
\[
    \text{Impact}_{\text{lin}} = \eta \cdot \frac{Q}{V},
\]
where $\eta$ is an impact coefficient calibrated to data.

\begin{lstlisting}[caption={Square-root and linear impact models}]
pub fn sqrt_impact(volatility: f64, order_size: f64, daily_volume: f64) -> f64 {
    if daily_volume <= 0.0 { return 0.0; }
    let participation = order_size / daily_volume;
    volatility * participation.abs().sqrt()
}

pub fn linear_impact(eta: f64, order_size: f64, daily_volume: f64) -> f64 {
    if daily_volume <= 0.0 { return 0.0; }
    eta * (order_size / daily_volume)
}
\end{lstlisting}

\subsection{Execution Cost}

\marginnote{%
\textbf{Half-spread = cost of immediacy}\\[0.3em]
A market taker buys at the ask and sells at the bid, paying the full
spread. The maker on the other side earns half of it on average. So
the taker's spread cost is $s/2$, not $s$: the other half is the
maker's compensation for providing liquidity.
}

The total cost of executing a market order combines the half-spread
(paid by the aggressor) with the square-root impact:
\[
    c_{\text{total}} = \underbrace{\frac{s}{2}}_{\text{half-spread}}
                     + \underbrace{\sigma\sqrt{\frac{Q}{V}}}_{\text{impact}}.
\]
Both terms are expressed as fractions of the mid price. Multiplying by
$10{,}000$ converts to basis points.

\begin{lstlisting}[caption={Execution cost combines spread and impact}]
pub fn execution_cost(
    order_size: f64,
    daily_volume: f64,
    volatility: f64,
    spread: f64,
) -> ExecutionCost {
    let spread_cost = spread / 2.0;
    let impact_cost = sqrt_impact(volatility, order_size, daily_volume);
    ExecutionCost {
        spread_cost,
        impact_cost,
        total_cost: spread_cost + impact_cost,
    }
}
\end{lstlisting}

\section{Numerical Demonstration}

We exercise the crate on a synthetic book and a 5-snapshot OFI series
(\texttt{cargo run -p quant-microstructure --example lob\_demo} and
\texttt{--example ofi\_demo}).

\paragraph{Limit order book demo.}
The book is seeded with three ask levels ($30, 50, 20$ shares at
$101, 102, 103$) and two bid levels ($40, 60$ at $99, 98$). A market
buy of $75$ shares sweeps the best ask fully ($30 @ 101$) and bites
$45$ out of the second level ($45 @ 102$). The post-trade best ask is
$102$ with $5$ shares remaining; the average fill price is
$101.60$.

\paragraph{OFI demo.}
The five snapshots show bid quantity growing from $50$ to $75$ and
then shrinking to $55$, while ask quantity shrinks from $40$ to $25$
and then grows to $45$. The OFI series is $[+15, +25, -15, -25]$, with
zero cumulative OFI at the end --- symmetric buy-then-sell pressure.

\paragraph{Execution cost.}
For a $2{,}000$-share order against a $\$1\text{M}$ daily volume,
$\sigma = 2\%$, spread $= 2$ bps:
\begin{center}
\begin{tabular}{lr}
\toprule
Component & Cost (bps) \\
\midrule
Half-spread & 1.00 \\
Square-root impact ($0.02\sqrt{0.002}$) & 8.94 \\
Total & 9.94 \\
\bottomrule
\end{tabular}
\end{center}

The impact dominates the spread by a factor of nearly $9\times$ at
this participation rate. For a $200$-share order the impact falls to
$2.83$ bps (a factor of $\sqrt{10}$ lower), and the spread becomes the
larger cost --- the square-root law is what makes small orders cheap
and large orders expensive.

\section{Complexity and Design Notes}

\paragraph{Why BTreeMap?} A \texttt{BTreeMap<u64, Vec<Order>>} gives
us ordered price levels with $O(\log L)$ insert, delete and best-price
lookup. A \texttt{HashMap} would not preserve order, and a sorted
\texttt{Vec} would shift data on every insert. The price is the key,
so the book never has to sort: ordering is structural.

\paragraph{Why integer prices?} Floating-point comparison is
non-associative and rounding-dependent. Limit order books must compare
prices exactly: an order at $100.10$ either matches the tick grid or
it does not. Representing prices as integer ticks makes that comparison
trivial and unambiguous.

\paragraph{Why hand-rolled?} Production trading systems use heavily
optimised LOB implementations (e.g.\ \texttt{lmax-disruptor} patterns,
lock-free queues). This crate is a teaching artefact: every line is
short enough to read in one sitting, and the price-time priority rule
is visible in the code rather than hidden behind a library API.

\section{Exercises}

\begin{enumerate}
    \item \textbf{Order book builder.} Write a function that takes a
      vector of \texttt{(Side, price, qty, ts)} tuples and returns a
      populated \texttt{OrderBook}. Verify the invariants
      \func{best\_bid}, \func{best\_ask}, \func{total\_volume} on the
      result.
    \item \textbf{Liquidity sweep.} Construct a book with $5$ ask
      levels totalling $1{,}000$ shares. Execute a market buy of
      $750$ shares and verify that the average fill price is less than
      or equal to the volume-weighted ask price of the swept levels.
    \item \textbf{OFI and returns.} Generate $100$ snapshots from a
      stochastic process where bid and ask quantities are driven by a
      latent ``buy pressure'' signal. Compute the OFI series and the
      mid-price return series; verify they are positively correlated.
    \item \textbf{Impact scaling.} For a fixed $\sigma$ and $V$,
      compute \func{sqrt\_impact} for $Q \in \{100, 200, 500, 1000,
      2000, 5000\}$ and plot impact versus $Q$ on a log-log scale.
      Confirm the slope is $1/2$.
    \item \textbf{Net frontier.} Take the Markowitz tangency portfolio
      from \cref{ch:portfolio} and an assumed turnover of $50\%$ per
      rebalance. Estimate the per-rebalance cost in bps using
      \func{execution\_cost} with $Q = 0.5 \times V_{\text{daily}}$
      and subtract it from the expected return. Plot the net efficient
      frontier.
\end{enumerate}

\section{Key Takeaways}

\begin{itemize}
    \item The limit order book is a price-time priority queue with
      integer tick prices; a \texttt{BTreeMap} per side plus an
      id-index gives a compact, correct implementation.
    \item Market orders walk the opposite side in price order and fill
      level-by-level in FIFO order; partial fills are the rule, not
      the exception.
    \item OFI = $\Delta q^{bid} - \Delta q^{ask}$ is a signed top-of
      -book pressure signal; trade imbalance normalises signed tape
      volume into $[-1, 1]$.
    \item The square-root law $\sigma\sqrt{Q/V}$ is the most robust
      empirical regularity in market impact: doubling size multiplies
      impact by $\sqrt{2}$, not by $2$.
    \item Total execution cost = half-spread $+$ impact. For small
      orders the spread dominates; for large orders the impact
      dominates.
\end{itemize}
\textit{Chapter content pending --- Phase 13 implementation.}

\chapter{AFML: From Research to Production}
\label{ch:afml}
% \chapter{AFML: Backtesting with Triple Barriers and Purged CV}
\label{ch:afml}

\begin{companioncode}{quant-backtest}
Code: \texttt{crates/quant-backtest/}\\
Dependencies: \crate{quant-core}, \crate{quant-timeseries}, \crate{quant-portfolio}, \crate{thiserror}\\
Run: \texttt{cargo run -p quant-backtest --example triple\_barrier}\\
Run: \texttt{cargo run -p quant-backtest --example purged\_cv}\\
Run: \texttt{cargo run -p quant-backtest --example kelly\_sizing}
\end{companioncode}

\section{Learning Objectives}

This chapter leaves behind the frictionless world of portfolio theory
and the single-trade world of microstructure, and enters the engine
room of systematic strategy research: how do you train a model on
time-series labels, cross-validate it without leaking information, and
size the resulting bets? After completing it you can:

\begin{itemize}
    \item Explain why standard k-fold cross-validation leaks
      information on financial time series with overlapping events
    \item Label a price series with the triple-barrier method
      (profit-taking, stop-loss, time barrier) and convert the labels
      to a binary outcome
    \item Compute the concurrency of overlapping events and derive
      sample weights from event uniqueness
    \item Generate purged k-fold splits with an embargo period that
      removes training samples overlapping --- and immediately
      following --- each test fold
    \item Compute the full Kelly criterion from a win probability and
      a win/loss ratio, and apply a fractional multiplier for risk
      control
    \item Run the end-to-end AFML pipeline that combines all of the
      above into a single backtest report
\end{itemize}

\begin{keyinsight}
Traditional backtesting fails on three fronts: labels are arbitrary
(fixed-horizon return sign), samples overlap (events that span
multiple bars share information), and cross-validation leaks (the
test fold is not independent of the training fold). The AFML framework
fixes all three: triple-barrier labels are path-dependent, sample
weights discount overlap, and purged k-fold with embargo removes
leakage. The result is an out-of-sample Sharpe that you can actually
trust.
\end{keyinsight}

\section{Why Traditional Backtesting Fails}

A naive backtest labels each bar by the sign of the return over a
fixed horizon $h$ --- for example, the 5-bar forward return --- and
trains a classifier on those labels. Three problems compound:

\begin{enumerate}
    \item \textbf{Arbitrary labels.} A 5-bar forward return says
      nothing about the path: a position that gained 4\%, lost 3\%,
      and gained 4\% again is labelled the same as one that monotonically
      drifted up. Yet the lived experience of holding the position
      is very different --- the first path probably hit a stop loss.
    \item \textbf{Overlapping samples.} Consecutive 5-bar labels share
      four of their five bars. A classifier that treats them as
      independent over-counts the information in the sample by a factor
      of $\sim h$, inflating the apparent edge.
    \item \textbf{Leakage in CV.} Standard k-fold CV assigns bars to
      folds uniformly at random. A training bar that is part of a
      5-bar event ending inside the test fold has already seen the
      label of that test sample --- the validation is no longer
      out-of-sample.
\end{enumerate}

The AFML framework of López de Prado (2018) addresses each in turn.

\section{The Triple-Barrier Method}

For each entry point $t_0$ the method places three barriers around
the entry price $P_{t_0}$:

\begin{itemize}
    \item \textbf{Upper} at $P_{t_0}(1 + \theta^+)$, the
      profit-taking threshold.
    \item \textbf{Lower} at $P_{t_0}(1 + \theta^-)$, the
      stop-loss threshold ($\theta^- < 0$).
    \item \textbf{Time} at $t_0 + T$, the maximum holding period.
\end{itemize}

We walk the price path forward from $t_0$ and label the event by
the \emph{first} barrier touched. The holding period is the number of
bars from entry to exit; the return is $(P_{\text{exit}} -
P_{t_0})/P_{t_0}$.

\begin{marginfigure}
    \centering
    \begin{tikzpicture}[scale=0.85, font=\scriptsize]
        % Price axis
        \draw[->] (0,0) -- (7,0) node[right] {time};
        \draw[->] (0,0) -- (0,4.5) node[above] {price};
        % Entry
        \draw[dashed] (0,2) -- (7,2);
        \node[left] at (0,2) {$P_{t_0}$};
        % Upper barrier
        \draw[blue, thick] (0,3) -- (7,3);
        \node[blue, right] at (7,3) {Upper $\theta^+$};
        % Lower barrier
        \draw[red, thick] (0,1) -- (7,1);
        \node[red, right] at (7,1) {Lower $\theta^-$};
        % Time barrier
        \draw[green!50!black, thick] (5,0) -- (5,4);
        \node[green!50!black, below] at (5,0) {$T$};
        % Price path
        \draw[gray] (0,2) .. controls (1,2.4) and (2,1.4) .. (3,1.7);
        \draw[gray] (3,1.7) .. controls (4,2.3) and (4.5,3.1) .. (5,3);
        \fill[blue] (5,3) circle (1.5pt);
        \node[blue, above] at (5,3.1) {Upper hit};
    \end{tikzpicture}
    \caption{Triple barriers around the entry price. The first
    barrier touched --- here the upper barrier at time $T$ ---
    determines the label.}
    \label{fig:triple-barrier}
\end{marginfigure}

\subsection{Configuration and Labels}

\marginnote{%
\textbf{Why three barriers?}\\[0.3em]
A path-dependent label captures the lived experience of the trade.
Two trades with the same 5-bar return but different paths get
different labels: one hits the upper barrier (a win), the other
hits the stop loss (a loss). The label is no longer a function of
the endpoint, but of the whole trajectory.
}

\begin{lstlisting}[caption={Triple-barrier configuration and label type}]
pub struct TripleBarrierConfig {
    pub upper_barrier: f64, // e.g. 0.02 = +2%
    pub lower_barrier: f64, // e.g. -0.02 = -2%
    pub time_barrier: usize, // max bars to hold
    pub min_return: f64, // threshold for Time label
}

pub enum TripleBarrierLabel { Upper, Lower, Time }

pub struct LabeledEvent {
    pub entry_index: usize,
    pub exit_index: usize,
    pub label: TripleBarrierLabel,
    pub return_pct: f64,
    pub holding_period: usize,
}
\end{lstlisting}

\subsection{Labeling a Single Event}

\func{label\_one} walks prices forward from the entry index, checking
at each bar whether the upper or lower barrier has been hit. If neither
is touched before the time barrier expires, the event is labelled
\texttt{Time} with the return at the last bar.

\marginnote{%
\textbf{No look-ahead in features}\\[0.3em]
The triple-barrier method \emph{intentionally} uses future prices to
label the event. That is not look-ahead bias in the model: the label
is the target, not a feature. Features must be computed strictly
from information available at $t_0$. The discipline is: train on
labels that depend on the future, predict them from features that
do not.
}

\begin{lstlisting}[caption={Walking the price path to find the first barrier hit}]
fn label_one(prices: &[f64], entry_index: usize,
             config: &TripleBarrierConfig) -> Result<LabeledEvent, BacktestError> {
    let p_entry = prices[entry_index];
    let upper_px = p_entry * (1.0 + config.upper_barrier);
    let lower_px = p_entry * (1.0 + config.lower_barrier);
    let last = (entry_index + config.time_barrier).min(prices.len() - 1);
    for (t, &p) in prices.iter().enumerate()
        .take(last + 1).skip(entry_index + 1) {
        if p >= upper_px {
            return Ok(LabeledEvent {
                entry_index, exit_index: t,
                label: TripleBarrierLabel::Upper,
                return_pct: (p - p_entry) / p_entry,
                holding_period: t - entry_index,
            });
        }
        if p <= lower_px {
            return Ok(LabeledEvent {
                entry_index, exit_index: t,
                label: TripleBarrierLabel::Lower,
                return_pct: (p - p_entry) / p_entry,
                holding_period: t - entry_index,
            });
        }
    }
    let p_exit = prices[last];
    Ok(LabeledEvent {
        entry_index, exit_index: last,
        label: TripleBarrierLabel::Time,
        return_pct: (p_exit - p_entry) / p_entry,
        holding_period: last - entry_index,
    })
}
\end{lstlisting}

\subsection{Binary Labels}

For a long-only strategy the label is naturally binary: Upper is a
win (1), Lower is a loss (0), and Time is a win iff the return at
the time barrier exceeds $\theta^+_{\min}$.

\begin{lstlisting}[caption={Converting a triple-barrier label to a binary outcome}]
pub fn to_binary_label_helper(label: TripleBarrierLabel,
                              event: &LabeledEvent,
                              min_return: f64) -> i32 {
    match label {
        TripleBarrierLabel::Upper => 1,
        TripleBarrierLabel::Lower => 0,
        TripleBarrierLabel::Time => {
            if event.return_pct >= min_return { 1 } else { 0 }
        }
    }
}
\end{lstlisting}

\section{Sample Uniqueness and Weighting}

When events overlap in time, each carries less unique information. A
training sample whose event $[t_0, t_1]$ shares all of its bars with
three other events is worth roughly a third of a sample that no other
event touches. We capture this with the concurrency count
$c_t$ = number of events active at bar $t$:

\begin{equation}
    c_t = \sum_i \mathbb{1}_{[t_0^{(i)},\, t_1^{(i)}]}(t).
\end{equation}

The average uniqueness of event $i$ is the mean of $1/c_t$ over its
duration, and the sample weight is the inverse of the average
concurrency:

\begin{equation}
    u_i = \frac{1}{t_1^{(i)} - t_0^{(i)}} \sum_{t=t_0^{(i)}}^{t_1^{(i)}} \frac{1}{c_t},
    \qquad
    w_i = \frac{1}{\bar{c}_i}.
\end{equation}

\marginnote{%
\textbf{Why inverse concurrency?}\\[0.3em]
If three events overlap on a bar, each one contributes a third of
the unique information there. Summing $1/c_t$ over the duration
averages that discount across the event's life. The resulting
weight is high for isolated events and low for crowded ones, so
the classifier down-weights redundant samples automatically.
}

\begin{lstlisting}[caption={Concurrency counts and average uniqueness}]
pub fn concurrent_events(events: &[LabeledEvent], n_bars: usize) -> Vec<usize> {
    let mut counts = vec![0usize; n_bars];
    for ev in events {
        let start = ev.entry_index.min(n_bars);
        let end = ev.exit_index.min(n_bars.saturating_sub(1));
        if start >= n_bars { continue; }
        for c in counts.iter_mut().take(end + 1).skip(start) { *c += 1; }
    }
    counts
}

pub fn average_uniqueness(events: &[LabeledEvent], n_bars: usize) -> Vec<f64> {
    let concurrent = concurrent_events(events, n_bars);
    events.iter().map(|ev| {
        let start = ev.entry_index;
        let end = ev.exit_index.min(n_bars.saturating_sub(1));
        if end < start { return 0.0; }
        let mut sum_inv = 0.0;
        let mut count = 0;
        for &c in concurrent.iter().take(end + 1).skip(start) {
            if c > 0 { sum_inv += 1.0 / c as f64; count += 1; }
        }
        if count == 0 { 0.0 } else { sum_inv / count as f64 }
    }).collect()
}
\end{lstlisting}

The final weights are the inverse of the average concurrency during
each event, normalised to sum to one across the sample:

\begin{lstlisting}[caption={Sample weights from inverse concurrency}]
pub fn sample_weights(events: &[LabeledEvent], n_bars: usize) -> Vec<f64> {
    let concurrent = concurrent_events(events, n_bars);
    let mut weights: Vec<f64> = events.iter().map(|ev| {
        let start = ev.entry_index;
        let end = ev.exit_index.min(n_bars.saturating_sub(1));
        if end < start { return 0.0; }
        let mut sum_c = 0.0;
        let mut count = 0;
        for &c in concurrent.iter().take(end + 1).skip(start) {
            if c > 0 { sum_c += c as f64; count += 1; }
        }
        if count == 0 { 0.0 } else { 1.0 / (sum_c / count as f64) }
    }).collect();
    let total: f64 = weights.iter().sum();
    if total > 0.0 { for w in &mut weights { *w /= total; } }
    weights
}
\end{lstlisting}

\section{Purged K-Fold Cross-Validation}

Standard k-fold CV leaks information whenever a training event
overlaps a test event. \emph{Purging} removes from the training
fold any event whose $[t_0, t_1]$ interval intersects the test fold
$[t_{\text{start}}, t_{\text{end}})$.

\begin{lstlisting}[caption={Overlap test for purging}]
pub fn event_overlaps(event: &LabeledEvent, start: usize, end: usize) -> bool {
    event.entry_index < end && event.exit_index > start
}
\end{lstlisting}

\marginnote{%
\textbf{Train must not see test}\\[0.3em]
If a training event ends inside the test window, its label ---
which depends on prices inside the test window --- has already
told the model what happens there. Purging removes exactly those
events, leaving the training fold honest.
}

\subsection{Embargo}

Purging is not enough under autocorrelation: a training event
\emph{starting} just after the test window ends still carries
test-window information through serially-correlated features. The
\emph{embargo} removes training events whose entry falls in
$[t_{\text{end}}, t_{\text{end}} + e)$.

\begin{lstlisting}[caption={Purged k-fold with embargo}]
pub fn purged_kfold_splits(events: &[LabeledEvent], n_bars: usize,
                           config: &PurgedKFoldConfig) -> Vec<PurgedSplit> {
    let fold_size = n_bars / config.n_folds;
    let mut splits = Vec::with_capacity(config.n_folds);
    for k in 0..config.n_folds {
        let t_start = k * fold_size;
        let t_end = if k == config.n_folds - 1 { n_bars }
                    else { (k + 1) * fold_size };
        let embargo_end = (t_end + config.embargo).min(n_bars);
        let mut train = Vec::new();
        let mut test = Vec::new();
        let mut purged = 0usize;
        let mut embargoed = 0usize;
        for (i, ev) in events.iter().enumerate() {
            if ev.entry_index >= t_start && ev.entry_index < t_end {
                test.push(i); continue;
            }
            if event_overlaps(ev, t_start, t_end) { purged += 1; continue; }
            if ev.entry_index >= t_end && ev.entry_index < embargo_end {
                embargoed += 1; continue;
            }
            train.push(i);
        }
        splits.push(PurgedSplit {
            train_indices: train, test_indices: test,
            purged_count: purged, embargoed_count: embargoed,
        });
    }
    splits
}
\end{lstlisting}

\marginnote{%
\textbf{Embargo calibrates autocorrelation}\\[0.3em]
Set the embargo to the lag at which the autocorrelation function
of the features drops below significance. For daily returns on
liquid equities a 1--5 bar embargo is typical; for noisy
high-frequency features a longer embargo is needed. An embargo of
zero reduces to plain purging.
}

\section{Bet Sizing: The Kelly Criterion}

A good label tells you \emph{when} to bet; it does not tell you
\emph{how much}. The Kelly criterion maximises the expected log
growth of capital. For a binary bet with win probability $p$, loss
probability $q = 1 - p$, and win/loss ratio $b = \bar{r}^+ /
|\bar{r}^-|$ (mean win divided by mean loss), the optimal fraction
of capital to risk on each bet is

\begin{equation}
    f^* = \frac{p\, b - q}{b} = p - \frac{q}{b}.
    \label{eq:kelly}
\end{equation}

The criterion is exact for a sequence of i.i.d.\ bets with
known parameters. In practice the parameters are estimated, so
\emph{fractional Kelly} --- using a fraction $\lambda \in (0, 1]$ of
$f^*$ --- is standard. Half Kelly ($\lambda = 0.5$) achieves most of
the growth with roughly half the variance.

\marginnote{%
\textbf{Kelly in one line}\\[0.3em]
$f^*$ is the edge divided by the odds. With no edge ($p = q =
0.5$, $b = 1$) it is zero --- you bet nothing. With a 60\% win
rate and 1:1 payoff it is $0.6 - 0.4 = 0.2$: risk 20\% of capital
per bet. Beyond that, the variance of log-wealth grows faster
than its mean.
}

\begin{lstlisting}[caption={Kelly criterion and fractional variant}]
pub fn kelly_fraction(win_prob: f64, win_loss_ratio: f64) -> f64 {
    let q = 1.0 - win_prob;
    win_prob - q / win_loss_ratio
}

pub fn fractional_kelly(win_prob: f64, win_loss_ratio: f64, fraction: f64) -> f64 {
    kelly_fraction(win_prob, win_loss_ratio) * fraction
}
\end{lstlisting}

To estimate Kelly from historical trade returns, compute the win
probability as the fraction of positive trades and the win/loss ratio
as the mean win divided by the mean loss:

\begin{lstlisting}[caption={Kelly from a series of trade returns}]
pub fn kelly_from_returns(returns: &[f64]) -> f64 {
    let n = returns.len() as f64;
    let wins: Vec<f64> = returns.iter().copied().filter(|&r| r > 0.0).collect();
    let losses: Vec<f64> = returns.iter().copied().filter(|&r| r < 0.0).collect();
    let p = wins.len() as f64 / n;
    let mean_win = wins.iter().sum::<f64>() / wins.len().max(1) as f64;
    let mean_loss = -losses.iter().sum::<f64>() / losses.len().max(1) as f64;
    if mean_loss <= 0.0 { return 0.0; }
    kelly_fraction(p, mean_win / mean_loss)
}
\end{lstlisting}

\subsection{Position Sizing}

The \func{compute\_position\_size} helper packages full and half Kelly
along with the empirical win probability and win/loss ratio so the
backtest can switch sizing rules without recomputing statistics.

\begin{lstlisting}[caption={Position sizing summary}]
pub struct PositionSize {
    pub kelly_full: f64,
    pub kelly_half: f64,
    pub win_probability: f64,
    pub win_loss_ratio: f64,
}

pub fn compute_position_size(trade_returns: &[f64]) -> PositionSize {
    let n = trade_returns.len() as f64;
    let wins: Vec<f64> = trade_returns.iter().copied().filter(|&r| r > 0.0).collect();
    let losses: Vec<f64> = trade_returns.iter().copied().filter(|&r| r < 0.0).collect();
    let p = wins.len() as f64 / n.max(1) as f64;
    let mean_win = wins.iter().sum::<f64>() / wins.len().max(1) as f64;
    let mean_loss = -losses.iter().sum::<f64>() / losses.len().max(1) as f64;
    let b = if mean_loss > 0.0 { mean_win / mean_loss } else { 0.0 };
    let kelly = kelly_fraction(p, b);
    PositionSize {
        kelly_full: kelly,
        kelly_half: 0.5 * kelly,
        win_probability: p,
        win_loss_ratio: b,
    }
}
\end{lstlisting}

\section{The AFML Pipeline}

The full pipeline ties the pieces together. Given a price series
and an entry step, it:

\begin{enumerate}
    \item Generates entry indices at every \texttt{entry\_step} bars.
    \item Labels each entry with the triple-barrier method.
    \item Computes sample weights from event uniqueness.
    \item Builds purged k-fold splits with embargo.
    \item Computes per-trade returns, the in-sample Sharpe, and the
      mean out-of-sample Sharpe across folds.
    \item Sizes each trade by the chosen \func{BetSizing} rule ---
      equal, full Kelly, half Kelly, or fixed fraction.
    \item Walks the trade returns through a compounded equity curve
      to compute total return and maximum drawdown.
\end{enumerate}

\begin{lstlisting}[caption={End-to-end AFML backtest}]
pub enum BetSizing { Equal, KellyFull, KellyHalf, Fixed(f64) }

pub struct AfmlBacktestConfig {
    pub barrier_config: TripleBarrierConfig,
    pub cv_config: PurgedKFoldConfig,
    pub use_weights: bool,
    pub bet_sizing: BetSizing,
    pub entry_step: usize,
}

pub struct AfmlBacktestResult {
    pub events: Vec<LabeledEvent>,
    pub weights: Vec<f64>,
    pub cv_splits: Vec<PurgedSplit>,
    pub in_sample_sharpe: f64,
    pub out_of_sample_sharpe: f64,
    pub position_size: PositionSize,
    pub total_return: f64,
    pub max_drawdown: f64,
}

pub fn afml_backtest(prices: &[f64],
                     config: &AfmlBacktestConfig)
    -> Result<AfmlBacktestResult, BacktestError>
{
    config.barrier_config.validate()?;
    // ... generate entries, label, weight, split, size, evaluate ...
    Ok(AfmlBacktestResult { ... })
}
\end{lstlisting}

\marginnote{%
\textbf{In-sample vs out-of-sample}\\[0.3em]
The in-sample Sharpe is computed on all trade returns. The
out-of-sample Sharpe is the mean of the per-fold Sharpe ratios
on each test set. A large gap between the two is the symptom of
overfitting that purged CV is designed to expose.
}

\section{Meta-Labeling (Brief)}

A primary model generates a direction signal (long/short/flat); a
secondary model --- the meta-model --- predicts whether the primary
signal will be correct. The triple-barrier label of the primary
model becomes the target of the meta-model. This separates
\emph{direction} from \emph{size}: the primary model bets often, the
meta-model decides how much to trust each bet. The result is a
betting strategy that is more selective than the primary alone and,
empirically, has a higher Sharpe.

\section{Exercises}

\begin{enumerate}
    \item \textbf{Custom barriers.} Generalise the triple-barrier
      method so that the upper and lower barriers are functions of
      rolling volatility rather than fixed percentages. How does the
      label distribution change on a noisy versus a calm series?
    \item \textbf{Embargo calibration.} Compute the autocorrelation
      function of fractional-differentiated returns (see
      \cref{ch:timeseries}) and set the embargo to the first lag
      where the ACF drops below significance. Compare the
      out-of-sample Sharpe with and without the embargo.
    \item \textbf{Volatility-scaled Kelly.} Replace the constant
      Kelly fraction with one that scales inversely with rolling
      volatility. Verify that the drawdown distribution tightens.
    \item \textbf{Meta-labeling.} Use the sign of a moving-average
      crossover as the primary signal, label it with the triple
      barrier, and train a logistic regression on the labelled events
      to predict correctness. Compare the meta-filtered strategy to
      the primary alone.
\end{enumerate}

\section{Key Takeaways}

\begin{itemize}
    \item \textbf{Triple-barrier labeling} converts price paths into
      supervised learning targets by setting profit-taking, stop-loss,
      and time barriers. It avoids the look-ahead bias of fixed-horizon
      labeling.

    \item \textbf{Sample weights from uniqueness} down-weight events
      that overlap heavily with others, reducing overfitting to
      clustered periods. The average uniqueness metric quantifies
      temporal independence.

    \item \textbf{Purged k-fold cross-validation} prevents information
      leakage by removing overlapping events from training folds and
      applying an embargo period. It exposes overfitting that standard
      k-fold misses in time-series data.

    \item \textbf{Walk-forward analysis} validates strategy robustness
      by measuring performance degradation from in-sample to
      out-of-sample. Walk-Forward Efficiency (WFE) below 0.5 indicates
      severe overfitting (see \cref{sec:walk-forward} in
      \cref{ch:backtest}).

    \item \textbf{Kelly criterion bet sizing} maximizes long-run
      geometric growth by sizing positions proportional to edge and
      inversely to volatility. Fractional Kelly trades growth for
      reduced drawdowns.

    \item \textbf{Meta-labeling} separates direction (primary model)
      from confidence (meta-model), producing more selective strategies
      with higher Sharpe ratios. The triple-barrier label of the
      primary becomes the target of the meta-model.

    \item \textbf{Trait-based architecture} enables composition of
      labeling, cross-validation, and bet sizing strategies without
      code duplication. The generic backtest engine combines any
      \texttt{Labeler}, \texttt{CrossValidator}, and \texttt{BetSizer}
      via trait bounds.

    \item \textbf{AFML workflow integration} orchestrates the full
      pipeline: label events $\to$ weight samples $\to$ cross-validate
      $\to$ size bets $\to$ evaluate. Each component is swappable via
      trait implementations, promoting experimentation and reuse.
\end{itemize}

\section{References}

\begin{itemize}
    \item López de Prado, M.\ (2018). \emph{Advances in Financial
      Machine Learning}. Wiley. Chapters 3 (triple-barrier), 4
      (sample weights), 7 (cross-validation), 10 (bet sizing).
    \item Kelly, J. L.\ (1956). ``A New Interpretation of Information
      Rate.'' \emph{Bell System Technical Journal}.
    \item \href{https://github.com/bitscrafts/quant-lab-rust}{quant-lab-rust}
      companion code, crate \crate{quant-backtest}.
\end{itemize}
\textit{Chapter content pending --- Phase 14 implementation.}

%----------------------------------------------------------------------------------------
%	APPENDICES
%----------------------------------------------------------------------------------------

\appendix

\pagelayout{wide}
\addpart{Appendices}
\pagelayout{margin}

\chapter{Mathematical Notation}
\label{app:notation}
\textit{Appendix content pending.}

\chapter{Rust Patterns for Finance}
\label{app:rust}
\textit{Appendix content pending.}

\chapter{Dataset Sources}
\label{app:datasets}
\textit{Appendix content pending.}

%----------------------------------------------------------------------------------------
%	BACK MATTER
%----------------------------------------------------------------------------------------

\backmatter
\setchapterstyle{plain}

%----------------------------------------------------------------------------------------
%	BIBLIOGRAPHY
%----------------------------------------------------------------------------------------

\defbibnote{bibnote}{References used in this book.\par\bigskip}
\printbibliography[heading=bibintoc, title=Bibliography, prenote=bibnote]

%----------------------------------------------------------------------------------------
%	INDEX
%----------------------------------------------------------------------------------------

\printindex

\end{document}
